%------------------------------------------------------------------------------%
% SELF-ADJOINT LINEAR OPERATORS                                                %
%------------------------------------------------------------------------------%
% File  : khMATH/src/sao.tex                                                   %
% Author: Klaus Hoeflinger, Beethovenstrasse 11, D-73117 Wangen                %
% Email : khoeflinger@t-online.de                                              %
%------------------------------------------------------------------------------%

%------------------------------------------------------------------------------%
% DOCUMENT CLASS und INCLUDE FILES                                             %
%------------------------------------------------------------------------------%
\documentclass[11pt,twoside]{article}
\usepackage{geometry}
\geometry{paperheight=241mm,
          paperwidth=152mm,
          top=22mm,
          bottom=17mm,
          left=20mm,
          right=20mm,
          textwidth=112mm,
          textheight=202mm,
          headheight=10mm,
          headsep=3mm,
          footskip=9mm,
          includehead,
          includefoot,
          marginparsep=0mm,
          marginparwidth=0mm}

\renewcommand{\baselinestretch}{0.945}\normalsize

\AtBeginDocument{\setlength\abovedisplayskip{2mm}}
\AtBeginDocument{\setlength\belowdisplayskip{2mm}}

%------------------------------------------------------------------------------%
% define title formats                                                         %
%------------------------------------------------------------------------------%
\usepackage{titlesec}

\renewcommand{\thesection}{\arabic{section}}
\renewcommand{\sfdefault}{FiraSans}
\renewcommand{\ttdefault}{pcr}
\renewcommand{\rmdefault}{antiqua}

\titleformat{\part}[display]
{\normalfont \large \rmfamily \bfseries \centering}
{\small{\thepart.}}{2pt}{}

\titleformat{\section}[runin]
{\normalfont \rmfamily \bfseries}
{\thesection}{1pt}{.\;}[.]

%\titleformat{\section}
%{\normalfont \bfseries \small \filcenter}
%{\thesection.\;}{0mm}{}

\titleformat{\subsection}[runin]
{\normalfont \itshape}
{}{0mm}{}

\titleformat{\subsubsection}[runin]
{\normalfont \itshape}
{}{}{}{}

\titleformat{\paragraph}[runin]
{\normalfont}
{}{}{}{}

\titleformat{\subparagraph}[runin]
{\normalfont}
{}{}{}{}

\titlespacing*{\part}          { 0pt}{ 0pt}{ 8pt}
\titlespacing*{\section}       { 0pt}{ 7pt}{ 4pt}
\titlespacing*{\subsection}    { 0pt}{14pt}{ 6pt}
\titlespacing*{\subsubsection} { 0pt}{ 6pt}{12pt}
\titlespacing*{\paragraph}     { 0pt}{ 6pt}{12pt}
\titlespacing*{\subparagraph}  { 0pt}{ 0pt}{12pt}

%------------------------------------------------------------------------------%
% define enumerate parameters                                                  %
%------------------------------------------------------------------------------%
\usepackage{enumitem}
\setlength{\parindent}{3pt}
\setlength{\parskip}  {0pt}

%\setlist{noitemsep}

\setlist[enumerate]{
         leftmargin=12pt,
         rightmargin=0pt,
         itemsep=1pt,
         itemindent=3pt,
         topsep=3pt,
         labelsep=4pt,
         partopsep=0pt,
         labelwidth=0pt,
         listparindent=0pt,
         parsep=0pt}

\setlist[itemize]{
         leftmargin=12pt,
         rightmargin=0pt,
         itemsep=0pt,
         itemindent=1pt,
         topsep=1pt,
         labelsep=4pt,
         partopsep=1pt,
         labelwidth=0pt,
         listparindent=0pt,
         parsep=0pt}

%------------------------------------------------------------------------------%
% define packages                                                              %
%------------------------------------------------------------------------------%
\usepackage{upref}
\usepackage{amssymb}
\usepackage{slantsc}
\usepackage{amsmath}
\usepackage{amsthm}
\usepackage{thmtools}
\usepackage{amscd}
\usepackage{verbatim}
\usepackage{latexsym}
\usepackage{multicol}
\usepackage{graphicx}
\usepackage{setspace}
\usepackage[ngerman,english]{babel}
\usepackage[protrusion=true,expansion=true]{microtype}
\usepackage{tikz}
\usetikzlibrary{arrows,arrows.meta,decorations.markings}
\usepackage{mathrsfs}
\usepackage{amsfonts}
\usepackage{datetime2}
\usepackage{textgreek}
\usepackage{hyperref}
\usepackage{pifont}
\usepackage{relsize}
%\usepackage{biblatex}
\relsize{-0.05}
\usepackage[skip=0mm]{parskip}
\usetikzlibrary{decorations.pathmorphing}

\usepackage{mlmodern}
\usepackage[T5,T1]{fontenc}
\usepackage{ragged2e}

%\usepackage{mathabx}
%\usepackage{times}

%\usepackage{fontspec}
%\setmainfont{Nimbus Roman No9 L}
%\usepackage[T1]{fontenc}

%------------------------------------------------------------------------------%
% define page layout                                                           %
%------------------------------------------------------------------------------%
\usepackage{fancyhdr}
\pagestyle{fancy}
\setlength{\headheight}{30.0pt}
\renewcommand{\headrulewidth}{0pt}

\AtBeginDocument{
  \addtolength\abovedisplayskip{-0.0\baselineskip}
  \addtolength\belowdisplayskip{-0.0\baselineskip}
}

\fancyhead{}
\fancyfoot{}
\addtolength{\headwidth}{0.02\marginparwidth}

\makeatletter
\let\ps@plain\ps@fancy
\makeatother

\renewcommand\sectionmark[1]
{
 \def\sectionname{#1}
 \markright{\thesection #1}
}

\makeatletter
\@addtoreset{section}{part}
\makeatother

%------------------------------------------------------------------------------%
% define footnote without number                                               %
%------------------------------------------------------------------------------%
\newcommand{\myfootnote}[1]{%
\renewcommand{\thefootnote}{}%
\footnotetext{#1}%
\renewcommand{\thefootnote}{\arabic{footnote}}%
}

%------------------------------------------------------------------------------%
% define numbering in equations                                                %
%------------------------------------------------------------------------------%
\numberwithin{equation}{section}

%------------------------------------------------------------------------------%
% define newtheorem                                                            %
%------------------------------------------------------------------------------%
\declaretheoremstyle[
headfont=\scshape, 
headindent = 0mm, %\parindent,
%postheadhook = {\hspace{0mm}\newline},
%postheadspace = \newline, 
spaceabove = 3mm, 
spacebelow = 3mm,
numberwithin=section,
name=Theorem]{khMATH}
\declaretheorem[style=khMATH]{theorem}

\declaretheoremstyle[
headfont=\bfseries, 
headindent = 0mm, %\parindent,
%postheadhook = {\hspace{0mm}\newline},
%postheadspace = \newline, 
spaceabove = 3mm, 
spacebelow = 3mm,
numberwithin=part,
name=Theorem]{khMATH1}
\declaretheorem[style=khMATH1]{theorem1}

\declaretheoremstyle[
headfont=\scshape, 
headindent = 0mm, %\parindent,
%postheadhook = {\hspace{0mm}\newline},
%postheadspace = \newline, 
spaceabove = 3mm, 
spacebelow = 3mm,
numberwithin=section,
name=Corollary]{khMATH1}
\declaretheorem[style=khMATH1]{corollary}

%            name         counter     label              numeration-mode
\theoremstyle{plain}
\newtheorem*{introduction}            {\sc Introduction}
\newtheorem {standard}                {}                 [section]
\newtheorem*{definition}              {\sc Definition}
\newtheorem{satz1}                    {\sc Satz}
\newtheorem{lemma}                    {\sc Lemma}
\newtheorem*{proposition}             {\sc Proposition}
%\newtheorem{corollary}               {\sc Corollary}
\newtheorem*{corollar}                {\sc Korollar}
\newtheorem*{remark}                  {\sc Remark}
\newtheorem*{remarks}                 {\sc Remarks}
\newtheorem*{example}                 {Example}
\newtheorem*{examples}                {Examples}
\newtheorem*{theo}                    {\sc Theorem}

%------------------------------------------------------------------------------%
% define tabular                                                               %
%------------------------------------------------------------------------------%
\usepackage{tabularx}
\newcolumntype{L}[1]{>{\raggedright\arraybackslash}p{#1}}
\newcolumntype{C}[1]{>{\centering\arraybackslash}  p{#1}} 
\newcolumntype{R}[1]{>{\raggedleft\arraybackslash} p{#1}} 

%------------------------------------------------------------------------------%
% define \sh, \zB                                                              %
%------------------------------------------------------------------------------%
\def\dsh{{d.\,h.}}
\def\ie{{i.\,e.}}
\def\zB{z.\,B.}
\renewcommand{\abstractname}{ }

%------------------------------------------------------------------------------%
% change default spacing for cases                                             %
%------------------------------------------------------------------------------%
\makeatletter
\renewcommand*\env@cases[1][1.6]{%
  \let\@ifnextchar\new@ifnextchar
  \left\lbrace
  \def\arraystretch{#1}%
  \array{@{}l@{\quad}l@{}}%
}
\makeatother

\newenvironment{rcases}
  {\left.\begin{aligned}}
  {\end{aligned}\right\rbrace}

%------------------------------------------------------------------------------%
% change default for \predisplaypenalty                                        %
%------------------------------------------------------------------------------%
\predisplaypenalty=990
\allowdisplaybreaks \numberwithin{equation}{section}

%------------------------------------------------------------------------------%
% general notations                                                            %
%------------------------------------------------------------------------------%
\def\*{\star}
\def\={\,=\,}
\def\+{\,+\,}
\def\xsubset{{\,\subset\,}}
\def\xtimes{{\,\times\,}}
\def\xeof{{\,\in\,}}
\def\eq{{\,=\,}}
\def\xto{{\,\to\,}}
\def\eqdef{\,:=\,}
\def\:{{\,:\,}}
\def\ele{{\,\in\,}}
\def\exists{\mbox{ exists }}
\def\and{\mbox{ and }}
\def\for{\mbox{ for }}
\def\fa{\mbox{ for all }}
\def\fe{\mbox{ for each }}
\def\qfa{\quad \mbox{ for all }}
\def\df{\,\dotfill}
\def\iff{if and only if }
\def\wrt{with respect to}

\makeatletter
\newcommand \Df {\leavevmode \cleaders \hb@xt@ .70em{\hss .\hss }\hfill \kern \z@}
\makeatother

\def\sql{\mathfrak{a}}               % sesquilinear form a
\def\DF{B}                           % Dirichlet form a
%------------------------------------------------------------------------------%
% number sets and special elements                                             %
%------------------------------------------------------------------------------%
\def\P{\mathbb{H}}                   % half plane of $\C$
\def\J{I}                            % interval I
\def\N{{\mathbb{N}}}                 % unsigned integers
\def\Q{{\mathbb{Q}}}                 % rational integers
\def\Z{{\mathbb{Z}}}                 % integers
\def\R{{\mathbb{R}}}                 % real numbers
\def\Rn{{\mathbb{R}}^n}              % real numbers in Rn
\def\Rd{{\mathbb{R}}^d}              % real numbers in Rn
\def\C{{\mathbb{C}}}                 % complex numbers
\def\F{{\mathbb{K}}}                 % arbitrary field
\def\Torus{{\mathbb{T}}}             % complex circle line
\def\T{\tau}                         % upper bound of interval (0,t)
\def\sign#1{\operatorname{sign}(#1)} % sign of a number

%------------------------------------------------------------------------------%
% set theory                                                                   %
%------------------------------------------------------------------------------%
\def\G{G}                            % domain $G$
\def\clG{\overline{\G}}              % closure of $G$
\def\bndG{\partial \G}               % boundary of $G$
\def\setA{\mathcal{A}}               % set A
\def\setB{\mathcal{B}}               % set B
\def\setC{\mathcal{C}}               % set C
\def\setM{M}                         % set M
\def\setN{N}                         % set N
\def\setS{\mathcal{S}}               % set S
\def\famF{\mathfrak{F}}              % family of sets
\def\indI{\mathcal{I}}               % index set I
\def\pot#1{\mathfrak{P}(#1)}         % power set

\def\relR{\mathbf{R}}                % relation R
\def\mapf{f}                         % mapping f
\def\mapg{g}                         % mapping g

\def\set#1#2{\{ \, #1 \,: \, #2 \, \}}
\def\tensor#1#2{#1 \otimes #2}
\def\Tens{\oplus}

\def\cal#1{{\mathcal{#1}}}
\def\aut#1{\operatorname{Aut}(#1)}

\def\cross#1#2{#1 {\,\times\,} #2}
\def\multicross#1#2{#1 \times  \ldots \times  #2}
\def\multitens#1#2 {#1 \otimes \ldots \otimes #2}

%------------------------------------------------------------------------------%
% group theory                                                                 %
%------------------------------------------------------------------------------%
\def\1{{\mathsf{1}}}                % unit element of a monoid
\def\0{{\mathsf{0}}}                % unit element of an Abelian monoid
\def\kernel#1{\mathfrak{N}(#1)}     % kernel of a homomorphism
\def\cokernel#1{\mathfrak{C}{(#1)}} % cokernel of a homomorphism
\def\Hom#1{\operatorname{Hom}(#1)}

%------------------------------------------------------------------------------%
% linear algebra                                                               %
%------------------------------------------------------------------------------%
\def\vecV{\mathit{V}}               % vector space V
\def\vecH{\mathit{H}}               % vector space H
\def\vecW{\mathit{W}}               % vector space W
\def\vecU{\mathit{U}}               % subspace U
\def\convC{\mathcal{C}}             % convex set C
\def\basH{\mathfrak{B}}             % Hamel base B

\def\LO#1{\mathscr{L}(#1) }         % algebra of linear transformations

\def\scalar#1#2{ \langle #1,#2 \rangle }
\def\trace#1{\operatorname{tr}(#1)}
\def\dim#1{{\operatorname{dim} \, #1 }}
\def\codim#1{{\operatorname{codim} \, #1 }}
\def\dim{\operatorname{dim}}
\def\codim{\operatorname{codim}}
\def\ind{\operatorname{ind}}

\def\genG{\cal{G}}

\def\algA{\cal{A}}
\def\algB{\cal{B}}
\def\algU{\cal{U}}

\def\idI{\cal{I}}

%------------------------------------------------------------------------------%
% topology                                                                     %
%------------------------------------------------------------------------------%
\def\compK{{K}}          % compact set C
\def\openO{\mathcal{O}}  % open set O
\def\openU{\mathcal{U}}  % open set U
\def\U{\mathcal{U}}      % neighbourhood U
\def\V{\mathcal{V}}      % neighbourhood V
\def\d{\mathit{d}}       % metric function d

\def\interior#1{\mathring{#1}}
\def\closure#1{{\overline{#1}}}

\def\topT{\mathcal{T}}   % topology T
\def\topX{X}             % topological space X
\def\topY{Y}             % topological space Y
\def\topZ{Z}             % topological space Z
\def\metrS{M}            % metric space M

\def\grpG{G}             % topological group G
\def\grpA{A}             % topological group A
\def\grpB{B}             % topological group B
\def\grpC{C}             % topological group C
\def\grpH{H}             % topological group H
\def\grpU{U}             % topological group U
\def\grpV{V}             % topological group V
\def\topG{G}             % topological group G
\def\topH{H}             % topological group H
\def\topU{U}             % topological subgroup U

\def\subS{\cal{S}}
\def\riR{R}              % ring R

%------------------------------------------------------------------------------%
% calculus                                                                     %
%------------------------------------------------------------------------------%
\def\supp#1{\operatorname{supp}(#1)}
\def\singsupp#1{\operatorname{sing \, supp}(#1)}

\def\re{\operatorname{Re}}
\def\im{\operatorname{Im}}
\def\grad{\operatorname{grad}}

%\def\abs#1 {\left|#1\right|}
\def\abs#1 {\lvert #1 \rvert}
%\def\abs#1 {\left|\,#1\,\right|}

%------------------------------------------------------------------------------%
% measure theory                                                               %
%------------------------------------------------------------------------------%
\def\salgA{\Sigma}               % sigma-algebra S
\def\salgB{\cal{B}}              % sigma-algebra B
\def\Borel#1{\cal{B}_{#1}}       % Borel-algebra 
\def\LB{\lambda}                 % Lebesgue-Borel measure

%------------------------------------------------------------------------------%
% function spaces                                                              %
%------------------------------------------------------------------------------%
\def\Lp{L^p(\G)}                 % spaces L_p
\def\Lq{L^q(\G)}                 % spaces L_q
\def\Lh{L^2(\G)}                 % spaces L_2
\def\Li{L^{\infty}(\G)}          % spaces L_infty
\def\Lc{L_{loc}^1(\G)}           % spaces L_1^{loc}
\def\Ck{C^k(\G)}                 % spaces C_k

\def\BV#1{BV(#1) }               % set of functions of bounded variation
\def\AC#1{AC(#1) }               % set of absolutely continuous functions

\def\MP#1{\mathscr{M}^+(#1) }    % set of positive measures
\def\MS#1{\mathscr{M}^-(#1) }    % vector space of finite signed measures
\def\MC#1{\mathscr{M}(#1) }      % vector space of complex measures
\def\MCR#1{\mathscr{M}_{r}(#1) } % vector space of regular complex measures

\def\SobW{W^{k,p}(\G)}                % Sobolev space W
\def\SobO{\mathring{W}^{k,p}(\G)}     % Sobolev space W
%\def\WN#1#2{\mathring{W}^{#1}(#2) }   % Sobolev space W
\def\WN#1#2{W_0^{#1}(#2) }   % Sobolev space W
%\def\HN#1#2{\mathring{H}^{#1}(#2) }   % Sobolev space H
\def\HN#1#2{H_0^{#1}(#2) }   % Sobolev space H
%------------------------------------------------------------------------------%
% functional analysis                                                          %
%------------------------------------------------------------------------------%
\def\snorm#1{\lVert #1 \rVert}
\newcommand{\norm}[1]{\left\lVert #1 \right\rVert}

\def\topV{V}                       % topological vector space V
\def\topW{W}                       % topological vector space W
\def\normS{X}                      % normed space

\def\banX{\mathit{X}}              % Banach space X
\def\banY{\mathit{Y}}              % Banach space Y
\def\banZ{\mathit{Z}}              % Banach space Z
\def\dbanX{\mathit{X^*}}           % dual space of Banach space X
\def\dbanY{\mathit{Y^*}}           % dual space of Banach space Y
\def\dbanZ{\mathit{Z^*}}           % dual space of Banach space Z

\def\banA{\mathit{A}}              % Banach algebra A
\def\banB{\mathit{B}}              % Banach algebra B
\def\banC{\mathit{C}}              % C*-algebra C


\def\domain#1{\mathfrak{D}(#1)}    % domain
\def\range#1{\mathfrak{R}(#1)}     % range

\def\isoJ{\mathfrak{J}}            % isometry J
\def\repA{{\cal{A}_{\sql}}}        % representation operator A
\def\linS{{S}}                     % linear operator S
\def\linG{{G}}                     % linear operator G
\def\linL{{L}}                     % linear operator L
\def\linT{{T}}                     % linear operator A
\def\linA{\mathit{A}}              % linear operator A
\def\linB{{B}}                     % linear operator B
\def\linM{{M}}                     % linear operator M
\def\linU{{U}}                     % linear operator U
\def\linI{{I}}                     % linear operator I
\def\A{\cal{A}}                    % linear differential operator A
\def\BDO{\cal{B}}                  % linear boundary differential operator B
\def\linU{{U}}                     % unitary operator U
\def\linP{{P}}                     % projection P
\def\compA{k}                      % compact operator k
\def\linFR{f}                      % finite rank operator f
\def\fredA{A}                      % Fredholm operator F

\def\HAM{\cal{H}}                  % Hamilton operator
\def\PA{\partial^{\alpha}}         % elementary differential operator
\def\DO{\cal{A}}                   % linear differential operator
\def\BC{\cal{B}}                   % boundary operator
\def\SLO{\cal{L}_{p,q}}            % Sturm-Liouville operator
\def\MO{\cal{L}_{M}}               % momentum operator

\def\HS#1{S(#1)           }        % algebra of Hilbert-Schmidt operators
\def\FR#1{F(#1)           }        % algebra of finite rank operators
\def\TC#1{N(#1)           }        % algebra of trace class operators
\def\BU#1{\mathscr{U}(#1) }        % algebra of unitary linear operators
\def\BO#1{\mathscr{L}(#1) }        % algebra of bounded linear operators
\def\CO#1{\mathscr{C}(#1) }        % algebra of compact linear operators
\def\CL#1{\mathscr{C}(#1) }        % algebra of compact linear operators
\def\FO#1{\Phi(#1) }               % set of Fredholm operators
\def\FK#1#2{\Phi_{#2}(#1) }        % set of Fredholm operators with index k

\def\hilH{\mathit{H}}              % Hilbert space H
\def\clU{\mathit{U}}               % closed subspace U

\def\orthO{\mathfrak{O}}           % orthonormal system O
\def\orthS{\mathfrak{S}}           % orthonormal system S
\def\basB{\mathfrak{B}}            % basis B of a Hilbert space

\def\GL#1 {\mathbf{GL}(#1)}        % linear group
\def\OG#1 {\mathbf{O}(#1)}         % linear group
\def\SOG#1 {\mathbf{SO}(#1)}       % linear group

%------------------------------------------------------------------------------%
% distribution theory                                                          %
%------------------------------------------------------------------------------%
\def\disF{F}                  % distribution F
\def\disG{G}                  % distribution G
\def\disH{H}                  % distribution H
\def\disU{U}                  % distribution U
\def\disT{T}                  % distribution T
\def\SF{C^{\infty}(\G)}       % space of smooth functions
\def\TF{C_c^{\infty}(\G)}     % space of compactly supported smooth functions
\def\TFRd{C_c^{\infty}(\R^d)} % space of test functions
\def\SRd{\mathscr{S}(\R^d)}   % Schwartz space
\def\SR1{\mathscr{S}(\R)}     % Schwartz space
\def\B1#1{C^1(#1) }           % space of continuously differentiable functions
\def\Bm#1{C^m(#1) }           % space of continuously differentiable functions
\def\Ck{C^k(\G)}
\def\TD{\mathscr{S}'(\R^d)}   % space of temperated distributions
\def\E{\cal{E}}               % space of distributions with compact support
\def\D{\mathscr{D}}           % D
\def\FT{\mathcal{F}}          % Fourier transformation F
\def\FT{\mathscr{F}}          % Fourier transformation F

%------------------------------------------------------------------------------%
% other definitions                                                            %
%------------------------------------------------------------------------------%
\def\Proof{\textsc{Proof}}
\def\FRECHET{Fr\'{e}chet}
\def\ARZELA{Arzel\'{a}}
\def\GARDING{G\r{a}rding}
\def\HOELDER{H\"older}
\def\SCHROEDINGER{Schr\"odinger}
\def\JOERGENS{J\"orgens}
\def\WUEST{W\"uest}
\def\KAEHLER{K\"aehler}
\def\HOEFLINGER{H\"oflinger}
\def\POINCARE{Poincar\'{e}}
\def\FA{Functional Analysis}

\def\Author   {K.\,{\HOEFLINGER}}
\def\AuthorUC {K.\,HOEFLINGER}
\def\Version  {1.0}
\def\Email    {khoeflinger@t-online.de}


%------------------------------------------------------------------------------%
%                              DOCUMENT SETTINGS                               %
%------------------------------------------------------------------------------%
\def\Title  {Self-Adjoint Linear Operators}
\def\Author {Klaus H\"{o}flinger}
\def\Version{1.0}
\def\Date   {30.10.2021}
\def\Email  {khoeflinger@t-online.de}

%------------------------------------------------------------------------------%
%                                BEGIN OF DOCUMENT                             %
%------------------------------------------------------------------------------%
\begin{document}
\thispagestyle{empty}

\titleformat{\part}[display]
{\normalfont \large \rmfamily \bfseries \centering}
{\small{Part \thepart}}{2pt}{}

\titleformat{\section}%[runin]
{\normalfont \bfseries \filcenter}
{\S \thesection.\;}{0mm}{}

\titlespacing*{\section}{0pt}{16pt}{ 4pt}

%------------------------------------------------------------------------------%
% FOOTNOTE                                                                     %
%------------------------------------------------------------------------------%
\myfootnote
{
  Version {\Version}, {\today};
  Email: \textit{\Email}
}

\centerline{\large{\textbf{\Title}}}
\vspace{4pt}

\centerline{\small{\textsc{\Author}}}
\vspace{9pt}

\pagenumbering{roman}
\fancyhead[C]{\small{\textsc{\Title}}}
\fancyhead[R]{\small{\thepage}}
\fancyhead[EC]{\small{\textsc{Introduction}}}
\renewcommand*{\thepage}{(\roman{page})}
\setcounter{page}{1}

The theory of \textit{self-adjoint linear operators} in Hilbert space
proves to be one of the first well-elaborated subjects
within the mathematical area of Functional Analysis.
Its corner stones have mainly been developed by the mathematicians 
\textsc{D.\,Hilbert},
\textsc{F.\,Riesz},
\textsc{E.\,Schmidt},
\textsc{J.\,von\,Neumann},
\textsc{M.\,H.\,Stone} and others in the first half of the 20th century.

\subparagraph*{}
A great portion of motivation to develop
an abstract theory of self-adjoint linear transformations came
from questions and problems arising in mathematical physics.
For example,
the mathematical formalism for quantum mechanical systems leads
in a very natural way to self-adjoint operators acting in $L^2(\Rd)$,
the Hilbert space of all 
measurable and square-integrable functions $f\:\Rd \xto \C$.
Other problems,
like the Sturm-Liouville eigenvalue problem and
the Dirichlet and the Neumann boundary value problem
for the negative Laplacian $-\Delta$,
also lead to the theory of self-adjoint operators
in the respective Hilbert spaces $L^2(\J)$ and $L^2(\G)$,
where $\J \xsubset \R$ is a finite interval 
and $\G$ is a bounded domain in some Euclidean space $\Rd$.

\subparagraph*{}
Self-adjoint operators are very suitable to be studied
by \textit{spectral theory},
hence the spectral properties of such an operator are well-understood.
Moreover,
the spectral theorem for unbounded self-adjoint operators is one
of the most celebrated assertions in Functional Analysis
and can be seen as the highlight of our subject.
But besides this, 
a further fact will be of interest:
if a self-adjoint operator is even \textit{half-bounded} in a sense,
then it appears as the infinitesimal generator
of a strongly continuous and analytic semigroup,
respectively,
which is of great importance in the solution theory of evolution equations
of hyperbolic and parabolic type
governed by a half-bounded self-adjoint operator.

\paragraph*{}
Troughout we will work with the sets $\R$ and $\C$ as the underlying fields
of the normed linear spaces under consideration.
Furthermore,
the Hilbert spaces mentioned in the sequel are always supposed to be
\textit{complex} if nothing else is said.
This is profitable,
since spectral theoretical methods can only be applied
to operators acting in complex Banach and Hilbert spaces.

\paragraph*{}
The content of this treatise is divided into three parts.
The first one describes some fundamental tools and concepts
in Functional Analysis and will be presented without proofs,
but often with references to them,
which can be found in any standard textbooks on Functional Analysis
(see bibliography).
The second part contains the abstract theory of symmetric and self-adjoint
linear operators in a separable Hilbert space
with the \textit{spectral theorem} as its main assertion.
Further topics,
like perturbation theory, extention theory and the functional calculus
for self-adjoint operators are included.
Finally,
in part III we apply the results to \textit{realizations}
of formally self-adjoint partial differential expressions
given in divergence form
\[
\DO(x,u) \eqdef
\sum_{\abs{\alpha} \leq m, \abs{\beta} \leq m} (-1)^{\abs{\alpha} } \,
\partial^{\alpha}(\,a_{\alpha \beta}(x) \, \partial^{\beta} u(x))
\quad (x \xeof \G),
\]
where formally self-adjointness of $\DO$ means
$a_{\alpha \beta}(x) \eq a_{\beta \alpha}$(x) for all $x \xeof \G$ and
all functions $a_{\alpha \beta}$.
Hence the realizations of such expressions 
are nothing but self-adjoint linear operators
in one of the spaces $L^2(\G)$ for $\G=\Rd$ or $\G \subsetneq \Rd$
and some $d \xeof \N$.

\newpage
\thispagestyle{empty}
\centerline{\textbf{\large{Table of Contents}}}
\vspace{6mm}
\begin{small}
\centerline{\small{Part I}}
\vspace{1mm}
\centerline{\textbf{Functional Analysis}}
\vspace{2mm}

\centerline
{
\begin{tabular}{R{5mm}L{100mm}R{4mm}}
%------------------------------------------------------------------------------%
 1. & Continuous Linear Operators                                    \Df &  1 \\
 2. & Inner Products and Hilbert Spaces                              \Df &  3 \\
 3. & Closed and Closable Linear Operators                           \Df &  6 \\
 4. & Adjoint Operators                                              \Df &  9 \\
 5. & Semi-Fredholm Operators                                        \Df & 10 \\
 6. & Resolvent Set and Spectrum                                     \Df & 12 \\
 7. & Compact Linear Operators and Fredholm Operators                \Df & 15 \\
 8. & Sesquilinear Forms and Associated Operators                    \Df & 18 \\
 9. & Linear Operator-Differential Equations                         \Df & 20 \\
%------------------------------------------------------------------------------%
\end{tabular}
}

\vspace{4mm}
\centerline{\small{Part II}}
\vspace{1mm}
\centerline{\textbf{Abstract Theory of Self-Adjoint Linear Operators}}
\vspace{2mm}

\centerline
{
\begin{tabular}{R{5mm}L{100mm}R{4mm}}
%------------------------------------------------------------------------------%
10. & Symmetric and Half-Bounded Operators                           \Df & 25 \\
11. & The Concept of Self-Adjointness                                \Df & 27 \\
12. & Self-Adjoint Extensions of Symmetric Operators                 \Df & 28 \\
13. & Perturbation Theory                                            \Df & 31 \\
14. & The Spectrum of Symmetric and Self-Adjoint Operators           \Df & 32 \\
15. & Projection-Valued Measures and Spectral Integrals              \Df & 35 \\
16. & The Spectral Theorem                                           \Df & 37 \\
17. & Self-Adjoint Operators with Pure Point Spectrum                \Df & 40 \\
18. & Evolution Equations Governed by Self-Adjoint Operators         \Df & 41 \\
19. & Positive Self-Adjoint Operators Generate Semigroups            \Df & 43 \\
%------------------------------------------------------------------------------%
\end{tabular}
}

\vspace{4mm}
\centerline{\small{Part III}}
\vspace{1mm}
\centerline{\textbf{Self-Adjoint Realizations of Differential Expressions}}
\vspace{2mm}

\centerline
{
\begin{tabular}{R{5mm}L{100mm}R{4mm}}
%------------------------------------------------------------------------------%
20. & Fourier-Plancherel Transformation                              \Df & 46 \\
21. & Hilbert-Sobolev Spaces                                         \Df & 47 \\
22. & Realizations of Linear Partial Differential Operators          \Df & 49 \\
23. & Formally Self-Adjoint Differential Expressions                 \Df & 50 \\
24. & The Momentum Operator on Various Intervals                     \Df & 51 \\
25. & Sturm-Liouville Operators                                      \Df & 53 \\
26. & The Laplacian in Bounded and Unbounded Domains                 \Df & 54 \\
27. & {\SCHROEDINGER} Operators                                      \Df & 56 \\
28. & Strongly Elliptic Linear Differential Operators of Order $2m$  \Df & 58 \\
%------------------------------------------------------------------------------%
\end{tabular}
}

\vspace{4mm}
\centerline
{
\begin{tabular}{R{5mm}L{100mm}R{4mm}}
%------------------------------------------------------------------------------%
    & Bibliography                                                   \Df & 53 \\
%------------------------------------------------------------------------------%
\end{tabular}
}
\end{small}

%------------------------------------------------------------------------------%
%                                   CONTENTS                                   %
%------------------------------------------------------------------------------%
\newpage
\pagenumbering{arabic}
\setcounter{page}{1}
\renewcommand*{\thepage}{\arabic{page}}

\thispagestyle{empty}
%------------------------------------------------------------------------------%
%                                    PART I                                    %
%------------------------------------------------------------------------------%
\part{Functional Analysis}
%------------------------------------------------------------------------------%
Let us start our review of Functional Analysis with a short exposition
of normed linear spaces.
Since we do not need very much knowledge on this general class of spaces,
we want to restrict ourselves to the definition of a normed linear space,
the standard topology on such a space and the notion of \textit{completeness}.

\subparagraph*{}
Throughout this text let $\F$ be one of the fields $\R$ and $\C$.
By a \textit{normed space} $\banX$ we mean a $\F$-linear space
that differs from other abstract linear spaces
by the presence of a particular mapping $\norm{\,\cdot\,}\: \banX \xto \R$,
denoted the \textit{norm} on $\banX$,
which fulfills for all $x,y \xeof \banX$ and $\lambda \xeof \F$

\begin{itemize}[leftmargin=22mm,itemsep=2pt]
\item[(N-1)]
$\norm{x} \, \geq \, 0$,

\item[(N-2)]
$\norm{\lambda \, x} \eq \abs{\lambda} \cdot \norm{x}$,

\item[(N-3)]
$\norm{x+y} \, \leq \, \norm{x} + \norm{y}$
\; \textit{(triangle inequality)},

\item[(N-4)]
$\norm{x} \eq 0$ \iff $x = \0$.
\end{itemize}

\subparagraph*{}
Linear spaces can of course be furnished with different norms.
Arbitrary norms $\norm{\,\cdot\,}_{\alpha}$ and $\norm{\,\cdot\,}_{\beta}$
defined on a linear space $\banX$
are called \textit{equivalent} to each another,
if there are constants $c_1,c_2 > 0$ such that
$
c_1 \norm{x}_{\alpha} \leq
    \norm{x}_{\beta}  \leq 
c_2 \norm{x}_{\alpha}$
for all $x \xeof \banX$.

\subparagraph*{}
Each normed linear space $\banX$ may be equipped
with a very natural topology $\topT$,
called the \textit{norm topology} on $\banX$.
The definition of $\topT$ is done either by the assertion
that $\topT := \topT_{\norm{\cdot}}$ is the smallest topology on $\banX$
such that the norm mapping remains continuous,
or by the metric topology $\topT_d$,
when $\banX$ is considered as a metric space with metric function
$d(x,y) := \norm{x-y} \fa x,y \xeof \banX$.
In fact,
both definitions lead to the same collection of so-called
"open" subsets in $\banX$,
so we have $\topT := \topT_{\norm{\cdot}} \eq \topT_d$,
which justifies to speak of \textit{the} topology of a normed linear space.
Hence a multiplicity of topological concepts,
like the \textit{closure} of a set,
the \textit{neighbourhood} of a point,
the \textit{convergence} of a series and many others are meaningful
in the setting of normed spaces.
Notice
that equivalent norms on $\banX$ are creating one and the same topology.
For example,
if $\banX$ is a finite-dimensional linear space,
then all norms are equivalent
and consequently there is exactly one norm topology on $\banX$.

\subparagraph*{} %- sequentially complete normed spaces, Banach spaces
Recall that a sequence $(x_n)$ in $\banX$ is called \textit{Cauchy},
if to each $\epsilon > 0$ there is $N_{\epsilon} \xeof \N$
such that $m,n > N_{\epsilon}$ implies $\norm{x_m-x_n} < \epsilon$.
Any normed linear space $\banX$ is called \textit{sequentially complete}
(or shortly \textit{complete}),
if for any Cauchy sequence $(x_n)$ in $\banX$
there is $x \xeof \banX$
such that $\lim_{\,n \xto \infty} \norm{x_n-x} \eq 0$.
Equivalently,
$\banX$ is complete
iff
for every absolutely convergent series
$\sum_{n=1}^{\infty} \norm{x_n} < \infty$
there is $x \xeof \banX$ such that $\norm{x-x_n} \xto 0$ for $n \xto \infty$.
A complete normed space is called a \textit{Banach space}.

\subparagraph*{}
It should also be noticed
that every uncomplete normed space $\banX$ may continuously be embedded
into a Banach space $\overline{\banX}$,
denoted the \textit{completion} of $\banX$.
The procedure to construct such an embedding is an abstract version
of the embedding of the rational numbers into the real ones.

\section{Continuous Linear Operators}
%------------------------------------
\fancyhead[EC]{\small{{\thesection.} \textsc{\sectionname}}}
In the following let $\banX$ and $\banY$ be $\F$-linear spaces.
By a \textit{linear trans\-formation} or \textit{linear operator}
(shortly denoted an \textit{operator})
is meant any mapping $ \linA\:\banX \xto \banY$ fulfilling
\[
\linA(\lambda x + \mu y) = \lambda \linA x + \mu \linA y
\quad \quad (x,y \xeof \banX, \; \lambda,\mu \xeof \F),
\]
called the \textit{linearity property} of $\linA$.
The spaces $\banX$ and $\banY$ are said to be the \textit{domain space}
and the \textit{range space} of $\linA$,
respectively.
In the special case, where $\banY \eq \banX$,
it is convenient to say that $\linA$ is an operator \textit{in $\banX$}.

\subparagraph*{}
We next review the concepts of continuity and boundedness
for a linear operator $\linA\:\banX \xto \banY$.
First,
the operator $\linA$ is called \textit{continuous},
if the preimage $\linA^{-1}(\openO)$ of each open set $\openO \subset \banY$
is also open in $\banX$.
Let $\BO{\banX,\banY}$ be the linear space
of all continuous linear operators $\linA\:\banX \xto \banY$,
which forms a subspace of $L(\banX,\banY)$,
the space of \textit{all} linear but not necessarily continuous mappings.
Second,
the operator $\linA$ is said to be \textit{bounded}
\iff it is \textit{norm-bounded},
that is,
there is $c \geq 0$
such that $\norm{\linA x} \leq c \cdot \norm{x}$ for each $x \xeof \banX$.
In this case we define the \textit{operator norm} of $\linA$ according to
\begin{equation}\label{OPERATOR_NORM}
\norm{\linA} \eqdef \sup_{\norm{x}=1} \norm{\linA x}.
\end{equation}

\paragraph*{}
The following result on continuous linear operators is well-known:

\begin{theorem}
Let $\linA\:\banX \xto \banY$ be a linear operator
between the normed linear spaces $\banX$ and $\banY$.
Then the subsequent statements are equivalent to each other:

\begin{itemize}[leftmargin=15pt]
\item[(a)] $\linA$ is continuous;
\item[(b)] $\linA$ is sequentially continuous in every $x \xeof \banX$:
           if $(x_n)$ is a sequence in $\banX$
           with $\lim_{\,n \xto \infty} x_n \eq x \xeof \banX$,
           then $\lim_{\,n \xto \infty} \linA x_n \eq \linA x$;
\item[(c)] $\linA$ is sequentially continuous at the origin $\0 \xeof \banX$;
\item[(d)] $\linA$ is uniformly continuous in $\banX$;
\item[(e)] $\linA$ is bounded.
\end{itemize}
\end{theorem}

\subparagraph*{}
The space $\BO{\banX,\banY}$ of bounded linear operators
from a normed linear space $\banX$ to a Banach space $\banY$
and endowed with the norm given by \eqref{OPERATOR_NORM} is complete,
since $\banY$ is supposed to be complete,
hence $\BO{\banX,\banY}$ forms a Banach space in its own right.

\subparagraph*{}
In addition,
in case of $\banY \eq \banX$
the space $\BO{\banX} := \BO{\banX,\banX}$
constitutes a \textit{normed algebra}
with composition $(\linA_1,\linA_2) \mapsto \linA_2 \circ \linA_1$
as its multiplicative operation
and owns the property
$\norm{\linA_1 \circ \linA_2} \leq \norm{\linA_1} \cdot \norm{\linA_2}$.
It is called the \textit{operator algebra}
associated with the normed linear space $\banX$.

\paragraph*{}
Finally,
we turn to the notion of \textit{embedding},
which is a special kind of a continuous linear operator
and leads to the concept of \textit{norm isomorphism}.

\subparagraph*{}
A linear mapping $\iota:\banX \xto \banY$ is called
a (continuous) \textit{embedding} from a Banach space $\banX$
into a Banach space $\banY$,
if it is injective (and continuous).
If $\iota$ is continuous and even surjective,
it is said to be a \textit{linear homeomorphism}
between $\banX$ and $\banY$
(in the sense of topological linear spaces).
Then \textsc{Banach}\textit{'s isomorphism theorem} claims
that the inverse map $\iota^{-1}$ is also continuous,
so it forms a norm isomorphism, too.
This statement is a direct consequence of the \textit{open mapping theorem}
(see \ref{OPEN_MAPPING_THEOREM}).
Furthermore,
if the mapping $\iota$ is \textit{norm-preserving},
that is,
\[
\norm{\iota x}_Y = \norm{x}_X \quad \quad (x \xeof \banX),
\]
we denote this mapping a \textit{Banach space isomorphism}
between $\banX$ and $\banY$,
which in this case are called \textit{norm-isomorphic} to each other
(notation: $\banX \cong \banY$).

\section{Hilbert Spaces}
%-----------------------
A \textit{Hilbert space} is a special kind of a Banach space,
where the norm is derived from an \textit{inner product}.
It is a matter of fact 
that certain investigations done in the setting of Hilbert spaces
are more transparent than the comparable ones in arbitrary Banach spaces.

\subparagraph*{}
Given a real or complex linear space $\hilH$,
the \textit{inner product} in $\hilH$ is a mapping
$\scalar{\,\cdot\,}{\,\cdot\,} \: \hilH \xtimes \hilH \xto \F$,
which is

\begin{itemize}[leftmargin=15mm]
\item[(I-1)]
linear in the first argument,

\item[(I-2)]
either linear or semilinear in the second argument (conditional upon,
whether $\F$ is the real or the complex number field),
\end{itemize}
together with the additional properties

\begin{itemize}[leftmargin=15mm]
\item[(I-3)]
$\scalar{y}{x} =           \scalar{x}{y}$  ($\F=\R$) and
$\scalar{y}{x} = \overline{\scalar{x}{y}}$ ($\F=\C$), respectively,
\item[(I-4)]
$\scalar{x}{x} \geq 0$ for all $x \neq 0$,
\item[(I-5)]
$\scalar{x}{x} = 0$ implies $x = \0$.
\end{itemize}

\subparagraph*{}
Property (I-3) is denoted the \textit{Hermitean symmetry} of the inner product,
while (I-4) and (I-5) together mean positive definiteness.
The presence of the inner product gives rise to define a norm in $\hilH$
according to
\begin{equation}\label{DERIVED_NORM}
\norm{x} \eqdef \scalar{x}{x} ^{1/2} \quad \quad (x \xeof \hilH),
\end{equation}
so each linear space endowed with an inner product is normed in a natural way.
Inner product and norm are related to each other
by the \textit{Cauchy-Schwarz} \textit{inequality}
\begin{equation}\label{CAUCHY_SCHWARZ_INEQUALITY}
\abs{\scalar{x}{y}} \, \leq \,
\norm{x} \cdot \norm{y} \quad \quad (x,y \xeof \hilH),
\end{equation}
thus the inner product is continuous
{\wrt} the product topology on $\hilH \xtimes \hilH$.
An inner product space $\hilH$ is called a \textit{Hilbert space},
if it is complete {\wrt} the norm defined by \eqref{DERIVED_NORM}.

\subparagraph*{}
The pleasant geometric properties of inner product spaces
may be attributed to the \textit{parallelogram equation}
\begin{equation}\label{PARALLELOGRAM_EQUATION}
2 \norm{x}^2 + 2 \norm{y}^2 \= \norm{x+y}^2 + \norm{x-y}^2
\qfa x,y \xeof \hilH.
\end{equation}
In general,
the norm $\norm{\,\cdot\,}$ on a linear space $\banX$
can be induced by an inner product
\iff \eqref{PARALLELOGRAM_EQUATION} holds.
In this case,
and if $\banX$ is a complex linear space,
the inner product $\scalar{\cdot}{\cdot}_{\norm{\cdot}}$
can be recovered from the norm by the \textit{polarization formula}
\[
\scalar{x}{y}_{\norm{\cdot}}
\=
\frac{1}{4} \,
(\,
\norm{x + y}^2 - \norm{x-y}^2 +i \norm{x+iy}^2 - i \norm{x-iy}^2
\,).
\]
If $\banX$ is a real linear space,
however,
the polarization formula simplifies to
\[
\scalar{x}{y}_{\norm{\cdot}} \=
\frac{1}{4} \,
(\,
\norm{x + y}^2 - \norm{x-y}^2
\,).
\]

\paragraph*{}
In the following we list a couple of pleasant properties of Hilbert spaces
which are not shared with more general Banach spaces.
Let $\hilH$ be a real or complex Hilbert space.

\begin{itemize}[leftmargin=5mm,itemsep=5pt]
\item[1.] \textit{Orthogonality}.
Any elements $x,y$ of a Hilbert space $\hilH$ are called \textit{orthogonal}
to each other,
if $\scalar{x}{y} \eq 0$ (notation: $x \bot y$).
A subset $\orthO \subset \hilH$ is said to be \textit{orthogonal}
or an \textit{orthogonal system},
if all elements in $\orthO$ are orthogonal to each other.
Finally,
given an arbitrary set $\setA \subset \hilH$,
the \textit{orthogonal space} $\setA^{\bot}$ contains
the elements $y \xeof \hilH$ being orthogonal to every $x \xeof \setA$.
This set always forms a linear subspace of $\hilH$.

\item[2.] \textit{The Projection Theorem}.
Let $\setA \subset \hilH$ be a \textit{closed} linear subspace.
Then for each element $x \xeof \hilH$ there is exactly one $x_0 \xeof \setA$
such that
\[
\norm{ x - x_0 } \= \inf \, \{ \, \norm {x-y}: y \xeof \setA \,\}.
\]
The element $x_0$ is called the \textit{projection} of $x$ onto $\setA$.
It minimizes the distance between $\setA$ and $x$.
The mapping
$\pi_{\setA}:x \xeof \hilH \mapsto x_0 \xeof \setA$
is also called the \textit{projection} of $\hilH$ onto $\setA$.
In general,
by a projection is meant a continuous linear operator $\pi:\hilH \xto \hilH$
satisfying the conditions $\pi \circ \pi \eq \pi$
and $\scalar{\pi x}{y} \eq \scalar{x}{\pi y}$ for all $x,y \xeof \hilH$.

\item[3.] \textit{The Decomposition Theorem}.
Using the notation of orthogonal space $\setA^{\bot}$,
the projection of $x$ onto $\setA$ is characterized as follows:
$x_0 \eq \pi_{\setA} x$ iff $x-x_0 \xeof \setA^{\bot}$.
Hence each $x \xeof \hilH$ can be uniquely represented by the formula
$x \eq x_0+x^{\bot}$,
where $x_0 \xeof \setA$ and $x^{\bot} \xeof \setA^{\bot}$.
This means
that the space $\hilH$ is isometrically isomorphic to the direct sum
$\setA \, \oplus \, \setA^{\bot}$
of the closed subspaces $\setA$ and $\setA^{\bot}$.
This fact is known as \textit{decomposition theorem} for Hilbert spaces.
In other words,
this theorem says
that each linear subspace $\setA$ of $\hilH$
has a topological complement\footnote{
\,A linear subspace $\setA \subset \banX$ of a normed space $\banX$
is said to be topologically complementable,
if there is a further closed subspace $\setB$
such that $\setA \cap \setB \eq \{0\}$ and $[\setA \cup \setB] \eq \banX$,
hence $\setB$ is the algebraic complement of $\setA$.},
which is of course the space $\setA^{\bot}$.

\item[4.] \textit{Orthonormal Bases}.
An orthogonal system $\orthO \subset \hilH$ is called
an \textit{orthonormal system},
if in addtion $\scalar{x}{x} \eq 1$ for all $x \xeof \orthO$.
If $\orthO$ contains only \textit{finitely} many elements,
then the span $[\orthO]$ forms a \textit{closed} subspace in $\hilH$,
and the projection $\pi_{[\orthO]}$ of any element $x$ onto $[\orthO]$
can be expressed by the formula
\[
\pi_{[\orthO]} \, x \= \sum_{i=1}^n \, \scalar{x}{e_i} \, e_i.
\]
However,
if $\orthO$ consists of arbitrarily many elements,
the inequality $\scalar{x}{e_{\alpha}} \neq 0$ can only be valid
for mostly countable many elements $e_{\alpha} \xeof \orthO$,
which enables us to summarize the coefficients $\scalar{x}{e_{\alpha}}$
about all elements of $\orthO$.

\subparagraph*{}
An orthonormal system $\orthO$ is called
an \textit{orthonormal base} in $\hilH$,
if its span $[ \orthO ]$ lies dense in $\hilH$,
that is,
$\closure{ [ \orthO ] } \eq \hilH$.
In this case it is possible to represent each $x \xeof \hilH$ as a formal sum
\begin{equation}\label{FOURIER_COEFFICIENTS}
x \= \sum_{e_{\alpha} \xeof \orthO} \, \scalar{x}{e_{\alpha}} \, e_{\alpha}.
\end{equation}
The inner products $\scalar{x}{e_{\alpha}}$ are said to be the
\textit{Fourier coefficients} of $x$ {\wrt} $\orthO$.
Notice that for each $x \xeof \hilH$
the right hand side of \eqref{FOURIER_COEFFICIENTS}
consists of mostly countable many non-zero summands and
the series is invariant by realignments of the elements $e_{\alpha} \xeof \orthO$.

\subparagraph*{}
There is the need to determine,
whether a given orthonormal system $\orthO$ forms an orthonormal base.
It can be shown that
the following statements are equivalent to each other:

\begin{itemize}[leftmargin=20pt]
\item[(i)]
$\orthO$ is an orthonormal base in $\hilH$.

\item[(ii)]
$\orthO$ is \textit{complete},
that is,
$\scalar{x}{e} \eq 0$ for all $e \xeof \orthO$ implies $x \eq \0$.

\item[(iii)]
$\sum_{e_{\alpha} \xeof \,\orthO} \;
\abs{ \scalar{x}{e_{\alpha}} } ^2
\eq
\norm{x}^2$ for all $x \xeof \hilH$ \textit{(Parseval}\textit{'s equation}).
\end{itemize}

\paragraph*{}
It is also natural to ask for the existence
of an orthonormal base in arbitrary Hilbert spaces.
First of all,
the collection of all orthonormal systems in a Hilbert space $\hilH$
forms a partial ordered set,
and under the assumption that \textsc{Zorn}{'s lemma}\,\footnote{\,
\,The \textit{axion of choice} claims that
if $\cal{F}$ is a disjoint family $F_{\alpha}$ ($\alpha \xeof \J$)
of non-empty sets,
then there is a set $A$ that shares exactly one element
with each other set $F_{\alpha} \xeof \cal{F}$.}
is valid,
there is a maximal orthonormal system $\orthO$ in $\hilH$
which cannot be increased.
Indeed,
such a set forms an orthonormal base,
which ensures that
every Hilbert space $\hilH \neq \{\0\}$ has an orthonormal base.

\subparagraph*{}
Furthermore,
if $\orthO_1$ and $\orthO_2$ are orthonormal bases in $\hilH$,
then there is a bijective mapping between them,
thus they own identical cardinal numbers.
This legitimates us to introduce the notation
of \textit{Hilbert space dimension} $\abs{\hilH} $
being the cardinal number of any orthonormal base in $\hilH$.

\subparagraph*{}
A Hilbert space $\hilH$ is called \textit{separable}
if there is a mostly countable subset $D \subset \hilH$
lying dense in $\hilH$.
It can be shown that $\hilH$ is separable
\iff
$\abs{ \hilH} \leq \abs{\N} $.
In fact,
each orthonormal base in $\hilH$ forms a mostly countable set in this case,
and $\hilH$ proves to be norm-isomorphic to the space $\ell^2$
of square summarizable sequences $f\:\N \xto \C$.

\item[5.] \textsc{Riesz}\textit{'s Representation Theorem}.
Recall that the \textit{conjugate space} $\banX^*$
of a normed $\F$-linear space $\banX$
consists of all \textit{continuous} (anti)linear functionals $f\:\banX \xto \F$,
depending whether $\F=\C$ or $\F=\R$.
In the complex case we have $f(\lambda x) \eq \overline{\lambda} f(x)$
for $\lambda \xeof \C$,
whereas in the real case $f(\lambda x) \eq \lambda f(x)$ holds
for $\lambda \xeof \R$.
It is well-known that $\banX^*$ is a Banach space again with operator norm
\begin{equation}\label{NORM_OF_FUNCTIONAL}
\norm{f} \eqdef
\sup_{x \xeof \banX, \, \norm{x}=1} \, \abs{f(x)}
\quad \quad (f \xeof \banX^*).
\end{equation}

\subparagraph*{}
Due to the properties of an inner product,
each element $x \xeof \hilH$ gives rise to
a continuous (anti)linear functional $f_x \xeof \hilH^*$
according to
\[
f_x(y) := \scalar{x}{y} \quad (y \xeof \hilH).
\]
The mapping $\Phi:\hilH \xto \hilH^*$ defined by $x \mapsto f_x$
is called the \textsc{Riesz} \textit{map} of $\hilH$.
It is linear on $\hilH$,
and since the inner product is positive definite,
$\Phi$ turns out to be injective.
\textsc{Riesz}\textit{'s representation theorem} claims
that $\Phi$ is actually surjective,
so there is a canonical correspondence
between the elements of a Hilbert space $\hilH$
and those of its conjugate space $\hilH^*$:

\begin{theorem}\label{RIESZ_REPRESENTATION_THEOREM}
If $\hilH$ is a Hilbert space,
then $\Phi: \hilH \xto \hilH^*$ is bijective and norm-preserving,
that is,
$\norm{\Phi(x)} \eq \norm{x}$ for all $x \xeof \hilH$.
\end{theorem}

\subparagraph*{}
This theorem implies that
to each $f \xeof \hilH^*$ there is exactly one element $x_f \xeof \hilH$,
which fulfills the equation
$\scalar{x_f}{y} \eq f(y)$ for every $y \xeof \hilH$.

\subparagraph*{}
The \textsc{Riesz} map also gives rise to introduce an inner product
on the conjugate space $\hilH^*$ by means of
\[
\scalar{f}{g} \eqdef \scalar{\Phi^{-1}(f)}{\Phi^{-1}(g)\,}
\quad \quad (f,g \xeof \hilH^*),
\]
hence it is not only norm-preserving, but also \textit{isometric}.
\end{itemize}

\section{Closed and Closable Linear Operators}
%---------------------------------------------
It is the purpose of this section to extend
the notation of continuous linear transformation to a more general concept,
which then will be called a \textit{closed} linear operator.
This extended definition is motivated by the fact
that realizations of linear partial differential expressions defined
on certain function spaces may indeed not be continuous,
but turn out to be closed or at least \textit{closable}.

\subsubsection*{} % not everywhere defined linear operators
Let $\banX$ and $\banY$ be real rsp. complex Banach spaces.
By a \textit{linear operator}
from the domain space $\banX$ to the range space $\banY$ is meant a mapping
$
\linA\:\domain{\linA} \xsubset \banX \xto \banY,
$
where the \textit{domain of definition} $\domain{\linA}$
of the operator $\linA$ forms a linear subspace of $\banX$
and $\linA$ is \textit{linear} on $\domain{\linA}$,
{\ie}
\[
\linA(\lambda x + \mu y) \= \lambda \, \linA x \, + \, \mu \, \linA y
\qfa x,y \xeof \domain{\linA}, \; \lambda,\mu \xeof \F.
\]
In case of $\banY \eq \banX$,
we denote $\linA$ to be an operator \textit{acting in $\banX$}
or shortly to be an operator \textit{in $\banX$}.

\subparagraph*{}
Of course,
this definition includes that of section 1,
where the domain of definition $\domain{\linA}$ is assumed
to be the entire space $\banX$.
It our extended definition it is often required that 
the domain of definition $\domain{\linA}$ lies dense in $\banX$,
that is,
each $x \xeof \banX$ can be approximated
by a sequence $(x_n)$ of elements in $\domain{\linA}$.
However,
if the domain of definition $\domain{\linA}$ is \textit{not} dense in $\banX$,
it is nearby to consider the closure
$\banX_0 := \overline{\domain{\linA} } \xsubset \banX$ 
and to pay attention to the slightly modified but densely defined operator
$\linA_0:\domain{\linA} \xsubset \banX_0 \xto \banY$.

\subparagraph*{}
Rather the domains of definition for operators $\linA_1, \linA_2$
with domain space $\banX$ do not necessarily form the whole of $\banX$,
the \textit{sum} $\linA_1 + \linA_2$ and
the \textit{composition} $\linA_1 \circ \linA_2$ are well-defined
with respective domains of definition $\domain{\linA_1 + \linA_2}$
and $\domain{\linA_1 \circ \linA_2}$ given by
\[
\begin{cases}[1.2]
\quad 
\domain{\linA_1 + \linA_2} \eqdef \domain{\linA_1} \cap \domain{\linA_2},\\
\quad
\domain{\linA_1 \circ \linA_2} \eqdef
\{\, x \xeof \domain{\linA_2}: \linA_2 \, x \xeof \domain{\linA_1}\,\}.\\
\end{cases}
\]
Especially,
if $\banY \eq \banX$ and $\linA_2$ is a \textit{dilatation} in $\banX$,
{\ie}
$\linA_2 \, x \eq \lambda x$ for $x \xeof \banX$ and some $\lambda \xeof \F$,
the operator $\linA_1 + \linA_2 \eq \linA_1 + \lambda \1$ 
has domain of definition $\domain{\linA_1 + \lambda} \eq \domain{\linA_1}$.

\subparagraph*{}
Since the domain of definition $\domain{\linA}$ of a linear operator $\linA$
may be a proper linear subspace of $\banX$,
the concept of \textit{extension} is meaningful.
Let $\linA_1$ and $\linA_2$ be linear operators in $\banX$
with domains $\domain{\linA_1}$ and $\domain{\linA_2}$,
respectively.
We introduce the partial ordered relation
\[
\linA_1 \prec \linA_2
\; : \Longleftrightarrow \;
\domain{\linA_1} \, \subset \, \domain{\linA_2}
\, \mbox{ and } \;
\linA_1 \, x \= \linA_2 \, x
\fa x \xeof \domain{\linA_1}
\]
and call $\linA_2$ an \textit{extension} of $\linA_1$
iff $\linA_1 \prec \linA_2$.
In addition,
the operators $\linA_1$ and $\linA_2$ are said to be \textit{equal}
(denoted $\linA_1 \eq \linA_2$)
iff $\linA_1 \prec \linA_2$ and $\linA_2 \prec \linA_1$,
or equivalently
\[
\domain{\linA_1} \= \domain{\linA_2}
\; \mbox{ and } \;
\linA_1\,x = \linA_2\,x
\qfa x \xeof \domain{\linA_1}.
\]

\subparagraph*{}
Reversely,
let $Y_0$ be a linear subspace of the target space $\banY$.
We then may restrict the operator $\linA$ to a linear operator
$\linA_0$ defined by
$\domain{\linA_0} \eq \{x \xeof \domain{\linA}: \linA x \xeof \banY_0\}$
and $\linA_0 x := \linA x$ for $x \xeof \domain{\linA_0}$.
Clearly,
$\linA_0$ has target space $\banY_0$,
and we say that $\linA_0$ is the \textit{part of $\linA$ in $\banY_0$}.

\subparagraph*{}
Finally,
with every linear operator $\linA\:\domain{\linA} \xsubset \banX \xto \banY$
acting between the linear spaces $\banX$ and $\banY$
there are associated three fundamental linear subspaces:

\begin{itemize}[leftmargin=12pt]
\item[1.]
the \textit{kernel} $\kernel{\linA} \eq \{x \xeof \domain{\linA}: \linA x \eq \0\}
\subseteq \banX$;

\item[2.]
the \textit{range} $\range{\linA} \eq
\{\, \linA x: x \xeof \domain{\linA}\,\} \subseteq \banY$;

\item[3.]
the \textit{cokernel} $\cokernel{\linA}$,
being the algebraic complement of the range $\range{\linA}$ within $\banY$,
that is,
$\banY \eq \range{\linA} \oplus \cokernel{\linA}$.
\end{itemize}

\subsubsection*{} % continuous linear operators
Recall that a linear operator $\linA$ is called \textit{continuous}
if $x_n \xto \0$ implies $\linA x_n \xto \0$
for any sequence $(x_n)$ settled in the domain of definition $\domain{\linA}$.
If this set is dense in $\banX$ and $\linA$ is continuous,
then $\linA$ can be extended to its closure $\overline{\domain{\linA}}$,
which is continuous as well.
This is the claim of the so-called \textit{B.L.T.\,theorem},
which is an abbreviation for "Bounded Linear Transformation theorem"
(see \cite{RS1}, page 6).

\subsubsection*{} % closed operators
Unfortunately,
many linear operators in applications and especially the differential operators
studied in the theory of partial differential equations 
are often uncontinuous.
But the concept of \textit{closed operator} stands for
a certain compensation for the missing property of continuity.
If $\banX$ and $\banY$ are Banach spaces,
then a linear operator $\linA\:\domain{\linA} \xsubset \banX \xto \banY$
is called \textit{closed},
if the linear subspace
\[
\Gamma_{\linA} \eqdef
\{\,(x,\linA x): x \xeof \domain{\linA}\,\} \; \subset \; \banX \xtimes \banY,
\]
denoted the \textit{graph} of $\linA$,
is closed {\wrt} the product topology $\topT_{\banX \xtimes \banY}$
in $\banX \xtimes \banY$,
which especially implies that the kernel of $\linA$ is a closed set in $\banX$.

\subparagraph*{}
The closeness of an operator can be characterized as follows:
a linear operator $\linA\:\domain{\linA} \xsubset \banX \xto \banY$ is closed
\iff

\begin{itemize}[leftmargin=12pt]
\item[1.]
the convergence of a sequence $(x_n)$ to an element $x \xeof \banX$
and the convergence of $(\linA x_n)$ to an element $y \xeof \banY$ implies
$x \xeof \domain{\linA}$ and $\linA x \eq y$, briefly written
\[
\begin{rcases}%[1.2]
 x_n       \xto x \quad \\
 \linA x_n \xto y \quad \\
\end{rcases}
\; \Longrightarrow \;
x \xeof \domain{\linA} \, \mbox{ and } \, \linA x = y.
\]
\item[2.]
the linear subspace $\domain{\linA}$ endowed with the \textit{graph norm}
\begin{equation}\label{GRAPH_NORM}
\norm{x}_{\linA} \eqdef \norm{x}_X \, + \, \norm{\linA x}_Y,
\quad
x \xeof \domain{\linA}
\end{equation}
is complete,
that is,
$\domain{\linA}$ furnished with norm $\norm{\,\cdot\,}_{\linA}$ is Banach.
\end{itemize}

\subsubsection*{} % closable operators
It may happen that a linear operator $\linA$ fails to be closed.
However,
$\linA$ is said to be \textit{closable},
if the closure of $\Gamma_{\linA}$ is the graph
of an extension $\overline{\linA}$ of $\linA$,
called the \textit{closure} of $\linA$.
Any dense subspace $D \xsubset \banX$ is called a \textit{core}
of a closed operator $\linA$,
if $D \xsubset \domain{\linA}$ and
the closure of $\linA_{\mid D}$ resembles $\linA$.

\subparagraph*{}
A similar criterion as for closeness is available for closability:
the linear operator $\linA\:\domain{\linA} \xsubset \banX \xto \banY$ is closable
\iff
each sequence $(x_n)$ in $\domain{\linA}$ with $x_n \xto \0$
and $\linA x_n \xto y \xeof \banX$ implies $y=\0$,
shortly written
\[
\begin{rcases}%[1.2]
 x_n       \xto \0 \quad \\
 \linA x_n \xto y  \quad \\
\end{rcases}
\; \Longrightarrow \;
y \= \0.
\]

\paragraph*{}
There are of course relationships between the presented concepts.
In fact,
if $\domain{\linA}$ is closed and $\linA$ is continuous, then it is closed,
and if $\linA$ is closed, then it is closable,
thus we obtain the following chain of implications:
\[
\linA \mbox{ bounded    } \Longleftrightarrow
\linA \mbox{ continuous } \Longrightarrow
\linA \mbox{ closed }     \Longrightarrow
\linA \mbox{ closable}.
\]

\paragraph*{}
We are going on with some further questions concerning the correspondence
between continuity and closeness of linear operators,
which are meaningful in what follows.

\subparagraph*{} %- open mapping theorem
Let us start with the \textit{open mapping theorem}, 
that describes a remarkable property of continuous linear operators
acting between Banach spaces.
In fact,
\textsc{Baire}\textit{'s category theorem} coming from point-set topology 
leads to this fundamental theorem in {\FA}:

\begin{theorem}\label{OPEN_MAPPING_THEOREM}
If $\banX$ and $\banY$ are Banach spaces
and the linear operator $\linA\:\banX \xto \banY$ is continuous and surjective,
then $\linA$ is {open}.
\end{theorem}

\paragraph*{} %- closed graph theorem
The \textit{closed graph theorem},
which is equivalent to the {open mapping theorem},
establishes the relationship between continuous and closed operators
under our specific assumption on the domain of definition:

\begin{theorem}\label{CLOSED_GRAPH_THEOREM}
If $\topX, \topY$ are Banach spaces
and the domain of definition $\domain{\linA}$ of a closed linear operator
$\linA\:\domain{\linA} \xsubset \banX \xto \banY$ is closed,
then $\linA$ is continuous.
\end{theorem}

%\input{PROOF_CLOSED_GRAPH_THEOREM}

\paragraph*{}
This theorem especially implies
that if $\linA$ is everywhere defined (so $\domain{\linA} \eq \banX$),
then closeness and continuity of $\linA$ are equivalent to each other.

\paragraph*{} %- uniform boundedness principle
A collection $\famF$ of linear operators between normed spaces $\banX, \banY$
is called \textit{pointwise bounded},
if for each $x \xeof \banX$ there is $c_x > 0$
such that
$\norm{\linA x} \leq c_x \norm{x}$ for all $\linA \xeof \famF$.
The following statement,
called the \textit{uniform boundedness principle},
can be proved either by theorem \ref{OPEN_MAPPING_THEOREM}
or by theorem \ref{CLOSED_GRAPH_THEOREM}.

\begin{theorem}\label{UNIFORM_BOUNDEDNESS_PRINCIPLE}
Let $\banX$ be a Banach space and $\banY$ be a normed space.
Then a family $\famF$ of pointwise bounded linear operators
$\linA_{\alpha}:\banX \xto \banY$ is actually \textit{uniformly bounded},
that is,
there is a constant $c > 0$
such that $\norm{\linA}_{\alpha}\leq c$ for each $\linA_{\alpha}\xeof \famF$.
\end{theorem}

\paragraph*{}
We now turn to the problem of invertibility of operators.
A linear operator $\linA$ is said to be \textit{invertible},
if it is bijective and the inverse $\linA^{-1}$ is continuous.

\subparagraph*{}
If $\linA$ is not only continuous and surjective,
but also injective,
then theorem \ref{OPEN_MAPPING_THEOREM}
is known as the \textit{bounded inverse theorem}
or as \textsc{Banach}'s \textit{isomorphism theorem}.
It quotes that under the just mentioned conditions
the inverse (linear) operator
$\linA^{-1}:\banY \xto \banX$ is even bounded,
so $\linA$ and $\linA^{-1}$ both are homeomorphisms
between $\banX$ and $\banY$.

\subparagraph*{} %- invertible operators
Similarly,
if $\linA$ is closed and bijective,
then the operator $\linA^{-1}:\banY \xto \domain{\linA} \xsubset \banX$
is closed, too.
This assertion can be proved much easier
than the invertibility of a continuous linear operator.
In addition,
the closed graph theorem implies
that the inverse $\linA^{-1}$ of a closed linear operator $\linA$
is even continuous,
which is of outstanding importance {\wrt} well-posedness
of the linear operator equation $\linA x \eq y$.
In fact,
$\linA^{-1}$ being continuous is nothing but the statement
that the solution $x$ of this equation depends continuously
on the right hand side $y$. 

\section{Adjoint Operators}
%--------------------------
Given a linear operator $\linA$ in a Hilbert space $\hilH$
with densely defined domain of definition $\domain{\linA}$,
we want to introduce the concept of the \textit{adjoint operator} $\linA^*$.
This notation can also be defined for densely defined linear operators
between ordinary and possibly different Banach spaces,
but it is sufficient for our purposes to restrict the considerations
to the Hilbert space setting.

\subparagraph*{}
By the domain of definition $\domain{\linA^*}$ of the operator $\linA^*$
is meant the collection of elements $y \xeof \hilH$,
for which there is $y^* \xeof \hilH$ such that
\[
\scalar{\linA x}{y} = \scalar{x}{y^*}
\qfa x \xeof \domain{\linA}.
\]
The set $\domain{\linA^*}$ constitutes a linear subspace of $\hilH$.
Since $\linA$ is assumed to be densely defined,
the element $y^*$ is uniquely determined by the element $y$
(this is a consequence of \textsc{Riesz}'s representation theorem),
thus one may define the operator $\linA^*$ according to
$\linA^* y :=  y^*$ for $y \xeof \domain{\linA^*}$.
Then,
given the pair of operators $(\linA,\linA^*)$,
\[
\scalar{\linA x}{y} \= \scalar{x}{\linA^* y}
\]
holds for all $x \xeof \domain{\linA}$ and $y \xeof \domain{\linA^*}$.

\subparagraph*{}
It is easy to verify that the operator $\linA^*$ is also linear.
Moreover,
if $\linB$ is an extension of $\linA$,
then $\linA^*$ is an extension of $\linB^*$,
that is,
$\linA \prec \linB$ implies $\linB^* \prec \linA^*$.
Indeed,
let $y \xeof \domain{\linB^*}$,
which means that there is $y^* \xeof \hilH$ such that
\[
\scalar{\linB x}{y} \= \scalar{\linA x}{y} \= \scalar{x}{y^*}
\qfa x \xeof \domain{\linB}\,,
\]
then it is easy to see
that $y \xeof \domain{\linA^*}$ and $\linA^* y \eq \linB^* y$.

\subparagraph*{}
From another point of view,
the notion of adjoint operator is of interest when studying
the continuity of the sesquilinear form 
$(x,y) \mapsto \scalar{\linA x}{y}$.
In fact,
the domain of definition $\domain{\linA^* }$of the adjoint operator
consists exactly of those $y \xeof \hilH$, 
for which the functional $f_y(x) := \scalar{\linA x}{y}$ is continuous.
Since $\domain{\linA}$ is supposed to be dense in $\hilH$,
the mapping $f_y$ can be extended to the whole of $\hilH$,
and by \textsc{Riesz}'s representation theorem
\ref{RIESZ_REPRESENTATION_THEOREM} there is $y^* \xeof H$
such that $f_y(x) \eq \scalar{x}{y^*}$ for all $x \xeof \hilH$.

\subparagraph*{}
We especially have
$(\linA - \lambda \1)^* \eq \linA^* - \overline{\lambda} \1$
for $\lambda \xeof \C$,
where $\1_H$ is the identity operator in $\hilH$.
From now on we use the abbreviation $\linA-\lambda$
instead of $\linA-\lambda \1_H$.
Then the domain of definition of the operator $(\linA - \lambda)^*$ is given by
\[
\domain{(\linA  - \lambda)^*} =
\domain{\linA^* - \overline{\lambda}} =
\domain{\linA^*}.
\]
Moreover,
for $\lambda \xeof \C \setminus \{0\}$ the equality
$(\lambda \linA)^* \eq \overline{\lambda} \linA^*$ is valid,
but this surely is not true in case of $\lambda \eq 0$.

\paragraph*{}
There are fundamental relationships between the kernels and ranges
of the linear operators $\linA$ and $\linA^*$,
which read as
\begin{equation}\label{ORTHOGONALITY_RELATIONS}
\overline{\range{\linA^*} } \= \kernel{\linA}^{\bot}
\quad \mbox{and} \quad
\overline{\range{\linA} } \= \kernel{\linA^*}^{\bot}.
\end{equation}
They are very useful
to investigate the possible property of surjectivity of the operator $\linA$.
Indeed,
if one can prove that $\linA^*$ is injective,
then the range of $\linA$ is dense in $\hilH$,
and if in addition $\linA$ has closed range,
for example if $\linA$ is bounded below,
then we can conclude that $\linA$ is even surjective.

\subparagraph*{}
Each adjoint operator is closed,
but it is not necessarily densely defined.
But if $\linA$ itself is closed,
then $\linA^*$ is densely defined and $A^{**} := (A^*)^* \eq A$ holds,
so in this case taking the adjoint is a self-dual operation.
Furthermore,
if $\linA$ is closable,
but not closed,
then we have $A^{**} \eq \overline{A}$,
that is,
the closure of $\linA$ can be obtained by taking the adjoint twice.

\section{Semi-Fredholm Operators}
%--------------------------------
Since the kernel $\kernel{\linA}$ and the range $\range{\linA}$
of a linear operator $\linA\:\domain{\linA} \xsubset \banX \xto \banY$
are linear subspaces of $\banX$ and $\banY$,
respectively,
it is possible to define the value $\alpha(\linA)$ and $\beta(\linA)$
to be the Hamel dimension
of $\kernel{\linA}$ and $\cokernel{\linA}$,
respectively.
Both together measure the deviation of $\linA$ from being bijective.
Then the so-called \textit{index}
\[
\chi(\linA) \eqdef
\begin{cases}[1.2]
\quad + \infty, & \alpha(\linA) = \infty\\
\quad \alpha(\linA) \, - \, \beta(\linA), &
\alpha(\linA) < \infty, \, \beta(\linA) < \infty\\
\quad - \infty, & \beta(\linA) = \infty\\
\end{cases}
\]
of $\linA$ is well-defined and takes values in the set
$\{- \infty\} \cup \N \cup \{+ \infty\}$.
In the special case,
where $\alpha(\linA) < \infty$, 
$\beta(\linA) < \infty$
and consequently $\chi(\linA) \xeof \Z$,
the operator $\linA$ is said to have \textit{finite index}.

\subparagraph*{}
We denote a densely defined linear operator $\linA$ with range space $\banY$
to be \textit{normally solvable}
if its range $\range{\linA}$ is closed in $\banY$.
A closed and normally solvable linear operator $\linA$
is said to be \textit{semi-Fredholm},
if the kernel $\kernel{\linA} $ 
or the cokernel $\cokernel{\linA} $ is \textit{finite-dimensional},
or equivalently,
if $\alpha(\linA) < \infty$ or $\beta(\linA) < \infty$.

\subparagraph*{}
We refine the concept of semi-Fredholm operator as follows:
a normally solvable linear operator
$\linA\:\domain{\linA} \xsubset \banX \xto \banY$ is called

\begin{itemize}[leftmargin=12pt]
\item[1.]
\textit{upper semi-Fredholm},
if $\alpha(\linA) < \infty$ and $\chi(\linA) < \infty$, respectively,

\item[2.]
\textit{lower semi-Fredholm},
if $\beta(\linA) < \infty$ and $\chi(\linA) > - \infty$, respectively.

\item[3.]
\textit{Fredholm},
if it is lower and upper semi-Fredholm as well, thus $\chi(\linA) \xeof \Z$.
\end{itemize}

\subsubsection*{} % operators having closed range
We continue with a further class of linear operators,
which is of importance {\wrt} solvability of equations
of the form $\linA x \eq y$,
where $x,y$ are elements of the Banach spaces $\banX$ and $\banY$.

\subparagraph*{}
A linear operator $\linA\:\domain{\linA} \xsubset \banX \xto \banY$
is said \textit{to have closed range}, 
if $\range{\linA}$ is closed in $\banY$.
Moreover,
$\linA$ said to be \textit{bounded below}
if there is $\theta > 0$ such that 
\begin{equation}\label{BOUNDED_FROM_BELOW_ESTIMATE}
\norm{\linA x} \, \geq \, \theta \norm{x} \qfa x \xeof \domain{\linA}.
\end{equation}

\subparagraph*{}
The following assertion is one of the fundamental theorems in operator theory:

\begin{theorem}\label{BOUNDED_FROM_BELOW_THEOREM}
A closed linear operator $\linA\:\domain{\linA} \xsubset \banX \xto \banY$
is bounded below
\iff
$\linA$ is injective and has closed range.
\end{theorem}

\paragraph*{}
For closed operators $\linA$ being bounded below,
the inverse operator $\linA^{-1}$
is then defined on the Banach space $\range{\linA}$
and proves to be continuous by the {closed graph theorem},
so the linear equation $\linA x \eq y$ is well-posed for $y \xeof \range{\linA}$.

\paragraph*{}
The \textit{closed range theorem} (in its Hilbert space version)
due to \textsc{S.\,Banach} claims
that the property of a closed linear operator $\linA$ having closed range
is strongly related to that property of the adjoint operator $\linA^*$.
Moreover,
it also quotes that,
concerning the orthogonality relations \eqref{ORTHOGONALITY_RELATIONS},
the closures of the ranges of $\linA$ and $\linA^*$ 
may be substituted by the ranges themselves,
supposed one of the operators $\linA$ and $\linA^*$ has closed range.

\begin{theorem}\label{CLOSED_RANGE_THEOREM}
Let $\linA$ be a closed and densely defined linear operator
in Hilbert space $\hilH$.
Then the following assertions are equivalent:

\begin{itemize}[leftmargin=12pt]
\item[1.]
$\linA$ has closed range.
\item[2.]
$\linA^*$ has closed range.
\item[3.]
$\range{\linA^*}$ is the annihilator of $\,\kernel{\linA}$, {\ie}
$\range{\linA^*} \eq \kernel{\linA} ^{\bot} $.
\item[4.]
$\range{\linA} $ is the pre-annihilator of $\,\kernel{\linA^*}$, {\ie}
$\range{\linA} \eq ^{\bot}\kernel{\linA^*} $.
\end{itemize}
\end{theorem}

% waiting for implementation: Atkinson's theorem

\paragraph*{}
We next quote a feature of semi-Fredholm operators {\wrt} the functions
$\xi \mapsto \alpha(\linA-\xi)$ and
$\xi \mapsto \beta(\linA-\xi)$,
where $\xi \xeof \C$ is close to zero.
It is known as \textit{Gohberg}'s \textit{punctured neighbourhood theorem}.

\begin{theorem}\label{PUNCTURED_NEIGHBOURHOOD_THEOREM}
Let the operator $\linA$ be upper semi-Fredholm.
Then:

\begin{itemize}[leftmargin=12pt]
\item[1.]
the mappings $\xi \mapsto \alpha(\linA-\xi) =:c_l$ and
             $\xi \mapsto \beta(\linA-\xi) =: c_u$
are constant for $\xi \neq 0$ and $\xi$ belonging to
a sufficiently small punctured environment of $0 \xeof \C$,
\item[2.]
$c_l \leq \alpha(\linA)$ and $c_u \leq \beta(\linA)$ holds,
\item[3.]
the operator $\linA - \xi$ is upper semi-Fredholm with index that of $\linA$.
\end{itemize}
\end{theorem}

\paragraph*{} %- numerical range
In the Hilbert space setting the concept of \textit{numerical range}
is useful to study the semi-Fredholm property of operators.
We define the {numerical range} of a linear operator $\linA$ acting
in a Hilbert space $\hilH$ and with domain of definition $\domain{\linA}$ by
\[
\Theta(\linA) \eqdef
\{\,\scalar{\linA \,x}{x}: x \xeof \domain{\linA}, \norm{x}=1 \,\}.
\]
It is a known fact that
$\Theta(\linA)$ is a convex set within the complex plane.
Moreover,
let $\Delta(\linA) := \C \setminus \overline{\Theta(\linA)}$
be the complement of the closure of $\Theta(\linA)$.
If $\Delta(\linA) \neq \emptyset$,
then $\Delta(\linA)$ consists of one or two connected components.
To be precise,
if this set has two components,
then the numerical range $\Theta(\linA)$ must be a vertical or horizontal strip.

\subparagraph*{}
For $\xi \xeof \Delta(\linA)$ we of course have $\scalar{\linA x}{x} \neq \xi$
for $x \xeof \domain{\linA}$ and $\norm{x}=1$.
Then 
\[
0 \, < \, \delta(\xi,\linA) \, \leq \,
\abs{\scalar{\linA x}{x} - \xi} \=
\abs{\scalar{(\linA-\xi)x}{x}} \, \leq \,
\norm{(\linA -\xi) x},
\]
where $\delta(\xi,\linA)$ is the positive distance between the point $\xi$
and the closed set $\overline{\Theta(\linA)}$.
Hence the operator $\linA -\xi$ is bounded below,
which implies that $\linA-\xi$ is semi-Fredholm with $\alpha(\linA-\xi) \eq 0$.
By this reason the set $\Delta(\linA)$ is sometimes called
the \textit{field of external regularity} of $\linA$.
By theorem \ref{PUNCTURED_NEIGHBOURHOOD_THEOREM},
the index of $\linA-\xi$ is constant
in each connected component of $\Delta(\linA)$.

\section{Resolvent Set and Spectrum}
%-----------------------------------
Throughout this section let $\banX$ be a \textit{complex} Banach space 
and $\1_X$ the identity operator in $\banX$.
Notice that each real Banach space $\banX$ can be made
into a complex Banach space $\banX_{\C}$,
and the procedure for this is denoted the \textit{complexification} of $\banX$.
It can be done as follows:
let $\banX_{\C} := \banX + i \banX \eq \{x+iy: x,y \xeof \banX\}$ and furnished
with the addition and multiplication derived from those in $\banX$.
The norm in $\banX_{\C}$ is defined according to
$\norm{x+iy} :=
\sup\,\{\norm{x} \cos \theta + \norm{y} \sin \theta: \theta \xeof [0,2\pi)\}$,
in which this space is complete.

\subparagraph*{}
\textit{Spectral theory} is a branch in operator theory
and deals with the obstacles for a linear operator
$\linA - \lambda \1:\domain{\linA} \xsubset \banX \xto \banX$ to be bijective.
Its purpose is to obtain assertions for the solvability of the equation
\[
(\linA - \lambda \1) \, x \= y
\quad \quad (x \xeof \domain{\linA}, \; y \xeof \banX)
\]
conditional upon the complex parameter $\lambda \xeof \C$.
Thus one may consider spectral theory as a far-reaching generalization
of the eigen\-value problem for complex $(d,d)$-matrices ($d \xeof \N$)
to linear operators acting between infinite-dimensional Banach spaces.
At the same time it is a useful tool in the functional analytic study
of linear evolution equations.

\subsubsection*{} % regular value and spectral value
In the sequel we shall use consistently the abbreviation $\linA - \xi$
instead of $\linA -\xi \1$.
Given a densely defined and closed linear operator
$\linA\: \domain{\linA} \xsubset \banX \xto \banX$,
each complex number $\xi \xeof \C$ is called a

\begin{itemize}[leftmargin=12pt]
\item[1.]
\textit{regular value} of $\linA$,
if $\linA - \xi: \domain{\linA} \xto \banX$ is bijective.
The set $\rho(\linA)$ of regular values is denoted
the \textit{resolvent set} or \textit{domain of regularity} of $\linA$;

\item[2.]
\textit{spectral value} of $\linA$,
if $\xi \notin \rho(\linA)$.
The set $\sigma(\linA)$ of spectral values of $\linA$ is called
the \textit{spectrum} of $\linA$,
thus $\sigma(\linA) \eq \C \setminus \rho(\linA)$.
\end{itemize}

\subparagraph*{}
If $\xi \xeof \rho(\linA)$, then the inverse operator
$
R(\linA,\xi) := (\linA - \xi)^{-1}: \banX \xto \domain{\linA}
$
is everywhere defined in $\banX$.
Due to the closeness of $\linA$,
the operators $\linA - \xi$ and $R(\linA,\xi)$ are closed, too,
and the {closed graph theorem} \ref{CLOSED_GRAPH_THEOREM}
implies that $R(\linA,\xi)$ is continuous,
that is,
$R(\linA,\xi) \xeof \BO{\banX}$.
The operator $R(\linA,\xi)$ is said to be the \textit{resolvent} of $\linA$
{\wrt} $\xi \xeof \rho(\linA)$,
and the operator-valued mapping
\begin{equation}\label{RESOLVENT_FUNCTION}
\begin{cases}[1.2]
\quad R(\linA,\cdot): \rho(\linA) \xto \BO{\banX}\,,\\
\quad \xi \mapsto (\linA - \xi)^{-1}\\
\end{cases}
\end{equation}
is denoted the \textit{resolvent function} of $\linA$.

\subparagraph*{}
The question,
whether a given complex number $\xi$ belongs to
the resolvent set or the spectrum of $\linA$,
is answered by the following criterion:
let $(\xi_n)$ be a sequence in
$\rho(\linA)$ with $\lim_{n \xto \infty} \xi_n \eq \xi$.
Then $\xi \xeof \rho(\linA)$ \iff
\[
\sup_{n \xeof \N} \norm{R(\linA,\xi_n)} < \infty.
\]
On the other hand,
if $\xi$ lies in the boundary of $\rho(\linA)$,
then $\norm{R(\linA,\xi_n)} \xto \infty$ for each sequence $\xi_n \xto \xi$.

\subparagraph*{}
Resolvents related to a given operator are not independent from each other.
In particular,
for all $\xi_1,\xi_2 \xeof \rho(\linA)$
the \textit{first resolvent equation}
\begin{equation}\label{RESOLVENT_EQUATION}
R(\linA,\xi_1) - R(\linA,\xi_2)
\=
(\xi_2-\xi_1) \, R(\linA,\xi_1) \circ R(\linA,\xi_2)
\end{equation}
holds,
which implies that
the resolvents $R(\linA,\xi_1)$ and $R(\linA,\xi_2)$ commute with each other.
Furthermore,
if $\xi_0 \xeof \rho(\linA)$,
then we have $\xi \xeof \rho(\linA)$ for all $\xi \xeof \C$ satisfying
\[
\abs{\xi-\xi_0 } < \norm{R(\linA,\xi_0)}^{-1},
\]
which implies that $\rho(\linA)$ is open in $\C$.
For values $\xi$ in this disk
one may represent the resolvent as a power series
with coefficients in the Banach algebra $\BO{\banX}$:
\[
R(\linA,\xi) \=
\sum_{n=0}^{\infty} \, R(\linA,\xi_0)^{n+1} \, (\xi_0-\xi)^n.
\]
In fact,
the mapping $R(\linA,\cdot)$ as defined in \eqref{RESOLVENT_FUNCTION}
is locally holomorphic in $\rho(\linA)$.

\paragraph*{}
We now turn to the properties of the spectrum of $\linA$.
First of all,
the set $\sigma(\linA)$ is closed,
since $\rho(\linA)$  is open.
If $\linA$ is closable but not closed,
it may happen
that the spectrum consists of the entire set $\C$.
For a closed linear operator $\linA$,
however,
the resolvent set is non-empty,
so $\sigma(\linA)$ cannot be the whole of $\C$.
It is also possible that $\sigma(\linA)$ degenerates to the empty set.
But in the case,
where $\linA$ is everywhere defined and \textit{bounded},
the spectrum of $\linA$ is non-empty.

\subsubsection*{} % subdivision of the spectrum
We now consider the spectral values of a closed linear operator $\linA$
and study possible reasons,
which for $\lambda \xeof \sigma(\linA)$
prevent the operator $\linA - \lambda$ from being bijective.

\begin{itemize}[leftmargin=12pt]
\item[1.]
Recall from linear algebra
that a complex number $\lambda$ is called an \textit{eigenvalue}
of a linear operator $\linA$,
if there is an element $x \xeof \banX$ such that
$x \neq 0$ and $\linA x \eq \lambda x$.
In this case the element $x$ is called an \textit{eigenvector} of $\linA$
{\wrt} $\lambda$,
and all these eigenvectors form a linear subspace of $\banX$,
called the \textit{eigenspace} of $\linA$ for $\lambda$.
The \textit{point spectrum} of $\linA$ is defined to be
the set $\sigma_p(\linA)$ of all eigenvalues of $\linA$.
Clearly,
for $\lambda \xeof \sigma_p(\linA)$ the operator
$\linA - \lambda$ is not injective,
since its kernel is non-trivial.

If $\linA - \lambda$ is injective,
but not surjective,
then the inverse operator may be defined
on the range of $\linA - \lambda$.
We distinguish between two cases:

\item[2a.]
if $\range{\linA - \lambda}$ is dense in $\banX$,
then $\lambda$ is said to be an element of the set $\sigma_c(\linA)$,
called the \textit{continuous spectrum} of $\linA$,
otherwise

\item[2b.]
the number $\lambda$ belongs to the \textit{residual spectrum}
$\sigma_r(\linA)$.
There is an interesting identity between the spectrum
of a closed and densely defined linear operator $\linA$
and that of its transposed $\linA^t$, namely
$\sigma_r(\linA) \eq \sigma_p(\linA^t)$ (see \cite{EN}, p.\,161).
\end{itemize}

\subparagraph*{}
In summary,
for each densely defined operator
$\linA\:\domain{\linA} \xsubset \banX \xto \banX$
the set of complex numbers forms a disjoint decomposition
of the resolvent set and the different kinds of spectra, namely
\[
\C \= \rho    (\linA) \, \sqcup
        \, \sigma_p(\linA) \, \sqcup
        \, \sigma_c(\linA) \, \sqcup
        \, \sigma_r(\linA).
\]

\subparagraph*{}
The \textit{approximate point spectrum} $\sigma_{ap}(\linA)$ is
defined to be the set of $\lambda \xeof \C$,
for which $\linA-\lambda$ is not injective or $\range{\linA-\lambda}$
is not closed in $\banX$.
The elements of $\sigma_{ap}(\linA)$ are called
\textit{approximate eigenvalues} of $\linA$.
In fact,
any number $\lambda$ belongs to the set $\sigma_{ap}(\linA)$
\iff
there is a sequence $(x_n) \xsubset \domain{\linA}$ such that
$\norm{x_n} \eq 1$ for all $n$ and
$\lim_{n \xto \infty} \norm{\linA x_n - \lambda x_n} \eq 0$.
For a closed linear operator $\linA$,
the boundary $\partial \sigma(\linA)$ of the spectrum 
forms a subset of the approximate point spectrum of $\linA$.

\subparagraph*{}
The approximate point spectrum generalizes
the point spectrum $\sigma_p(\linA)$,
but it can be empty only in the case,
where $\sigma(\linA) \eq \emptyset$ or $\sigma(\linA) \eq \C$.

\paragraph*{}
It is also possible to breakup the spectrum into several parts
by using the notion of semi-Fredholm operator.
The \textit{essential spectrum} $\sigma_{ess}(\linA)$
of a closed operator $\linA$ consists of those points $\lambda$
that are either accumulation points of $\sigma(\linA)$
or isolated eigenvalues withe infinite-dimensional eigenspaces.
The essential spectrum $\sigma_{ess}(\linA)$ forms
a closed subset of $\sigma(\linA)$.

\subparagraph*{}
Semi-Fredholm operators can also be used to characterize
the essential spectrum $\sigma_{ess}(\linA)$
of a closed operator $\linA$:
a number $\lambda \xeof \C$ belongs to $\sigma_{ess}(\linA)$
\iff $\linA - \lambda$ is not semi-Fredholm.
The main property of the essential spectrum is given by the fact
that it does not change under compact perturbations of $\linA$.

\subparagraph*{}
On the other hand,
the relative complement
$
\sigma_d(\linA) := \sigma(\linA) \setminus \sigma_{ess}(\linA)
$
is called the \textit{discrete spectrum} of $\linA$
and purely contains eigenvalues of finite multiplicity.

\paragraph*{}
For fixed $\xi \xeof \C$,
the spectrum of the resolvent $(\linA - \xi)^{-1}$ 
is related to the spectrum of $\linA$ itself by the identity
\[
\sigma(R(\linA,\xi)) \setminus \{0\} \=
\{(\lambda-\xi)^{-1}: \lambda \xeof \sigma(\linA)\},
\]
known as \textit{spectral mapping theorems} for resolvents,
see \cite{EN}, pp.\,161-162.
Similar identities are also valid for sets $\sigma_p(\linA)$,
$\sigma_{ap}(\linA)$ and $\sigma_r(\linA)$:
\[
\sigma_p(R(\linA,\xi)) \setminus \{0\} \=
\{(\lambda-\xi)^{-1}: \lambda \xeof \sigma_p(\linA)\}.
\]
\[
\sigma_{ap}(R(\linA,\xi)) \setminus \{0\} \=
\{(\lambda-\xi)^{-1}: \lambda \xeof \sigma_{ap}(\linA)\}.
\]
\[
\sigma_r(R(\linA,\xi)) \setminus \{0\} \=
\{(\lambda-\xi)^{-1}: \lambda \xeof \sigma_r(\linA)\}.
\]

\subparagraph*{} %- spectral radius
If $\linA$ is bounded,
we may define its \textit{spectral radius} by
$r(\linA) := \sup \, \{\lambda \xeof \sigma(\linA)\}$.
It can be shown that $r(\linA) \leq \norm{\linA}$,
and since each resolvent $R(\linA,\xi)$ is bounded,
we have
\begin{equation}\label{RESOLVENT_ESTIMATE}
\norm{R(\linA,\xi)} \, \geq \,
r(R(\linA,\xi))     \, \geq \,
\delta(\xi,\sigma(\linA))^{-1}
\end{equation}
for any $\xi \xeof \rho(\linA)$,
where $\delta(\xi,\sigma(\linA))$ is the minimal distance between
$\xi$ and the spectrum of $\linA$.
Note that inequality \eqref{RESOLVENT_ESTIMATE} is a consequence
of the spectral mapping theorem.

\paragraph*{}
We finally describe the basic sets of spectral theory 
for adjoint operators in Hilbert space. 
Let $\linA\:\domain{\linA} \xsubset \hilH \xto \hilH$
be closed and densely defined.
Then the following identities hold:

\begin{itemize}[leftmargin=12pt]
\item[1.] $\rho(\linA^*) =
          \{\, \overline{\xi} \xeof \C: \xi \xeof \rho(\linA) \,\}$.
\item[2.] $\sigma(\linA^*) =
          \{\, \overline{\lambda} \xeof \C: \lambda \xeof \sigma(\linA) \,\}$.
\item[3.] $\sigma_c(\linA^*) =
          \{\, \overline{\lambda} \xeof \C: \lambda \xeof \sigma_c(\linA) \,\}$.
\item[4.] If $\lambda \xeof \sigma_p(\linA) \setminus \sigma_c(\linA)$,
          then $\overline{\lambda} \xeof (\sigma_p(\linA^*) \setminus
               \sigma_c(\linA^*)) \cup \sigma_r(\linA^*)$.
\item[5.] If $\lambda \xeof \sigma_r(\linA)$,
          then $\overline{\lambda} \xeof \sigma_p(\linA^*) \setminus
                                       \sigma_c(\linA^*)$.
\end{itemize}

\paragraph*{} %- numerical range and spectrum
We finally clarify the relations
between the spectrum and the numerical range
rsp. the external field of regularity
of operators acting in a Hilbert space.

\begin{theorem}
Let $\hilH$ be a Hilbert space and
$\linA\:\domain{\linA} \xsubset \hilH \xto \hilH$ be a linear operator.
Then the following assertions are true:

\begin{itemize}[leftmargin=15pt,itemsep=5pt]
\item[(a)]
If $\linA$ is bounded,
then $\sigma(\linA) \subseteq \overline{\Theta(\linA)}$.

\item[(b)]
$\sigma_p(\linA) \subseteq \Theta(\linA)$.

\item[(c)]
$\sigma_c(\linA) \subseteq \overline{\Theta(\linA)}$.

\item[(d)]
If $\zeta \xeof \Delta(\linA)$,
then
$\norm{R(\linA,\zeta)} \, \leq \, {\mbox{dist}(\zeta,\Theta(\linA)}^{-1}$.

\item[(e)]
If $U \subseteq \Delta(\linA)$ is connected
and $U \cap \rho(\linA) \neq \emptyset$,
then $U \subseteq \rho(\linA)$.
\end{itemize}
\end{theorem}

\subparagraph*{}
The last assertion of this theorem is nothing but to say
that if a subset $U$ of $\Delta(\linA)$ contains at least one element
of $\rho(\linA)$,
then all elements of $U$ belong to $\rho(\linA)$.

\section{Compact Linear Operators and Fredholm Operators}
%--------------------------------------------------------
A linear operator $\compA\:\banX \xto \banY$ between Banach spaces
$\banX$ and $\banY$ is called \textit{compact},
if the range of each bounded set under $\compA$ is relative compact
{\wrt} the norm topology in $\banY$,
or equivalently,
if for every bounded sequence $\{x_n\}$ in $\banX$
the sequence $\{\compA \, x_n\}$ in $\banY$ has a convergent subsequence.

\subparagraph*{}
For example,
a linear operator $\compA\:\banX \xto \banY$ is compact,
if at least one of the spaces $\banX$ and $\banY$ is finite-dimensional.
In case of $\banX \eq \banY$,
the identity operator $\1_X \: \banX \xto \banX$ is compact
\iff $\banX$ is finite-dimensional.

\subsubsection*{} % properties of compact linear operators
Compact operators have the following general properties:

\begin{itemize}[leftmargin=12pt]
\item[1.]
they are closed under taking finite linear combinations of them,
and the limit of a norm convergent sequence of compact operators
is compact again;
\item[2.]
the composition $\linA_1 \circ \linA_2$ of bounded linear operators
is compact,
if at least one of them is so;
\item[3.]
each compact linear operator is continuous, 
which can easily be shown by contradiction;
\item[4.]
$\CO{\banX}$, the set of compact operators $\compA\:\banX \xto \banX$
constitutes a \textit{closed} ideal within the operator algebra $\BO{\banX} $;
\item[5.]
$\CO{\banX}$ is complete and thus a Banach algebra in its own right.
\end{itemize}

\paragraph*{}
A linear operator $F \: \banX \xto \banY$ is said to have \textit{finite rank},
if the range $\range{F}$ is finite-dimensional,
which of course implies compactness of $F$.
Assume that $(F_n)$ is a sequence of finite-rank operators
converging to a further operator $\linA$ {\wrt} the operator norm,
then $\linA$ also proves to be compact.
Thus compact operators are part of
the closure of the set $\mathscr{F}(\banX,\banY)$ of finite-rank operators,
which itself lies within the set $\mathscr{C}(X,Y)$ of compact operators.

\subparagraph*{}
Reversely,
in the setting of Hilbert spaces 
each compact operator $\compA$ appears in such a way,
that is,
$\compA$ is the norm-limit of a convergent sequence of finite-rank operators.
But this assertion may not be true
for compact operators acting between general Banach spaces.

\paragraph*{}
With respect to the question of compactness of a linear operator and its
transposed,
\textsc{Schauder}\textit{'s theorem} is valid:

\begin{theorem}
A continuous linear operator $\linA\:\banX \xto \banY$
between Banach spaces $\banX$ and $\banY$ is compact
\iff $\linA^t$ is compact.
\end{theorem}

\subparagraph*{}
For a proof of this theorem, see for example \cite{Yo}, pp. 282-283.

\paragraph*{}
We now turn to the spectral properties of compact linear operators.
Historically,
the spectral theory of such an operator was developed
by \textsc{F.\,Riesz} and \textsc{J.\,Schauder}.
The essential results are known
as the \textsc{Riesz-Schauder} \textit{spectral theorem},
which resembles those of $(d,d)$-matrices
acting in the finite-dimensional spaces $\C^d$:

\begin{theorem}\label{RIESZ-SCHAUDER_SPECTRAL_THEOREM}
Let $\banX$ be a complex Banach space and
$\compA\:\banX \xto \banX$ be compact.
Then the following assertions hold:

\begin{itemize}[leftmargin=12pt]
\item[1.]
The spectrum $\sigma(\compA) \xsubset \C$ is mostly countable.

\item[2.]
$\xi \eq 0$ is the only possible limit point of $\sigma(\compA)$.

\item[3.]
If $\banX$ is infinite-dimensional, then $0 \xeof \sigma(\compA)$.

\item[4.]
Each spectral value $\lambda \neq 0$ is an eigenvalue of $\compA$
with finite multi\-plicity.

\item[5.]
Every $0 \neq \lambda \xeof \sigma(\compA)$
is a pole of the resolvent function
$R(\compA,\, \cdot\, ):\rho(\compA) \xto \BO{\banX}$.
\end{itemize}
\end{theorem}

\paragraph*{} %- Fredholm operators
In the second part of this section we shall deal with the class of
Fredholm operators.
A continuous linear operator $\linA\:\banX \xto \banY$
is called \textit{Fredholm},
if it is lower and upper semi-Fredholm at the same time,
or equivalently,
if $\chi(\linA) \xeof \Z$.
We collect a couple of non-trivial properties of Fredholm operators:

\begin{itemize}[leftmargin=15pt]
\item[(a)]
$\linA^*$ is Fredholm \iff $\linA$ is so.

\item[(b)]
If one of $\linA$ or $\linA^*$ is Fredholm,
then the dimensions of the kernels and ranges of $\linA$ and $\linA^*$
are related to each other by
$\alpha(\linA^*) \eq \beta(\linA)$ and
$\beta(\linA^*) \eq \alpha(\linA)$,
which implies $\chi(\linA^*) \eq - \chi(\linA)$.

\item[(c)]
Continuous Fredholm operators acting in $\banX$ and
belonging to one and the same connected component set
within the operator algebra $\BO{\banX}$ have identical index.
Thus,
if $\linA$ is Fredholm and $\linB$ can be reached by a continuous path
with starting point $\linA$,
then $\linB$ is Fredholm, too,
and $\chi(\linB) \eq \chi(\linA)$.
\end{itemize}

\paragraph*{} %- left and right parametrix
A Fredholm operator can be represented by means of
a \textit{left} and \textit{right parametrix},
respectively.

\subparagraph*{}
Preliminary,
if $\linA\:\banX \xto \banY$ is a continuous linear operator,
then a linear mapping $P_l:\banY \xto \banX$ is called a \textit{left parametrix}
if $\linA \circ P_l \eq \1_X + \compA_l$
for some compact operator $\compA_l \xeof \BO{\banX}$.
In a similar way,
a linear mapping $P_r:\banY \xto \banY$
is said to be a \textit{right parametrix} of $\linA$,
if $P_r \circ \linA \eq \1_Y + \compA_r$
for some compact operator $\compA_r \xeof \BO{\banY}$.
Every Fredholm operator owns a left and right parametrix,
where the involved compact operators $\compA_l,\,\compA_r$
are of \textit{finite-rank},
that is,
they have finite-dimensional ranges.

\paragraph*{} %- Fredholm's alternative theorem
We next cite the \textsc{Fredholm} \textit{alternative theorem}
as the outstanding result in Fredholm theory
{\wrt} the solvability of the linear equation $\linA x \eq y$.
Preliminary,
the \textit{rank-nullity theorem} known from linear algebra claims that
if a linear space $\vecV$ is finite-dimensional,
then a linear mapping $\linA\:\vecV \xto \vecV$ is surjective
iff it is injective.
\textsc{Fredholm}'s {alternative theorem}
is a considerable generalization to general Banach spaces
and Fredholm operators with vanishing index:

\begin{theorem}\label{FREDHOLM_ALTERNATIVE_THEOREM}
Let $\banX,\,\banY$ be Banach spaces and
the continuous linear operator $\linA\:\banX \xto \banY$ be Fredholm
with index $\chi(\linA) \eq 0$.
According to the solvability of the linear equation $\linA \, x \eq y$,
exactly one of the following cases occurs:

\begin{itemize}[leftmargin=15pt]
\item[(a)]
$\linA \, x \eq y$ is uniquely solvable for all $y \xeof \banX$,
so $\linA$ is bijective,

\item[(b)]
$\linA \, x \eq y$ is not for all $y \xeof \banX$ solvable,
and in this case the operator $\linA$ is neither injective nor surjective.
If the equation is solvable for an element $y_0 \xeof \banX$,
then the set of solutions forms a finite-dimensional subspace of $\banX$.
In addition,
the equation $\linA \, x \eq y$ owns a solution $x \xeof \banX$
\iff
$y \xeof \,^{\bot}\kernel{\linA^*}$.
\end{itemize}
\end{theorem}

\paragraph*{} %- usage of Fredholm operators
A great portion of {Fredholm theory} deals
with the study of linear operators of type $\1+\compA$, 
where $\compA:\banX \xto \banY$ is compact.
Such operators are indeed Fredholm,
to which \textsc{Fredholm}'s alternative theorem can be applied.
On the other hand,
a continuous semi-Fredholm operator can be represented by means of
a \textit{left} and \textit{right parametrix},
respectively.
Given a continuous linear operator $\linA\:\banX \xto \banY$,
a linear mapping $P_l:\banY \xto \banY$ is called a \textit{left parametrix}
if $\linA \circ P_l \eq \1_X + \compA_l$
for some compact operator $\compA_l \xeof \BO{\banX}$.
Similarly,
a linear mapping $P_r:\banY \xto \banY$
is said to be a \textit{right parametrix} of $\linA$,
if $P_r \circ \linA \eq \1_Y + \compA_r$
for some compact operator $\compA_r \xeof \BO{\banY}$.

\subparagraph*{} %- equivalence of linear equations
We emphasize the usage of Fredholm theory by the following consideration,
which is helpful for applications in partial differential equations.
Consider the linear equation $\linA x \eq y$,
where $\banX$ is Banach
and $\linA\:\banX \xto \banX$ is a continuous linear operator.
We assume that there is $\zeta \xeof \C$
such that
$\linA - \zeta$ is invertible and the mapping $(\linA - \zeta)^{-1}$ is compact.
Then we obtain
\[
\linA \,x \= (\linA -\zeta)\,x + \zeta \,x \= y,
\]
and applying the mapping $(\linA -\zeta)^{-1}$ from the left yields
\[
x \= \zeta (\linA + \zeta)^{-1}\,x \= \tilde{y} := (\linA -\zeta)^{-1}\,y
\]
or equivalently $x + \tilde{\compA} \,x \eq \tilde{y}$,
where $\tilde{\compA} := - \zeta (\linA-\zeta)^{-1}$ is compact.

\subparagraph*{}
The equations $\linA \,x \eq y$ and $(\1 + \tilde{\compA})\,x \eq \tilde{y}$
are called \textit{equivalent to each other},
since $x$ is a solution of the first one
\iff
it is a solution of the second one.
By theorem \textsc{Atkinson}'s theorem \ref{ATKINSON_THEOREM},
the operator $\1 + \tilde{\compA}$ is Fredholm with vanishing index,
so theorem \ref{FREDHOLM_ALTERNATIVE_THEOREM} is applicable
and provides assertions also for the solution of the equation $\linA x \eq y$.
We summarize this statement to the following theorem:

\begin{theorem}\label{FREDHOLM_ALTERNATIVE_THEOREM_II}
Let $\linA\:\domain{\linA} \xsubset \banX \xto \banX$ be an operator,
for which there is $\zeta \xeof \C$ such that $\linA -\zeta$ is invertible.
Then the equation $\linA x \eq y$ is equivalent to a equation
of the form $(\1+\compA)x=y$,
and \textsc{Fredholm}'s alternative theorem for $\linA x=y$ is valid.
\end{theorem}

\paragraph*{} %- Fredholm property of closed linear operators
We have discussed the (semi-)Fredholm property of linear operators
under the basic assumption
that they are continuous and everywhere defined in their domain spaces.
If $\linA$ is only closed instead of continuous
however,
we define the possible Fredholm property by regarding
$\linA$ as a continuous linear operator $\linA\:\domain{\linA} \xto \banY$,
where $\domain{\linA}$ is endowed
with the graph norm $\norm{\,\cdot\,}_{\linA}$ defined by \eqref{GRAPH_NORM}.

\paragraph*{} %- solvability of linear equations
We close this section with a collection of concepts regarding solvability
of linear equations of the form $\linA x \eq y$,
where $\linA$ is a closed linear operator acting between Banach spaces
$\banX$ and $\banY$.
Following \cite{Kr} and \cite{Br2},
the equation $\linA x \eq y$ is called

\begin{itemize}[leftmargin=15pt]
\item[(a)]
\textit{uniquely solvable} (\cite{Kr}),
if $\linA$ is injective;

\item[(b)]
\textit{everywhere solvable} (\cite{Kr}),
if $\linA$ is surjective;

\item[(c)]
\textit{solvable} (\cite{Br2}),
if it is uniquely and everywhere solvable;

\item[(d)]
\textit{well-posed},
if it is solvable and $\linA^{-1}$ is continuous;

\item[(e)]
\textit{completely solvable} (\cite{Br2}),
if it is solvable and $\linA^{-1}$ is compact;

\item[(f)]
\textit{densely solvable} (\cite{Kr}),
if $\overline{\range{\linA}} \eq \banY$;

\item[(g)]
\textit{normally solvable} (\cite{Kr}),
if $\overline{\range{\linA}} \eq \range{\linA}$;

\item[(h)]
\textit{correctly solvable} (\cite{Kr}),
if $\linA$ is bounded below;

\item[(i)]
\textit{semi-solvable} (\cite{Br2}), if $\linA$ is Fredholm;

\item[(j)]
\textit{almost solvable} or \textit{regularly solvable} (\cite{Br2}),
if $\linA$ is semi-solvable and $\chi(\linA) \eq 0$.
\end{itemize}

\subparagraph*{}
To the end,
remember again the important result based on the closed graph theorem:
if $\linA$ is continuous and closed,
then (a)\,+\,(b) implies (d).

\section{Sesquilinear Forms and Associated Operators}
%----------------------------------------------------
Remember that a \textit{sesquilinear form} (\textit{bilinear form})
in a complex Hilbert space $\vecV$ is a mapping
$\sql:\vecV \xtimes \vecV \xto \C$
which is linear in the first and antilinear in the second argument.
Such forms are in one-to-one relationship with mappings
$\mathbb{A}:\vecV \xto \vecV^*$,
where $\vecV^*$ is the conjugate space
of continuous and (anti)linear functionals $\phi:\vecV \xto \F$.
In fact,
we define $\mathbb{A}(x) := f_x \xeof \vecV^*$
with $f_x(y) := \sql(x,y)$ for $y \xeof \vecV$,
and the form $\sql$ can be recovered from the operator $\mathbb{A}$
in a similar way.
One often says that $\mathbb{A}$ is the \textit{representation operator}
associated with the form $\sql$.

\subparagraph*{}
Moreover,
any form $\sql$ is called \textit{bounded},
if there is $\Theta \geq 0$ such that
\[
\abs{\sql(x,y)} \, \leq \, \Theta \norm{x}{y} \quad \quad (x,y \xeof \vecV).
\]
It can be shown by the \textit{uniform boundedness principle},
that $\sql$ is continuous \iff $\sql$ is bounded.
A fundamental representation theorem for continuous forms claims
that there is a bounded linear operator $B \xeof \BO{\vecV}$ such that
$\sql(x,y) \eq \scalar{B x}{y}$ for all $x,y \xeof \vecV$.

\paragraph*{}
The concept of \textit{coerciveness} is of outstanding meaning in what follows.
Let $\sql$ be a continuous sesquilinear form in Hilbert space $\vecV$.
Then $\sql$ is called \textit{coercive},
if there is $\theta > 0$ such that
\[
\re \sql(x,x) \, \geq \, \theta \norm{x}^2 \quad \quad (x \xeof \vecV).
\]
This property especially implies
\[
\norm{B x}\norm{x} \geq
\abs{\scalar{B x}{x}} =
\abs{\sql(x,x)} \geq
\re \sql(x,x) \geq
\theta \norm{x}^2 \quad (x \xeof \vecV),
\]
hence the operator $B$ is bounded below with lower bound $\theta$.

\paragraph*{} %- Lax-Milgram theorem
The \textsc{Lax-Milgram} \textit{lemma} is one of the key theorems
in applied functional analysis
and ensures the well-posedness of the variational problem
\begin{equation}\label{VARIATIONAL_EQUATION}
\sql(x,y) \= \phi(y) \qfa y \xeof \vecV.
\end{equation}
for any right-hand side $\phi \xeof \vecV^*$.

\begin{theorem}\label{LAX-MILGRAM_THEOREM}
If the form $\sql:\vecV \xtimes \vecV \xto \C$ is coercive,
then problem \eqref{VARIATIONAL_EQUATION} has exactly
one solution $x_{\phi} \xeof \vecV$ for each $\phi \xeof \vecV^*$.
In addition,
$\norm{\,x_{\phi}\,} \leq \theta^{-1} \norm{\phi}$ holds.
\end{theorem}

\subparagraph*{}
In other words:
if $\sql$ is coercive,
then the representation operator
$\mathbb{A}: \vecV \xto \vecV^*$ is bjective,
the inverse operator $\mathbb{A}^{-1}:\vecV^* \xto \vecV$ is bounded
and the solution $x_{\phi}$ of \eqref{VARIATIONAL_EQUATION}
continuously depends on $\phi$.

\begin{comment}
\paragraph*{}
\Proof:
Since $\linB$ is bounded below,
theorem \ref{BOUNDED_FROM_BELOW_THEOREM_I} implies
that $\linB$ is injective and $\range{\linB}$ is closed.
It is still to show that $\linB$ is surjective.
If $z \xeof \range{\linB}^{\bot}$,
then for $\linB z \xeof \range{\linB} $ we have
\[
0 \=
\scalar{\linB z}{z}
\=
\sql(z,z)
\, \geq \,
\theta \norm{z}^2
\, \geq \, 0,
\]
which implies $z \eq 0$ and therefore $\overline{\range{\linB} } \eq \vecV$.
But $\range{\linB}$ is closed,
so $\range{\linB} \eq \vecV$.
Finally,
the inequality $\norm{x_f} \leq \theta^{-1} \norm{f}$
follows from the fact that $\norm{\linB^{-1}} \leq \theta^{-1}$.
\qed
\end{comment}

\paragraph*{} %- pairs of Hilbert spaces
The presented theory of coercive sesquilinear forms in a Hilbert space $\vecV$
is often not sufficent for applications.
Thus it is indicated to study coercive forms in the setting of
\textit{pairs} $(\vecV,\hilH)$ of Hilbert spaces,
where the space $\vecV$ is continuously and densely embedded
into a second space $\hilH$.
More precisely,
we assume that

\begin{itemize}[leftmargin=40pt]
\item[(T-1)]
there is a linear mapping $\iota_1:\vecV \xto \hilH$
such that $\overline{\range{\iota_1} } \eq \hilH$;

\item[(T-2)]
there is $c>0$ such that
$\norm{\iota_1 v}_{\hilH} \leq c \norm{v}_{\vecV}$ for all $v \xeof \vecV$.
\end{itemize}
By means of the embedding $\iota_1$,
one can also think of $\vecV$ as a dense linear subspace
within the larger space $\hilH$.

\subparagraph*{} %- triples of Hilbert spaces
Moreover,
each Hilbert pair $(\vecV,\hilH)$ gives rise to introduce
the \textit{transposed pair} $(\hilH^*,\vecV^*)$,
where the conjugate dual space $\hilH^*$
is continuously and densely embedded into
the conjugate dual space $\vecV^*$
by a further embedding $\iota_2:\hilH^* \xto \vecV^*$.
Then,
by identifying the spaces $\hilH$ and $\hilH^{\times}$
via the \textsc{Riesz} map,
we combine both Hilbert pairs to obtain the triple
\begin{equation}\label{GELFAND_TRIPLE}
\vecV \, \overset{\iota_1}{\hookrightarrow} \,
\hilH \, \overset{\iota_2}{\hookrightarrow} \, \vecV^*,
\end{equation}
sometimes called a \textit{Gelfand triple},
an \textit{evolution triple} or a \textit{rigged Hilbert space}.
The center space $\hilH$ is called
the \textit{pivot space} of \eqref{GELFAND_TRIPLE}.

\subparagraph*{}
Triples of Hilbert spaces own a remarkable property.
Let $(\,\cdot\,,\,\cdot\,)$ and $\scalar{\,\cdot\,}{\,\cdot\,}$
be the inner products in $\vecV$ and $\hilH$, respectively,
then for fixed chosen $x \xeof \hilH$
the action of the antilinear linear functional
$f_x := \iota_2 \, x \xeof \vecV^*$ to any element $v \xeof \vecV$
can be represented by the inner product of $x$ and $v$ in $\hilH$,
{\ie} $f_x(v) \eq \scalar{x}{\iota_1 \, v}$ for $v \xeof \vecV$.

\paragraph*{} %- quasi-coercive forms
It is quite natural to examine sesquilinear forms
also {\wrt} pairs $(\vecV,\hilH)$ of Hilbert spaces.
Moreover,
in applications it is often the case
that sesquilinear forms are merely continuous but not coercive.
In the setting of a Hilbert pair $(\vecV,\hilH)$,
however,
we may try to make a form $\sql$ coercive
by adding an additional term,
which is a multiplicity of the inner product in $\hilH$.

\subparagraph*{}
More precisely,
if $\sql$ is given on $\vecV$ and the slightly modified form
\[
\sql_{\lambda}(u,v) \eqdef
\sql(u,v) \, + \, \lambda \scalar{\iota_1 u}{\iota_1 v}\, \quad
(u,v \xeof \vecV)
\]
is coercive for some $\lambda \xeof \R$,
that is,
there is $\theta > 0$ such that
\begin{equation}\label{QUASI_COERCIVE}
\re \sql(u,u)  \, + \,
\lambda \norm{\iota_1u}_H^2 \, \geq \, \theta \norm{u}_V^2
\quad (u \xeof \vecV),
\end{equation}
then $\sql$ is called
\textit{quasi-coercive},
\textit{$\hilH$-elliptic}
or simply a \textit{\GARDING} \textit{form}.
It should be clear that
if $\sql_{\lambda_0}$ is coercive for some $\lambda_0 \xeof \R$,
then $\sql_{\lambda}$ is so for $\lambda > \lambda_0$.

\paragraph*{} %- operators associated with coercive forms
Quasi-coercive sesqilinear forms can be used
to derive densely defined and closed linear operators
in the pivot space $\hilH$ of the underlying Hilbert triple.

\subparagraph*{}
Let $\sql$ be a sesquilinear form defined on a Hilbert space $\vecV$,
which is the left hand side of a Hilbert triple $(\vecV,\hilH,\vecV^*)$
with continuous embeddings
$\iota_1:\vecV \xto \hilH$ and
$\iota_2:\hilH \xto \vecV^*$.
We introduce a linear operator $\linA_{\sql}$ with domain space $\hilH$
and range space $\hilH$
by restriction of the form operator $\repA:\vecV \xto \vecV^*$ of $\sql$
to those elements $u \xeof \vecV$,
for which $\repA \, u \xeof \hilH$.
Indeed,
the domain $\domain{\linA_{\sql}}$ is defined to be the subspace
of those elements $u \xeof \vecV$,
for which $\repA \, u \xeof \iota_2(\hilH) \xsubset \vecV^*$ holds,
briefly denoted
\[
\domain{\linA_{\sql}} \eqdef \{ \, u \xeof \vecV: \repA \, u \xeof \hilH \, \}.
\]
Clearly,
this set can be regarded as a subset of $\hilH$
by means of the embedding $\iota_1$.

\subparagraph*{}
Then for $u \xeof \domain{A}$ there exists $y \xeof \hilH$
such that $\iota_2 y \eq \repA \, u$.
We set $x := \iota_1 u$ and $\linA_{\sql} x := y$
for $x \xeof \domain{\linA_{\sql}}$
and denote $\linA_{\sql}$ the \textit{operator in $\hilH$
associated with $(\vecV,\sql)$},
or shortly,
the \textit{variational operator} derived from $\sql$.
Another phrase is to say that
the operator $\linA_{\sql}$ is the \textit{part of $\repA$ in $\hilH$}.

\subparagraph*{} %- main theorem on quasi-coercive forms
The following theorem collects a couple of nice properties 
of the operator $\linA_{\sql}$ associated with a quasi-coercive form $\sql$,
which besides the \textsc{Lax-Milgram} lemma
is our second important theorem in this section:

\begin{theorem}\label{QUASI-COERCIVE_FORMS_AND_ASSOCIATED_OPERATORS}
Let $(\vecV,\hilH,\vecV^*)$ be a Hilbert triple
and the sesquilinear form $\sql$ be quasi-coercive on $\vecV$
with parameter $\lambda_0 \xeof \R$.
Then the following assertions are true:

\begin{itemize}[leftmargin=12pt]
\item[1.]
$\linA_{\sql}$ is closed.

\item[2.]
$\linA_{\sql}$ is densely defined in $\hilH$.

\item[3.]
If the embedding of $\vecV$ into $\hilH$ is \textit{compact},
then the operator $\linA_{\sql} + \lambda$ has a compact inverse
$(\linA_{\sql} +\lambda)^{-1}$ for each $\lambda > \lambda_0$.
\end{itemize}
\end{theorem}

\begin{comment}
\subparagraph*{}
\Proof:
To prove that $\linA_{\sql}$ is closed,
let $(x_n)$ be a sequence in $\domain{\linA_{\sql}}$
with limit element $x \xeof \hilH$
and $(\linA_{\sql} x_n$ be the sequence mapped by $\linA_{\sql}$
with limit element $y \xeof \hilH$.
It is to show that $x \xeof \domain{\linA_{\sql}}$ and $\linA_{\sql} x=y$.
Let $\abs{x} $ and $\norm{x}$ be the norms in $\vecV$ and $\hilH$,
respectively.
Clearly,
we have
\[
\norm{(\linA_{\sql}+\lambda) x_n - y - \lambda x} \xto 0 
\quad \for n \xto \infty.
\]
In the following we shall use the fact
that for $\lambda > \lambda_0$ the operators
$\linA_{\sql}+\lambda$ and $(\linA_{\sql}+\lambda)^{-1}$
forms linear homeomorphism between the normed linear spaces
$(\hilH,\norm{\,\cdot\,})$
and $(\domain{\linA_{\sql}},\abs{\,\cdot\,} )$,
so the latter one forms a closed subspace
within the normed space $(\vecV,\abs{\,\cdot\,} )$.
By this reason,
\[
x_n = (\linA_{\sql}+\lambda)^{-1} (\linA_{\sql}+\lambda) x_n
\to 
(\linA_{\sql}+\lambda)^{-1} (y + \lambda x)
\quad \for n \xto \infty
\]
{\wrt} the norm $\abs{\,\cdot\,} $ in $\domain{\linA_{\sql}}$.
Since $x$ belongs to $\hilH$ and $\domain{\linA_{\sql}}$ is closed,
the element $(\linA_{\sql}+\lambda)^{-1} x$
belongs to $\domain{\linA_{\sql}}$ with $\linA_{\sql} x \eq y$.

\subparagraph*{}
To prove that $\domain{\linA_{\sql}}$ is dense in $\hilH$,
suppose that $y \xeof \hilH$ fulfills $\scalar{x}{y} \eq 0$
for all $x \xeof \domain{\linA_{\sql}}$.
It is to show that $y=\0$ in this case.
Since $\linA_{\sql}+\lambda$ is surjective,
there must exist $x_0 \xeof \domain{\linA_{\sql}}$ such that
$(\linA_{\sql}+\lambda) x_0 \eq y$.
Then
\[
0 \=
\scalar{\linA_{\sql} x_0}{x} \, + \, \scalar{\lambda x_0}{x} \=
\sql(x_0,x) \, + \, \lambda \scalar{x_0}{x}.
\]
Especially for $x \eq x_0$ the last equation leads to
\[
0 \=
\scalar{\linA_{\sql} x_0}{x_0} \, + \, \scalar{\lambda x_0}{x_0} \,\geq\,
\theta \norm{x_0}^2
\]
for some $\theta > 0$,
due to to quasi-coercivity of $\sql$,
hence $x_0 \eq \0$ and finally $y=\0$.

\subparagraph*{}
Finally,
if $\iota_1:\vecV \xto \hilH$ is compact,
then $\iota_1 \circ (\linA_{\sql}+ \lambda)^{-1}:\hilH \xto \hilH$
is compact,
too,
since the composition of a compact linear operator
and a continuous linear operator proves to be compact.
\qed.
\end{comment}

\section{Linear Operator-Differential Equations}
%-----------------------------------------------
A \textit{vector-valued} function $u\:[a,b] \xto \banX$
with values in a Banach space $\banX$
is said to be \textit{strongly or classical differentiable} in $t_0 \xeof (a,b)$,
if there is an element $u'(t_0) \xeof \banX$ such that
\[
\lim_{t \xto t_0} \,
\norm{ \, \frac{u(t)-u(t_0)}{t-t_0}\, - u'(t_0) \, } \= 0.
\]
The function $u$ is called \textit{differentiable in} $(a,b)$, 
if $u$ is so in each $t \xeof (a,b)$.
The space $C^1(a,b;\banX)$ consists of functions
having continuous classical derivatives in $(a,b)$.
Moreover,
the spaces $C^k(a,b;\banX)$ ($k \xeof \N$) are defined recursively as usual,
whereas the space $C^{\infty}(a,b;\banX)$ is 
the intersection of all the spaces $C^k(a,b;\banX)$ for $k \xeof \N_0$.

\subsubsection*{} % linear evolution equations
By a linear \textit{evolution equation} in a Banach space $\banX$,
also  called a \textit{linear operator-differential equation},
we mean an ordinary differential equation of the form
\begin{equation}\label{OPERATOR_DIFFERENTIAL_EQUATION}
\quad u'(t) \= \linA \, u(t) + f(t) \quad (0 < t < \T),
\end{equation}
where $\linA$ is a densely defined linear operator
$\linA\:\domain{\linA} \xsubset \banX \xto \banX$
and \eqref{OPERATOR_DIFFERENTIAL_EQUATION} must be valid
for each $t \xeof (0,\T)$.
We emphasize
that the differential equation only makes sense
if $u'(t) \xeof \domain{\linA}$ for every $t \xeof (0,\T]$.
We shall call \ref{OPERATOR_DIFFERENTIAL_EQUATION} \textit{homogeneous},
if $f \eq 0$.

\subsubsection*{} % homogeneous Cauchy problem
The \textit{homogeneous Cauchy problem}
is to find a solution for every $t \xeof [0,\T]$
of the initial-value problem
\begin{equation}\label{HOMOGENEOUS_CAUCHY_PROBLEM}
\begin{cases}[1.2]
\quad u'(t) \= \linA \, u(t) \\
\quad u(0)=u_0 \xeof \banX. \\
\end{cases}
\end{equation}

\subparagraph*{}
A function $u \xeof C([0,\T],\domain{\linA}) \, \cap \, C^1((0,\T],\banX)$
is called a \textit{classical solution} of \eqref{HOMOGENEOUS_CAUCHY_PROBLEM},
if it satisfies the initial condition and
the differential equation in each point $t \xeof (0,\T]$.
Moreover,
a continuous function $u\:[0,\T) \xto \banX$ is called a \textit{mild solution}
of \eqref{HOMOGENEOUS_CAUCHY_PROBLEM} if
\[
\mbox{(1) } \, \tilde{u}(t) := \int_{\,0}^t u(s) \, ds \, \xeof \domain{\linA}
\quad \mbox{ and } \quad
\mbox{(2) } \, \linA \, \tilde{u}(t) = u(t)-u_0
\]
holds for all $0 < t \leq \T$.
One is led to the notion of mild solution in a natural way
by integrating both sides of the equation $u'=\linA \, u$,
which then becomes into
\[
\int_{\,0}^t u' ds \=
u(t)-u(0) \=
\int_{\,0}^t \linA \, u(s) \, ds.
\]
If $\linA$ is assumed to be closed,
then the action of $\linA$ and performing the integration commutes, that is,
\[
\int_{\,0}^t \linA \, u(s) \, ds  \= \linA \, \int_{\,0}^t u(s) \, ds,
\]
and to the end we achieve the equation
\[
u(t) \= u(0) \, + \,  \linA \, \int_{\,0}^t u(s) \,ds
\quad \quad (t \xeof (0,\T]).
\]
Notice that each classical solution is a mild solution,
and reversely a mild solution $u$ is classical
\iff
$u$ is continuously differentiable in $(0,\T]$,
thus classical and mild solutions of homogeneous Cauchy problems
only differ {\wrt} this regularity condition.

\subparagraph*{}
The homogeneous Cauchy problem \eqref{HOMOGENEOUS_CAUCHY_PROBLEM} 
is called \textit{correctly posed},
if there is a unique classical solution
for any initial value $x_0 \xeof \domain{\linA}$.
In addition,
it is denoted \textit{well-posed},
if the solution depends continuously on the initial values,
or equivalently,
if for each $\T > 0$ there is $c_{\T} > 0$ independent of $u_0$ such that
\[
\norm{u(t)} \leq c_{\T} \norm{u_0} \qfa t \xeof (0,\T]
\]
and all $u_0 \xeof \banX$,
for which a classical solution exists
(the set of such elements may be the domain of definition $\domain{\linA}$
of $\linA$ or any other dense subspace in $\banX$).

\paragraph*{}
The abstract Cauchy problem of \textit{second order}
is to solve the initial-value problem
\begin{equation}\label{HOMOGENEOUS_SECOND_ORDER_CAUCHY_PROBLEM}
\begin{cases}[1.2]
\quad u''(t) \, + \, \linA \, u(t) \= 0\\
\quad u(0) \= u_0, \; u'(0) \= u_1,\\
\end{cases}
\end{equation}
where $u_0 \xeof \domain{\linA}$,
but $u_1$ may be an element of an enlarged subspace within $\banX$.

\subparagraph*{}
Similar to the first order problem,
a function $u \xeof C([0,\T],\domain{\linA}) \, \cap \, C^2((0,\T],\banX)$
is called a \textit{classical solution}
of \eqref{HOMOGENEOUS_SECOND_ORDER_CAUCHY_PROBLEM},
if it satisfies the initial conditions and
the differential equation in each point $t \xeof (0,\T]$.

\subsubsection*{} %- semigroups
The notion of \textit{semigroup} for bounded linear operators
has already been introduced by \textsc{J.\,A.\,S\'{e}guier} in 1904.
Later some applications to partial differential equations were detected,
but the first systematic and comprehensive treatement of semigroup theory
was done in the pioneering monograph \cite{HP}
by \textsc{E.\,Hille} and \textsc{R.\,Phillips}.
Nowadays semigroup theory provides a powerful method
in solving linear operator-differential equations
with a multiplicity of applications to linear partial differential equations.

\subparagraph*{} %- strongly continuous semigroups of bounded linear operators
Given a Banach space $\banX$,
a \textit{C$_0$-semigroup} or
a \textit{strongly continuous semigroup of bounded linear operators}
is an operator-valued mapping $\linS:[0,\infty) \xto \BO{\banX}$,
shortly denoted $\{\linS(t)\}_{t \geq 0}$
and distinguished by the properties

\begin{itemize}[leftmargin=28pt]
\item[(G-1)]
$\linS(t{+}s) \eq \linS(t) {\circ} \linS(s)$ for all $t,s \geq 0$
(the \textit{semigroup property});

\item[(G-2)]
$\lim_{\,t \downarrow 0} \, \linS(t)\,x \eq x$ for all $x \xeof X$.
\end{itemize}

\subparagraph*{}
It follows from (G-2),
which is nothing but strong continuity from the right-hand side in $t \eq 0$,
that $\linS(0)$ is the identity mapping in $\banX$.
Notice also
that the conditions (G-1) and (G-2) together
imply continuity from the right of the semigroup even in every $t > 0$.
In case of $\banX=\R$,
semigroups share both properties
with the family $\{\cal{E}(t)\}_{\, t \geq 0}$ of exponential functions
$\cal{E}(t): \lambda \mapsto e^{\lambda t}$, ($\lambda \xeof \R$),
so this family serves as an elementary example
of a strongly continuous semigroup.

\subparagraph*{} %- classification of strongly continuous semigroups
We next perform a classification on the totality of strongly
continuous semigroups.
For each such semigroup $\{\linS(t)\}_{t \geq 0}$
there are constants $M \geq 0$ and $\omega \xeof \R$
such that
\begin{equation}\label{SEMIGROUP_GROWTH}
\norm{\linS(t)} \, \leq \, M e^{\omega t}
\quad \for t > 0.
\end{equation}
In fact,
let $M \geq 1$ such that $\norm{S(t)} \leq M$ for all $s \xeof [0,1]$.
For $t > 0$ there is $s \xeof [0,1)$ and $n \xeof N$ such that $t \eq n+s$.
Then $\norm{\linS(n)}\leq M \norm{S(n-1)}$ for $n > 1$ and
\[
\norm{\linS(t)} = \norm{\linS(n+s)} = \norm{\linS(n) \circ \linS(s)} \leq
\]
\[
\norm{\linS(n)} \cdot \norm{\linS(s)} \leq M^{n+1} = M e^{n \cdot \ln{M}}.
\]

\paragraph*{} %- orbit maps
By the \textit{orbit map} or \textit{trajectory} $\eta_x$ of a semigroup
associated with an element $x \xeof \banX$
we mean the vector-valued mapping $\eta_x:[\,0,\infty) \xto \banX$
given by $\eta_x(t) := S(t) \, x$.
Due to condition (G-1),
strong continuity of $\{\linS(t)\}_{t \geq 0}$
implies continuity of the functions $\eta_x$.
But $\eta_x$ is even differentiable under a very mild assumption:
if $\{\linS(t)\}_{t \geq 0}$ is a C$_0$-semigroup,
then for each $x \xeof \banX$
the mapping $\eta_x$ is differentiable in $(0,\infty)$
\iff
$\eta_x$ is right differentiable in $t \eq 0$.

\paragraph*{} %- infinitesimal generators
If $\{\linS(t)\}_{t \geq 0}$ is a strongly continuous semigroup,
we define its \textit{infinitesimal generator}
(briefly denoted the \textit{generator} of the semigroup)
as a mapping $\linA$ in $\banX$ by means of the graph
\[
\Gamma_{\linA} \eqdef
\{\,(x,y) \xeof X \xtimes X: \;
y \= \lim_{\,h \downarrow 0} \, \frac{1}{h}(\linS(h)x-x) \exists \}.
\]
Hence $\domain{\linA}$ of $\linA$ consists of those elements $x \xeof \banX$,
for which the orbit map $\eta_x$ is differentiable from the right in $t \eq 0$,
that is,
$\domain{\linA} \eq \{\,x \xeof \banX: \eta_x \mbox{ is differentiable}\,\}$.
It is easy to show that
$\linA$ is linear,
so it is a linear operator acting in $\banX$.

\subparagraph*{}
The meaning of strongly continuous semigroups is confirmed
by the solution theory of homogeneous Cauchy problems
governed by linear operators,
which at the same time are infinitesimal generators of $C_0$-semigroups.
Indeed,
the \textit{main theorem of semigroup theory} reads as follows:

\begin{theorem}\label{MAIN_THEOREM_OF_SEMIGROUP_THEORY}
Let $\linA$ be a closed and densely defined linear operator
acting in a Banach space $\banX$.
Then the following assertions are equivalent to each other:

\begin{itemize}[leftmargin=12pt]
\item[1.]
The homogeneous Cauchy problem \eqref{HOMOGENEOUS_CAUCHY_PROBLEM}
is well-posed.
\item[2.]
The operator $-\linA$ generates a strongly continuous semigroup
of type $(\omega,M)$.
\item[3.]
The resolvent set of $\linA$ fulfills $(0,\infty) \xsubset \rho(\linA)$
and the resolvent function $R(\linA,\xi)$ satisfies
\begin{equation}\label{HILLE-YOSIDA_CONDITION}
\norm{R(\linA,\xi)^n} = \norm{(\linA - \xi)^{-n}}
\leq M (\xi - \omega)^{-n} \quad (n \xeof \N, \xi >0).
\end{equation}
\end{itemize}
\end{theorem}

\subparagraph*{}
The equivalence of the both last assertions is the content
of the famous \textsc{Feller-Miyadera-Hille-Yosida} \textit{theorem}.
Since the number of conditions on the operator $\linA$ is infinite,
this theorem has limited worth {\wrt} applications.
However,
the original theorem
for generators of \textit{contraction semigroups}
due to \textsc{Hille} and \textsc{Yosida} in 1948 quotes
that in this case only the first inequality
($n=1$, with $\omega \eq 0$ and $M \eq 1$)
of \eqref{HILLE-YOSIDA_CONDITION} has to be verified:

\begin{theorem}\label{HILLE-YOSIDA_THEOREM}
The linear operator $-\linA$ acting in a Banach space $\banX$
is the generator of a C$_0$-semigroup of \textit{contractions}
\iff
\begin{itemize}[leftmargin=12pt]
\item[1.]
$\linA$ is densely defined and closed,
\item[2.]
the interval $(-\infty,0)$ belongs to the resolvent set $\rho(\linA)$,
\item[3.]
$\norm{\xi \, R(\linA,\xi) } \leq 1$ for all $\xi < 0$.
\end{itemize}
\end{theorem}

\subparagraph*{}
In fact,
if $-\linA$ owns the mentioned properties,
then the semigroup $\{\linS(t)\}_{t \geq 0}$ may be constructed by
\[
\linS(t) \, x \eqdef
\lim_{\,n \xto \infty} (\1 + \frac{t}{n} \linA)^{-n} \,x 
\fe t > 0 \and x \xeof \banX.
\]

\paragraph*{} %- accretive operators
Our next aim is to describe the spectral properties
in theorem \ref{HILLE-YOSIDA_THEOREM} by equivalent conditions,
which leads to the concept of \textit{accretive} operator.

\subparagraph*{}
Let $\hilH$ be a complex Hilbert space,
then a linear operator $\linA$ with domain of definition $\domain{\linA}$
acting in $\hilH$ is called \textit{accretive}, if
\[
\re \, \scalar{\linA \,x}{x} \, \geq 0
\qfa x \xeof \domain{\linA},
\]
or equivalently, 
if the numerical range $\Theta(\linA)$ is part of the right complex half-plane.
Moreover,
a linear operator $\linA$ is said to be \textit{dissipative}, 
if $-\linA$ is accretive.

\subparagraph*{}
We first notice that if $\linA$ is accretive (dissipative),
then $\lambda \linA$ is so for each $\lambda > 0$.
Furthermore,
since 
\[
\norm{ \linA \,x + \lambda x}^2 - \norm{\lambda x}^2
\, \geq \,
2 \lambda \; \re \, \scalar{\linA \,x}{x}
\qfa x \xeof \domain{\linA},
\]
the operator $\linA$ is accretive \iff
\begin{equation}\label{ACCRETIVE_PROPERTY}
\norm{(\linA + \lambda) x} \, \geq \, \norm{\lambda x}
\qfa \lambda > 0 \and x \xeof \domain{\linA}.
\end{equation}
This property can be used to define coercivity and dissipativity
for operators also in normed linear spaces,
in which no inner product is defined.

\subparagraph*{}
The inequality \eqref{ACCRETIVE_PROPERTY} says that
for $\lambda > 0$
the operator $\linA + \lambda$ is bounded below and thus injective.
In addition,
if $\linA$ is closed,
then theorem \ref{BOUNDED_FROM_BELOW_THEOREM} implies that
$\linA$ has closed range,
and by the closed graph theorem the inverse operator
$(\linA + \lambda)^{-1}: \range{\linA + \lambda} \xto \domain{\linA}$
is bounded with norm
$
\norm{(\linA + \lambda)^{-1}} \leq  1 / {\lambda} ,
$
so it even forms a contraction for $\lambda \geq 1$.

\paragraph*{} %- m-accretive operators
In the following
we assume that for $\linA$ being accretive there is $\lambda_0 > 0$
such that $\linA + \lambda_0$ is surjective.
A linear operator having this property is said
to fulfill the \textit{range condition}.
Then there is an open interval $\J$ containing $\lambda_0$
such that $\linA + \lambda$ is surjective for all $\lambda \xeof \J$.
It can be proved that in this case
$\linA + \lambda$ is even surjective for all $\lambda > 0$.

\subparagraph*{}
A linear operator $\linA$ is called \textit{maximal accretive}
or shortly \textit{m-accretive}
if it is accretive and the range condition is fulfilled.
The following list contains some properties of such an operator $\linA$:

\begin{itemize}[leftmargin=12pt]
\item[1.] $\linA$ is densely defined in $\hilH$;
\item[2.] $\linA$ is closed;
\item[3.] $(-\infty,0) \xsubset \rho(\linA)$;
\item[4.] $\linA + \lambda$ is surjective for all $\lambda > 0$.
\end{itemize}

\subparagraph*{} %- characterization of m-accretive operators
Our next theorem characterizes the class of m-accretive operators:

\begin{theorem}
For $\linA$ being accretive,
the following assertions are equivalent to each other:

\begin{itemize}[leftmargin=12pt]
\item[1.]
$\linA$ is m-accretive;
\item[2.]
$\linA$ is accretive and there is no proper accretive extension of $\linA$;
\item[3.]
The resolvent set $\rho(\linA)$ contains an element of the open left half-plane;
\item[4.]
the adjoint operator $\linA^*$ is accretive, too.
\end{itemize}
\end{theorem}

\subparagraph*{}
In fact,
if $\linA$ owns the mentioned properties,
the semigroup $\{\linS(t)\}_{t \geq 0}$ may be constructed by
\begin{equation}\label{YOSIDA_APPROXIMATION}
\linS(t) \, x \eqdef \lim_{\,n \xto \infty} (\1 + \frac{t}{n} \linA)^{-n} \,x 
\fe t > 0 \and x \xeof \banX.
\end{equation}

\paragraph*{} %- Lumer-Phillips theorem
Recall that a m-accretive operator $\linA$ acting in a Hilbert space
distinguishes by the fact
that $\re \scalar{\linA \,x}{x} \geq 0$ for $x \xeof \domain{\linA}$
and $\xi \xeof \rho(\linA)$ for $\xi < 0$.
Using the concept of maximal accretivity,
the \textit{Lumer-Phillips} \textit{theorem} classifies the operators
being the infinitesimal generators of a contraction semigroup
without using any spectral theoretical concepts:

\begin{theorem}\label{LUMER-PHILLIPS_THEOREM}
Let $\linA$ be a closed linear operator in Hilbert space.
Then $-\linA$ generates a contractive C$_0$-semigroup
\iff
$\,\linA$ is m-accretive.
\end{theorem}

\subparagraph*{}
Notice that both conditions causes the operator $\linA$
to be densely defined and closed.
This is also true
if the operator $\linA$ is acting in a reflexive Banach space.
But in the setting of an arbitrary Banach space,
theorem \ref{LUMER-PHILLIPS_THEOREM} only holds for those operators,
which in addition are densely defined and closed.

\newpage
\thispagestyle{empty}
%------------------------------------------------------------------------------%
%                                   PART II                                    %
%------------------------------------------------------------------------------%
\part{Abstract Theory of Self-Adjoint Linear Operators}
%------------------------------------------------------------------------------%
In this second part we introduce and discuss some classes
of linear operators in Hilbert space,
which rely on the presence of the inner product and cannot be defined 
in arbitrary Banach spaces.

\subparagraph*{} %- normal operators
The concept of adjoint operator gives rise to a further class
of linear operators in Hilbert space,
which are called \textit{normal}.
More precisely,
a densely defined linear operator $\linA$
with domain of definition $\domain{\linA}$
and acting in $\hilH$ is denoted \textit{normal},
if
\[
\linA \circ \linA^* \= \linA^* \circ \linA.
\]
This definition of course implies the identity
of the sets $\domain{\linA}$ and $\domain{\linA^*}$.
Notice that a normal operator is closed (since $\linA^*$ is so),
$\domain{\linA} \eq \domain{\linA^*}$,
the normed space $(\domain{\linA^*},\norm{\,\cdot\,}_{\linA^*}$ is complete
and the norm $\norm{\,\cdot\,}_{\linA}$ is equivalent to the norm
$\norm{\,\cdot\,}_{\linA^*}$.

\paragraph*{}
The subsequent statements are equivalent to each other:
\begin{itemize}[leftmargin=9mm]
\item[1.] $\linA$ is normal;
\item[2.] $\linA^*$ is normal;
\item[3.] $\domain{\linA} \eq \domain{\linA^*}$ and
          $\norm{\linA \,x} \eq \norm{\linA^* x}$
          for each element $x \xeof \domain{\linA}$.
\end{itemize}

\subparagraph*{}
The class of normal and bounded operators in Hilbert space
includes a couple of important subclasses.
In summary,
an \textit{unitary} linear operator distinguishes 
by the property $\linA \circ \linA^* =1$,
a \textit{symmetric operator} owns the property $\linA \prec \linA^*$,
and a \textit{self-adjoint operator} is defined by $\linA \eq \linA^*$.
Especially the concept of self-adjoint operator have led
to a well-elaborated theory with many applications in quantum mechanics
and other topics of mathematical physics.

\section{Symmetric and Half-Bounded Operators}
%---------------------------------------------
Let $\hilH$ be a real or complex Hilbert space.
A linear operator acting in $\hilH$ and with domain of definition
$\domain{\linA} \xsubset \hilH$ is called \textit{symmetric}, if
\begin{equation}\label{SYMMETRY_PROPERTY}
\scalar{\linA x}{y} \= \scalar{x}{\linA y} \qfa x,y \xeof \domain{\linA}.
\end{equation}

\paragraph*{}
We first consider the symmetry property of operators 
being continuous and everywhere defined.
Let $\BO{\hilH}$ be the Banach space of \textit{continuous} linear operators
$\linA\:\hilH \xto \hilH$ furnished with the \textit{operator norm}
as defined in \eqref{OPERATOR_NORM}.
Then, if $\linA$ is symmetric,
we have
\[
- \norm{\linA} \, \leq \, \scalar{\linA x}{x} \, \leq \, \norm{\linA},
\]
for $x \xeof \hilH$ with $\norm{x}=1$,
or equivalently,
\[
\Theta(\linA) :=
\{\scalar{\linA x}{x}: x \xeof \domain{\linA}, \norm{x}=1\}
\subset [-\norm{\linA}, \norm{\linA}\,],
\]
where $\Theta(\linA)$ is the \textit{numerical range} of $\linA$.
Hence one may introduce in the class of bounded symmetric operators
acting in $\hilH$ the partial order
\[
\linA_1 \, \geq \linA_2
\;\; \mbox{\iff} \;\;
\scalar{(\linA_1 - \linA_2) \, x}{x} \, \geq \, 0
\quad \quad (x \xeof \hilH),
\]
which implies that
each symmetric linear operator $\linA \xeof \BO{\hilH}$ can be enclosed
by the simple multiplication operators $\linA_{-} := - \norm{\linA} \1$
and $\linA_+ := \norm{\linA} \1$.

\subparagraph*{}
On the other hand,
the \textsc{Hellinger-Toeplitz} \textit{theorem} quotes
that if $\linA$ is symmetric and $\domain{\linA}$ is the whole of $\hilH$,
then $\linA$ must be bounded,
which to the end is a consequence of the {closed graph theorem}.
In fact,
let $\linA$ be symmetric with $\domain{\linA} \eq \hilH$,
then for each convergent sequence $(x_n)$ tending to $0$
and with $(\linA x_n)$ tending to $x \xeof \hilH$ 
we have
\[ 
                    \scalar{x}        {x}       \= 
\lim_{n \xto \infty} \scalar{\linA x_n}{x}       \=
\lim_{n \xto \infty} \scalar{x_n}      {\linA x} \= 
                    \scalar{0}        {\linA x} \= 0,
\]
which implies $x \eq 0$, so $\linA$ is closed 
and even continuous by theorem \ref{CLOSED_GRAPH_THEOREM}.
Hence we may always assume
that the domain of definition $\domain{\linA}$
of a non-continuous symmetric linear operator $\linA$
is a \textit{proper} subspace of $\hilH$.

\paragraph*{}
The symmetry of a linear operator with domain of definition dense in $\hilH$
can be characterized by other properties than \eqref{SYMMETRY_PROPERTY}:

\begin{theorem}
Let $\linA$ be a densely defined linear operator with domain of definition
$\domain{\linA}$ in Hilbert space $\hilH$.
Then the following statements are equivalent:

\begin{itemize}[leftmargin=12pt]
\item[1.]
$\linA$ is symmetric;
\item[2.]
$\scalar{\linA x}{x} \xeof \R$ for all $x \xeof \domain{\linA}$;
\item[3.]
$\linA \prec \linA^*$.
\end{itemize}
\end{theorem}

\subparagraph*{}
The second property is of course also meaningful for those symmetric operators,
which are not densely defined in the entire space $\hilH$.

\begin{comment}
\paragraph*{}
We want to prove the equivalence of 1. and 2.

\subparagraph*{}
If $\linA$ is symmetric,
then $\scalar{\linA x}{x} \eq \scalar{x}{\linA x} \eq \overline{\scalar{\linA x}{x}}$,
that is, this expression is real-valued.
Reversely,
if $\linA$ is densely defined in $\hilH$ and
$\scalar{\linA x}{x} \xeof \R$ for all $x \xeof \domain{\linA}$ is true,
then $\linA$ must be symmetric.
This can be shown as follows: let $x,y \xeof \domain{\linA}$, then
\[
\begin{cases}[1.3]
\quad
\R \ni \scalar{\linA (x+y)}{x+y}   & = \; \scalar{\linA x}{x}   +   \scalar{\linA x}{y} +    \scalar{\linA y}{x} + \scalar{\linA y}{y},\\
\quad
\R \ni \scalar{\linA (ix+y)}{ix+y} & = \; - \scalar{\linA x}{x} + i \scalar{\linA x}{y} + -i \scalar{\linA y}{x} + \scalar{\linA y}{y},\\
\end{cases}
\]
which implies $\im \scalar{\linA x}{y} \eq - \im \scalar{\linA y}{x}$
and $\re \scalar{\linA x}{y} \eq  \re \scalar{\linA y}{x}$.
Then
\[
\scalar{\linA x}{y} = \re \scalar{\linA x}{y} + i \im \scalar{\linA x}{y} =
\re \scalar{\linA y}{x} - i \im \scalar{\linA y}{x} =
\overline{\scalar{\linA y}{x}} = \scalar{x}{\linA y},
\]
which shows that $\linA$ is symmetric.
\end{comment}

\paragraph*{}
Since the adjoint $\linA^*$ of a symmetric operator $\linA$
forms a closed extension of $\linA$,
the operator $\linA$ itself is closable.

\subparagraph*{}
Moreover,
the closure $\overline{\linA}$ turns out to be symmetric.
In fact,
to show that $\overline{\linA}$ is symmetric,
let $(x_n)$ and $(y_n)$ be sequences in $\domain{\linA}$
with $x_n \xto x \xeof \domain{\overline{\linA}}$,
$\linA x_n \xto \overline{\linA} x$
and $y_n \xto y \xeof \domain{\overline{\linA}}$,
$\linA y_n \xto \overline{\linA} y$.
Since $\scalar{\linA x_n}{y_n} \eq \scalar{x_n}{\linA y_n}$
for all $n \xeof \N$,
taking the limit for $n \xto \infty$ shows
that the closure of $\linA$ must be symmetric, too.

\paragraph*{}
Recall that a number $\lambda \xeof \C$ is called an \textit{eigenvalue}
of a linear operator $\linA$,
if there is $x \xeof \domain{\linA}$,
called the \textit{eigenvector} of $\linA$ {\wrt} $\lambda$,
such that $\linA x \eq \lambda x$.
The following properties of eigenvalues and eigenvectors
of a \textit{symmetric} linear operator $\linA$ are easy to prove:

\begin{itemize}[leftmargin=12pt]
\item[1.]
each eigenvalue $\lambda$ of $\linA$ is real;

\item[2.]
eigenvectors $x_1,x_2$
belonging to different eigenvalues $\lambda_1 \neq \lambda_2$
are \textit{orthogonal} to each other,
that is,
$\scalar{x_1}{x_2} \eq 0$.
\end{itemize}

\subsubsection*{} % half-bounded operators
We finally consider an important subclass of symmetric operators.
Any symmetric operator $\linA$
with domain of definition $\domain{\linA} \xsubset \hilH$
is called \textit{half-bounded}
if there is $\theta \xeof \R$ such that
\begin{equation}\label{HALF_BOUNDED_PROPERTY}
\scalar{\linA x}{x} \, \geq \, \theta \norm{x}^2
\qfa x \xeof \domain{\linA}.
\end{equation}
The operator $\linA$ is called \textit{positive} (\textit{strictly positive}),
if $\theta \eq 0$ ($\theta > 0$) can be chosen.
In case of $\theta > 0$,
however,
\eqref{HALF_BOUNDED_PROPERTY} together with the Cauchy-Schwarz inequality
\eqref{CAUCHY_SCHWARZ_INEQUALITY} implies
that $\linA$ is also bounded below,
that is,
\begin{equation}\label{BOUNDED_FROM_BELOW_PROPERTY}
\norm{\linA x} \, \geq \, \theta \norm{x} \qfa x \xeof \domain{\linA},
\end{equation}
so $\linA$ is injective by theorem \ref{BOUNDED_FROM_BELOW_THEOREM}.
Furthermore,
this theorem also claims
that if $\linA$ is \textit{closed}, 
then $\range{\linA}$ is closed in $\hilH$.
Thus, to prove bijectivity of $\linA$ it suffices to show that 
$\kernel{\linA^*} \eq \{0\}$,
since 
$\hilH \eq \kernel{\linA^*}^{\bot} \eq \overline{\range{\linA}} \eq \range{\linA}$
in this case.

\section{The Concept of Self-Adjointness}
%----------------------------------------
\textit{Self-adjointness} of a linear operator
acting in a Hilbert space is an extremely desirable property, 
since it makes the operator accessible for spectral analysis,
which is the study of the spectrum and the resolvent function,
respectively.

\subparagraph*{}
To motivate the definition of this concept,
let $\linA$ be a symmetric linear operator in Hilbert space $\hilH$.
Then every other symmetric linear operator $\linS$ with $\linA \prec \linS$
is called a \textit{symmetric extension} of $\linA$.
Notice 
that $\linA \prec \linS$ implies $\linS^* \prec \linA^*$,
which is even true if $\linA$ and $\linS$ are non-symmetric.
Furthermore,
let $\linS_1,\ldots,\linS_k$ be a finite sequence
of symmetric extensions of $\linA$,
where $\linS_{i+1}$ extends $\linS_i$ for $i=1,\ldots,k-1$.
Since $\linS_k \prec \linS_k^*$ and
$\linS_{i+1}^* \prec \linS_i^*$ for $i=1,\ldots,k-1$,
we finally obtain by concatenating all these extensions the chain
\[
\linA \, \prec \, \linS_1
\, \prec \,
\linS_2
\, \prec \,
\ldots
\, \prec \,
\linS_k
\, \prec \,
\linS_k^*
\, \prec \,
\ldots
\, \prec \,
\linS_2^*
\, \prec \,
\linS_1^*
\, \prec \,
\linA^*.
\]
Thus one may hope to find a \textit{maximal} symmetric extension
$\linS_{sa}$ of $\linA$ such that
\[
\linA \, \prec \, \linS_{sa} \= \linS_{sa}^* \, \prec \, \linA^*,
\]
which means that
$\linS_{sa}$ cannot be extended to a larger symmetric operator.
This expectation lead us to introduce the following notion:
$\linA$ is called \textit{self-adjoint},
if $\linA^* \prec \linA$,
or equivalently,
if $\linA \eq \linA^*$.
For \textit{bounded} and \textit{everywhere defined} linear operators,
however,
there is of course no difference between symmetry and self-adjointness.

\subparagraph*{}
It should be clear that a self-adjoint operator $\linA$ is closed, 
since $\linA^*$ is so.
On the other hand,
the closure of a symmetric linear operator is not necessarily self-adjoint.
Hence we denote a closable and symmetric operator $\linA$
\textit{essentially self-adjoint},
if $\overline{\linA}$ is self-adjoint.
It should be clear
that an essentially self-adjoint operator $\linA$ has exactly
one self-adjoint extension,
which is nothing but the closure of $\linA$.

\paragraph*{}
The \textit{range condition} serves as the standard criterion to verify
(essentially) self-adjointess of a symmetric linear operator $\linA$:

\begin{theorem}\label{RANGE_CONDITION}
A symmetric linear operator $\linA$
with dense domain of definition $\domain{\linA}$ in $\hilH$ is

\begin{itemize}[leftmargin=12pt,itemsep=5pt]
\item[1.]
self-adjoint \iff
$\,\range{\linA + i} \eq \range{\linA - i} \eq \hilH$,
\item[2.]
essentially self-adjoint \iff
$\overline{ \range{\linA + i} } \eq \overline{ \range{\linA - i} } \eq \hilH$.
\end{itemize}
\end{theorem}

\subparagraph*{}
Proving the first assertion,
let $\linA$ be self-adjoint. 
We make use of the fundamental relationship
between the range of $\linA-i$
and the kernel of $(\linA-i)^* \eq \linA^* +i$,
namely
\[
\overline{\range{\linA-i}} = \kernel{\linA^*+i}^{\bot}.
\]
Let $y \xeof \kernel{\linA^*+i}$,
that is, $\linA^* y \eq -i y$.
Then
\[
\scalar{y}{y} =
\im \scalar{-i y}{y} =
\im \scalar{\linA^* y}{y} =
\im \scalar{y}{\linA y} = 0,
\]
since $\scalar{y}{\linA y}$ is real-valued, which implies $y \eq 0$.
This shows that $\kernel{\linA^* +i}=\{0\}$
and $\range{\linA-i}$ is dense in $\hilH$.
But since $\linA-i$ is bounded below and thus has closed range,
we have $\range{\linA-i}=\hilH$.\\
A similar computation with $i$ replaced by $-i$ shows
that also $\range{\linA+i}=\hilH$.

\subparagraph*{}
Reversely,
let the range condition be valid.
We show self-adjointness of $\linA$ as follows:
Since $\linA \prec \linA^*$,
we have only to show $\linA^* \prec \linA$, or equivalently,
$y \xeof \domain{\linA^*}$ implies $y \xeof \domain{\linA}$.

Now,
since $(\linA-i)^*x \eq (\linA^* +i)x$ for $x \xeof \domain{\linA}$,
we have
\[
\scalar{\,(\linA-i)x}{y} = \scalar{\linA x}{y} -i\scalar{x}{y} =
\scalar{x}{\linA^*y} -\scalar{x}{-iy} = \scalar{x}{\,(\linA^* +i)y}.
\]
But since $\range{\linA+i} \eq \hilH$,
there exists $z \xeof \domain{\linA}$ such that
$
(\linA + i)z \eq (\linA^* +i)y.
$
Inserting this in the right-hand side of the last equation, we obtain
\[
\scalar{\,(\linA-i)x}{y} = \scalar{x}{(\linA + i)z} = \scalar{\,(A-i)x}{z}.
\]
But since in addition $\range{\linA-i} \eq \hilH$, we obtain
$\scalar{u}{y-z} \eq 0$ for every vector $u := (\linA-i)x \xeof \hilH$,
which implies $y=z$.
\qed

\paragraph*{}
Essentially self-adjointness of an operator can also be characterized
by means of its adjoint:
a linear operator $\linA$ is essentially self-adjoint
\iff
$\linA$ is symmetric and $\linA^*$ is symmetric,
and $\overline{\linA} \eq \linA^*$ holds in this case.

\section{Self-Adjoint Extensions of Symmetric Operators}
%-------------------------------------------------------
We review some aspects of \textsc{John\,v.\,Neumann}'s 
\textit{extension theory} for symmetric operators.
First notice that the closure of a symmetric linear operator is not
necessarily self-adjoint.
To ensure self-adjointess of a symmetric operator $\linA$,
the kernels of both operators $\linA^* \pm i$ are of special interest.
In fact,
if these are trivial,
then $\linA \pm i$ are surjective and
$\pm i$ are both regular values of $\linA$.

\subparagraph*{}
\textsc{F.\,Riesz} and \textsc{J.\,v.\,Neumann} 
have developed an extension theory,
which answers the question
in what way a symmetric and closed operator
can be extended to a self-adjoint operator,
that is,
under which conditions a symmetric and closed
but \textit{not} self-adjoint operator can be extended
in order the extension to be self-adjoint.
Moreover,
the uniqueness of such an extension if of fundamental interest.

\subsubsection*{}
Remember first
that for a densely defined and closed symmetric operator $\linA$
the relation $\domain{\linA} \xsubset \domain{\linA^*}$ holds.
But even more is valid,
since due to \textsc{von\,Neumann} 
the domain of definition $\domain{\linA^*}$ of the adjoint operator
can be represented as sum of three complementary subspaces \footnote{
\,The subspaces $\vecV_1,\vecV_2$ of $\vecV$
are called a complementary decomposition of $\vecV$
(notation: $\vecV \eq \vecV_1 \oplus \vecV_2)$,
if $\scalar{x}{y} \eq 0$ for all $x \xeof \vecV_1$, $y \xeof \vecV_2$
and $\vecV$ is generated by $\vecV_1 \cup \vecV_2$.}
:
\begin{equation}\label{NEUMANN}
\domain{\linA^*} \= \domain{\linA} \, \oplus \, \kernel{\linA^*+i}
                                   \, \oplus \, \kernel{\linA^*-i}.
\end{equation}
Thus the linear operators $\linA^*+i$ and $\linA^*-i$ are of special interest.
The numbers
\[
n_+ \eqdef \dim \kernel{\linA^*+i}\,,\quad \quad
n_- \eqdef \dim \kernel{\linA^*-i}
\]
are called the \textit{defect indices} of $\linA$,
where the value $n_{\pm} \eq \infty$ is admissible, too.
One easily derives from \eqref{NEUMANN}
that a symmetric and closed linear operator $\linA$
with defect indices $(n_+,n_-)$ can be extended to a self-adjoint operator
iff $n_+ \eq n_-$,
and that this operator is self-adjoint
iff in addition $n_+ \eq n_- \eq 0$ holds.

\paragraph*{}
To extend in case of $n_+ \eq n_-$ the symmetric linear operator $\linA$
to a self-adjoint one,
it suffices to have a isometric and linear mapping
\[
U: \kernel{\linA^*+i} \xto \kernel{\linA^*-i}.
\]
Then, if $U$ is given in such a way, we define $\linA_s x := \linA x$
for $x \xeof \domain{\linA}$ and
\[
\linA_s z \eqdef iz +-i U z \for z \xeof \kernel{\linA^* + i},
\]
which implies that $\linA_s$ is symmetric and hence forms a symmetric extension
of $\linA$.
One can show that $\linA_s \eq \linA^*$, so $\linA_s$ is even self-adjoint.
The set of all self-adjoint extensions is given by the set of isometric linear
operators $U$ as defined above.

\subsubsection*{} %- symmetric forms and associated operators
It is often not easy to find a self-adjoint extension
of a symmetric linear operator,
But it is helpful to start with the symmetric sesquilinear form
associated with a symmetric linear operator $\linA$
to derive a self-adjoint operator $\linA_{sa}$ extending $\linA$.

\subparagraph*{}
For $\sql$ being a symmetric sesquilinear form on a Hilbert space $\hilH$,
we define its \textit{adjoint form} $\sql^*$ by virtue of
\[
\sql^* (u,v) \eqdef \overline{\sql(v,u)}
\quad \quad (u,v \xeof \hilH).
\]
Clearly,
if $\sql$ is continuous,
then its adjoint $\sql^*$ is so.
Moreover,
a sesquilinear form $\sql$ is called \textit{symmetric},
if $\sql \eq \sql^*$.
A well-known example of a symmetric form is the inner product
$\scalar{\,\cdot\,}{\,\cdot\,}_H$ in a Hilbert space $\hilH$,
which differs from arbitrary symmetric sesquilinear forms
by {positive definiteness}.

\begin{theorem}
If $\sql$ is symmetric,
then $\linA_{\sql}$ is self-adjoint
and even half-bounded with lower bound $\lambda_0$.
\end{theorem}

\subsubsection*{}
In the following we want to elaborate
that each half-bounded linear operator can be extended to a self-adjoint one
by means of its associated sesquilinear form.

\paragraph*{}
We first consider a symmetric sesquilinear form
$\sql: \domain{\sql} \xtimes \domain{\sql} \xto \C$,
whose domain of definition $\domain{\sql}$ is a dense subspace in $\hilH$.
Recall that $\sql$ is said to be symmetric if
\[
\overline{\sql(y,x)} = \sql(x,y) \quad \fa x,y \xeof \domain{\sql},
\]
thus $\sql(x,x)$ is real-valued for $x \xeof \domain{\sql}$.
Just as for operators,
such a form is said to be \textit{half-bounded} or \textit{bounded below}
if
\[
\sql [x] \eqdef \sql(x,x) \, \geq \, \theta \norm{x}^2
\fa x \xeof \domain{\sql}
\]
and some $\theta \xeof \R$.
Clearly,
a half-bounded linear operator $\linA$ with domain of definition 
$\domain{\linA} \xsubset \hilH$ gives rise to a half-bounded sesquilinear form
$\sql_{\linA}$ by letting
\[
\begin{cases}[1.3]
\quad \domain{\sql_{\linA}} \eqdef \domain{\linA}\,,\\
\quad \sql_{\linA}(x,y)     \eqdef \scalar{\linA x}{y}
\quad \mbox{for } x,y \xeof \domain{\sql_{\linA}}\,.\\
\end{cases}
\]
But $\sql_{\linA}$ only becomes into an inner product
on the subspace $\domain{\sql_A}$,
if $\theta > 0$ or a further term is added
to make $\sql_{\linA}$ into a positive definite form.
In fact,
let $\linA$ be half-bounded with lower bound $\theta$,
then for each $\lambda > -\theta$
the subspace $\domain{\sql_{\linA}}$ becomes into a pre-Hilbert space
with inner product $[\cdot,\cdot]_{\sql,\lambda}$ given by
\[
[x,y]_{\sql_{\linA},\lambda} \eqdef
\scalar{\linA x}{y} \, + \, \lambda \, \scalar{x}{y}\,,
\quad x,y \xeof \domain{\sql_{\linA}}.
\]
and derived norm
\begin{equation}\label{ENERGETIC_NORM}
\norm{x}_{\sql_A,\lambda} \eqdef
\sqrt{[x,y]_{\sql_A,\lambda}} \=
\sqrt{\scalar{\linA x}{x} + \lambda \cdot \norm{x}^2}
\=
\sqrt{\scalar{\linA x + \lambda x}{x}}\,,
\end{equation}
called the \textit{energetic norm} induced by the operator $\linA$.
It should be noted
that if $\mu > -\theta$ is another real number,
then $\norm{x}_{\sql_A,\lambda}$ and $\norm{x}_{\sql_A,\mu}$
are equivalent norms,
that is,
they own the same Cauchy sequences in their common domain of definition.

\subparagraph*{}
However,
the domain $\domain{\sql_{\linA}}$ of the form $\sql_{\linA,\lambda}$
is not necessarily a complete normed space {\wrt} \eqref{ENERGETIC_NORM}.
Hence we define the space $\hilH_{\sql_{\linA},\lambda}$ to be the completion
of the domain $\domain{\sql_{\linA}}$ {\wrt} the topology 
induced by the norm $\norm{\cdot}_{\sql_{\linA},\lambda}$:
\[
\hilH_{\sql_A,\lambda} \eqdef
\overline{\domain{\sql_{\linA}}_{\norm{\cdot}_{\sql_{\linA},\lambda}}}.
\]

\subparagraph*{}
The sesquilinear form $[\cdot,\cdot]_{\sql_{\linA},\lambda}$,
defined on the subspace $\domain{\sql_A}$,
may be extended to the space $\hilH_{\sql_{\linA},\lambda}$ by
\begin{equation}\label{ENERGETIC_PRODUCT}
\overline{[x,y]_{\sql_{\linA},\lambda} }(x,y)
\eqdef
\lim_{n \xto \infty} \, [x_n,y_n]_{\sql_{\linA},\lambda}\,,
\end{equation}
where $(x_n)$ and $(y_n)$ are approximating sequences of
$(x,y) \xeof \hilH_{\sql_{\linA},\lambda} \xtimes \hilH_{\sql_{\linA},\lambda}$
{\wrt} the energetic norm
and with $x_n, y_n \xeof \domain{\sql_{\linA}}$ for all $n \xeof \N$.
Moreover,
it can be shown
that first the definition of the expression $\overline{[x,y]_{\sql{\linA},\lambda}}$
is independent from the choice of such sequences
and second $\hilH_{\sql_{\linA},\lambda} \eq \hilH_{\sql_{\linA},\mu}$ holds
for all $\lambda, \mu > -\theta$,
which gives rise to introduce
$\hilH_{\linA} := \hilH_{\sql_{\linA},\lambda}$ for some $\lambda > -\theta$.

\subparagraph*{}
The set $\hilH_{\linA}$ endowed with the inner product \eqref{ENERGETIC_PRODUCT}
forms a Hilbert space in its own right and
is called the \textit{energetic space} of the half-bounded operator $\linA$.
Its elements can be identified with elements in Hilbert space $\hilH$:
let $x_0 \xeof \hilH_{\sql_{\linA}}$,
then there is $(x_n)$ with elements $x_n \xeof \domain{\sql_{\linA}}$ such that
$(x_n)$ is Cauchy {\wrt} the energetic norm.
But then $(x_n)$ is also Cauchy {\wrt} the norm $\norm{\,\cdot\,}$
of the space $\hilH$,
thus there is $x_1 \xeof \hilH$ such that $\norm{x_1-x_n} \xto 0$.
The mapping $x_0 \xeof \hilH_{\linA} \mapsto x_1 \xeof \hilH$ proves to be injective
and hence forms a continuous embedding of $\hilH_A$ into $\hilH$,
so we can regard $(\hilH_A,\hilH)$ as a pair of Hilbert spaces
in the sense of section 12.

\paragraph*{}
The operator $\linA_F$
associated with the form $\sql_{\linA,\lambda}$
surely forms an extension of $\linA$ and proves to be self-adjoint
by theorem \ref{QUASI-COERCIVE_FORMS_AND_ASSOCIATED_OPERATORS}.
It is called the \textsc{Friedrichs} \textit{extension} of $\linA$.
Moreover,
the domain of definition of the operator $\linA_F$
is given by the linear subspace
$\domain{\linA_F} \eq \hilH_A \, \cap \, \domain{\linA^*}$.

\subparagraph*{}
The operator $\linA_F$ owns a remarkable property:
if one adds to the half-bounded operator $\linA$
the operator $\lambda \1$ ($\lambda \xeof \R$),
then the \textsc{Friedrichs} extensions
of $\linA$ and $\linA+\lambda \1$ are related to each other
by the equality
$(\linA + \lambda \1)_F \eq \linA_F + \lambda \1$.

\paragraph*{}
The \textsc{Friedrichs} extension method is an important tool
to construct self-adjoint exensions of half-bounded linear operators,
which can be applied to various linear partial differential operators
appearing in applications,
see part III for some examples of such constructions.

\section{Perturbation Theory}
%----------------------------
\textit{Perturbation theory} deals with operators of the form $\linA+\linB$,
where $\linA$ is assumed to be self-adjoint and $\linB$ to be symmetric
with domain of definition $\domain{\linB}$ being a superset of $\domain{\linA}$,
so $\linA + \linB$ is well defined with domain of definition $\domain{\linA}$.
It is the aim to decide
under which conditions on the operator $\linB$
the sum $\linA+\linB$ proves to be self-adjoint.
This question is especially of interest for applications in quantum mechanics.

\subsubsection*{} % A-boundedness
Let $\hilH$ be a complex Hilbert space and the operator
$\linA\:\domain{\linA} \xto \hilH$ be self-adjoint.
Then the following concept is useful:
a symmetric operator $\linB:\domain{\linS} \xto \hilH$
with $\domain{\linA} \xsubset \domain{\linB}$
is called $\linA$\textit{-bounded},
if there are constants $\theta, c \geq 0$ such that
\begin{equation}\label{A_BOUNDEDNESS}
\norm{\linB \,x} \, \leq \, \theta \norm{\linA \,x} \, + \, c \norm{x}
\qfa x \xeof \domain{\linA}.
\end{equation}
In order to ensure $\linA$-boundedness of $\linB$,
it is sufficient to ensure \eqref{A_BOUNDEDNESS}
only for elements $x$ in a core $\linA_0$ of $\linA$
(a \textit{core} of a closed operator $\linA$
is any restriction $\linA_0$ to a smaller domain,
whose closure coincides with $\linA$).

\subparagraph*{}
The number $\theta$ is called the $\linA$\textit{-bound} of $\linB$.
Notice that if $\linB$ is a bounded operator,
then it is $\linA$-bounded with $a<1$ for any $\linA$.
In addition,
the validness of the estimate
\[
\norm{\linB \,x}^2 \, \leq \,
\theta^2 \norm{\linA \,x}^2 \, + \, c^2 \norm{x}^2
\qfa x \xeof \domain{\linA}.
\]
is a sufficient condition for \eqref{A_BOUNDEDNESS} to be true.

\subsubsection*{} % Kato-Rellich theorem
Let $\linA$ be self-adjoint in Hilbert space $\hilH$ and $\linB$ 
a further symmetric operator having domain
$\domain{\linB} \supset \domain{\linA}$.
The following \textsc{Kato-Rellich}\textit{-theorem}
provides a criterion on the self-adjointness
of the operator $\linA+\linB$,
which has been independently found by \textsc{T.\,Kato}
and \textsc{F.\,Rellich}:

\begin{theorem}\label{KATO-RELLICH_THEOREM_I}
Let the linear operator $\linA$ be self-adjoint and 
the linear operator $\linB$ with $\domain{\linB} \supset \domain{\linA}$ 
be symmetric and $\linA$-bounded with $\theta < 1$.
Then $\linA+\linB$ with domain of definition $\domain{\linA}$ is self-adjoint.
\end{theorem}

\Proof:
Clearly, $\linA + \linS$ is symmetric.
For $\mu > 0$ we have
\[
\norm{(\linA - i \mu)^{-1} x} \, \leq \, \frac{1}{\mu} \, x,
\quad (x \xeof \hilH).
\]
Moreover, an elementary computation yields
$\norm{\linA x - i \mu x}^2 \eq \norm{\linA x}^2 + \abs{\mu} ^2 \norm{x}^2$.
Let $y := (\linA-i \mu) x$, then $x \eq (\linA -i \mu)^{-1} y$,
and we obtain
\[
\norm{y}^2 \= \norm{\linA (\linA -i \mu)^{-1} y}^2
           \, + \,
                   \abs{\mu} ^2 \norm{x}^2
           \, \geq \,
                   \norm{\linA (\linA -i \mu)^{-1} y}^2
\]
or equivalently $\norm{y} \geq \norm{\linA (\linA -i \mu)^{-1} y}$.
Now consider for any $y \xeof \hilH$ the operator equation
\begin{equation}\label{EQ_KATO-RELLICH}
z \, + \, \linS(\linA - i \mu)^{-1}z \= y.
\end{equation}
Since $\linS$ is $\linA$-bounded,
\[
\norm{\linS(\linA - i \mu)^{-1}z} \, \leq \,
\theta \norm{\linA(\linA - i \mu)^{-1}z} \, + c \, \norm{(\linA - i \mu)^{-1}z}.
\]
Inserting the estimations above into this equation, we obtain
\[
\norm{\linS(\linA - i \mu)^{-1}z} \,
\leq \,
\theta \norm{z} \, + \, \frac{c}{\abs{\mu} } \norm{z},
\]
which implies that the operator $\linS(\linA - i \mu)^{-1}$ is bounded.
Especially for $\mu$ small enough it follows that $\linS(\linA - i \mu)^{-1}$
is a contraction,
so by \textsc{Banach}\textit{'s fixpoint theorem}
the equation \eqref{EQ_KATO-RELLICH} has a unique solution $z$.
Finally,
letting $v := (\linA - i \mu)^{-1}z$ leads to
\[
(\linA + \linS -i \mu)v \=
\linA v - i\mu x \, + \, \linS x \= z \, + \, \linS(\linA-i \mu)^{-1} z
\= y,
\]
which implies that $\range{\linA + \linS -i\mu} \eq \hilH$.
But this is the range condition for the operator $\linA + \linS$.
\qed

\paragraph*{}
A similar result is obtained for the question,
whether \textit{essentially self-adjointness} of a linear operator $\linA$
is preserved by a symmetric and $\linA$-bounded perturbation $\linB$:

\begin{theorem}\label{KATO-RELLICH_THEOREM_II}
Let $\linA\:\domain{\linA} \xto \hilH$ be essentially self-adjoint
and assume that $\linB$ is symmetric and $\linA$-bounded,
that is,
$\domain{\linA} \xsubset \domain{\linB}$ and \eqref{A_BOUNDEDNESS} holds
with bound $\theta < 1$.
Then the operator $\linA + \linB$
is essentially self-adjoint on $\domain{\linA}$, where
$
\overline{\linA + \linB} \eq \overline{\linA} + \overline{\linB}.
$
\end{theorem}

\section{The Spectrum of Symmetric and Self-Adjoint Operators}
%-------------------------------------------------------------
A great portion of spectral theory deals with the study of the
resolvent set and the spectrum of certain classes
of closed or continuous linear operators.
Important examples are constituted by the classes of
symmetric and self-adjoint operators.

\subsubsection*{} % spectrum of symmetric operators
In order to investigate the spectrum of a symmetric operator,
we first start with the study of linear operators
of the form $\linA - \lambda$,
where $\linA$ is closed, densely defined and symmetric.

\subparagraph*{} %- \linA-\lambda is semi-Fredholm 
Remember that a number $\lambda \xeof \C$ with $\im \lambda \neq 0$
does not belong to the numerical range $\Theta(\linA)$ of $\linA$,
hence $\lambda$ is contained in the field $\Delta(\linA)$
of semi-Fredholm regularity.
%First notice that
%$(\linA-\lambda)^* = \linA^* - \overline{\lambda}$ for $\lambda \xeof \C$.
%Then a straight\-forward computation yields
%\[
%\norm{\linA \,x - \lambda x}^2
%\=
%\scalar{\linA \,x - \lambda x}{\linA \,x - \lambda x}
%\=
%\norm{\linA \,x}^2 -
%\scalar{\linA \,x}{\lambda x} -
%\scalar{\lambda x}{\linA \,x} +
%\abs{\lambda} ^2 \norm{x}^2
%\]
%\begin{equation}\label{RESOLVENT_IS_BOUNDED_FROM_BELOW}
%\=
%\norm{\linA \,x}^2 + \abs{\re \lambda} ^2 \norm{x}^2 -
%2 \,\re \lambda \cdot \scalar{\linA \,x}{x}
%+\abs{\im \lambda} ^2 \norm{x}^2.
%\end{equation}
%
%\subparagraph*{}
%Due to the Cauchy-Schwarz inequality and
%the estimate $2ab \leq \abs{a} ^2 + \abs{b} ^2$
%for $a,b \xeof \R$, we also have
%\[
%2 \, \abs{\re \lambda} \cdot \abs{\scalar{\linA \,x}{x}} \, \leq  \,
%2 \, \abs{\re \lambda} \cdot \norm{\linA \,x} \norm{x} \, \leq \,
%\norm{\linA \,x}^2 + \abs{ \re \lambda} ^2 \cdot \norm{x}^2,
%\]
%which together with inequality \eqref{RESOLVENT_IS_BOUNDED_FROM_BELOW}
%implies the estimate
%\begin{equation}\label{SYM1}
%\norm{\linA \,x - \lambda x} \, \geq \, \abs{\im \lambda} \, \norm{x}
%\qfa x \xeof \domain{\linA}.
%\end{equation}
This means that
$\linA - \lambda$ is bounded below,
and since we assume $\linA$ to be closed,
$\linA - \lambda$ is injective and normally solvable
by theorem \ref{CLOSED_RANGE_THEOREM},
thus upper semi-Fredholm with $\alpha(\linA-\lambda) \eq 0$.
By the punctured neighbourhood theorem \ref{PUNCTURED_NEIGHBOURHOOD_THEOREM},
the Fredholm index $\chi(\linA - \lambda)$
is locally constant for each fixed $\lambda \notin \R$.
In addition,
it can be shown that it is even globally constant
in the upper and lower complex half-plane,
respectively.

\subparagraph*{}
Before proceeding with the study of the range of $\linA-\lambda$,
notice that by using the orthogonal relations \eqref{ORTHOGONALITY_RELATIONS}
this space and the kernel of $\linA^*-\overline{\lambda}$
are coupled by
\[
\hilH \=
\range{\linA - \lambda}
\, \oplus \,
\kernel{\linA^* - \overline{\lambda}}.
\]

\subparagraph*{}
We now distinguish between two cases
{\wrt} surjectivity of $\linA - \lambda$,
or equivalently,
to the injectivity of $\linA^* - \overline{\lambda}$:

\begin{itemize}
\item[a)]
either $\range{\linA - \lambda} \eq \hilH$,
which implies $\kernel{\linA^* - \overline{\lambda}} \eq \{0\}$,
then $\linA - \lambda$ is invertible,
$\lambda$ is a regular value of $\linA$,
and $\mathbb{H}^+$ and $\mathbb{H}^-$ belongs to $\rho(\linA)$,
respectively,

\item[b)]
or $\range{\linA - \lambda} \neq \hilH$,
then $\kernel{\linA^* - \overline{\lambda}} \neq \{0\}$
and $\lambda$ belongs to the residual spectrum $\sigma_r(\linA)$,
so $\mathbb{H}^+$ and $\mathbb{H}^-$ are subspaces of $\sigma_r(\linA)$,
respectively.
\end{itemize}

\paragraph*{}
As it has already been indicated,
both cases may each happen for the upper and the lower complex half plane 
$\mathbb{H}^+$ and $\mathbb{H}^-$.
So in combination there are four possibilities
how the spectrum $\sigma(\linA)$ of a closed symmetric operator $\linA$
looks like:

\begin{itemize}[leftmargin=12pt,itemsep=5pt]
\item[1.]
$\sigma(\linA) \eq \C$
%the entire complex plane $\C$,
in case of $\mathbb{H}^+ \cup \mathbb{H}^- \xsubset \sigma_r(\linA)$,
\item[2.]
$\sigma(\linA) \eq \overline{\mathbb{H}^+}$
%the entire complex plane $\C$,
%the closure of the upper half-plane $\mathbb{H}^+$,
in case of $\mathbb{H}^+ \xsubset \sigma_r(\linA)$ and
$\mathbb{H}^- \xsubset \rho(\linA)$,
\item[3.]
$\sigma(\linA) \eq \overline{\mathbb{H}^-}$
%the entire complex plane $\C$,
%the closure of the lower half-plane $\mathbb{H}^-$,
in case of $\mathbb{H}^- \xsubset \sigma_r(\linA)$ and
$\mathbb{H}^+ \xsubset \rho(\linA)$,
\item[4.]
%a subset of the real line,
$\sigma(\linA) \xsubset \R$
in case of $\mathbb{H}^+ \cup \mathbb{H}^- \xsubset \rho(\linA)$.
\end{itemize}

\paragraph*{}
In summary,
we have
$\sigma(\linA) \in
\{\,\C,\overline{\mathbb{H}^+},\overline{\mathbb{H}^-},S\,\}$,
where $S$ is any subset of the real line.
But if the resolvent set $\rho(\linA)$
contains at least one element $\xi \xeof \R$,
then only the last case in the previous enumeration can occur.
In fact,
the property $\sigma(\linA) \xsubset \R$ leads us again
to the notion of \textit{self-adjointness},
since both properties prove to be equivalent to each other.

\subsubsection*{} % spectrum of self-adjoint operators
Our next subject is the study of the spectrum of a self-adjoint operator.
The first result is the claim
that self-adjoint operators differ from other operators
by the specific location of their spectra:

\begin{theorem}\label{CRITERION_FOR_SELF-ADJOINTNESS}
A closed linear operator $\linA$ is self-adjoint \iff
$\sigma(\linA) \subseteq \R$,
and in this case we have for $\im \xi \neq 0$ the resolvent estimation
\[
\norm{R(\linA,\xi)} \leq \frac{1}{\abs{\im \xi} }.
\]
\end{theorem}

\subparagraph*{}
Indeed,
$\sigma(\linA)$ being a subset of the real line means that
$\mathbb{H}^+$ and $\mathbb{H}^-$ belong to the resolvent set of $\linA$,
which is true if and only if the range condition is fulfilled,
that is, 
$\linA \pm i$ are both surjective.
The resolvent estimation is a consequence of the inequality
$\norm{R(\linA,\xi)} \leq \delta(\xi,\sigma(\linA))^{-1}$,
where $\delta(\xi,\sigma(\linA)$ denotes the positive distance between 
$\xi \xeof \rho(\linA)$ and the spectrum of $\linA$.

\subsubsection*{} % residual spectrum of self-adjoint operators
Remember the subsets of the spectrum $\sigma(\linA)$
of a linear operator $\linA$ defined in section 7.
They have been introduced by the criterion
how the perturbed operator $\linA - \lambda$ fails to be bijective.
This approach has led us to the notion of
\textit{point spectrum} $\sigma_p(\linA)$,
\textit{continuous spectrum} $\sigma_c(\linA)$ and
\textit{residual spectrum} $\sigma_r(\linA)$,
respectively.

\subparagraph*{}
The next result shows
that the division of the spectrum for self-adjoint operators
is simpler than in the arbitrary case:

\begin{theorem}\label{RESIDUAL_SPECTRUM_OF_SELF-ADJOINT_OPERATOR_IS_EMPTY}
If $\linA$ is self-adjoint,
then the residual spectrum $\sigma_r(\linA)$ is empty,
thus $\sigma(\linA) \eq \sigma_p(\linA) \cup \sigma_c(\linA)$.
\end{theorem}

\subparagraph*{}
This can be concluded by the orthogonality relation 
$\kernel{\linA} \eq \range{A}^{\bot}$,
so no spectral value may belong to the residual spectrum $\sigma_r(\linA)$.

\subsubsection*{} % pure point spectrum and pure continuous spectrum
When dealing with a self-adjoint linear operator $\linA$,
it is advantageous to breakup the set $\sigma(\linA)$
into other subsets as for arbitrary linear operators.

\subparagraph*{}
First,
the \textit{pure point spectrum} $\sigma_{pp}(\linA)$ is
\[
\sigma_{pp}(\linA) \eqdef
\{ \lambda \xeof \C:
\range{\linA - \lambda} \= \overline{\range{\linA - \lambda} }
\, \neq \, \hilH \}.
\]

\subparagraph*{}
Second,
by \textit{points embedded in continuum}
are meant the elements of the set
\[
\sigma_{pec}(\linA) \eqdef
\{ \lambda \xeof \C:
\range{\linA - \lambda} \, \neq \, \overline{\range{\linA - \lambda} }
\, \neq \, \hilH
\}.
\]
Both definitions do not require $\linA-\lambda$ to be injective,
so $\sigma_{pp}(\linA)$ and $\sigma_{pec}(\linA)$ 
possibly contain eigenvalues of $\linA$.
We join both subsets together and redefine the point specturm
$\sigma_p(\linA) := \sigma_{pp}(\linA) \sqcup \sigma_{pec}(\linA)$,
that is,
$\sigma_p(\linA)$ for $\linA$ being self-adjoint
consists of all $\lambda \xeof \C$,
for which $\overline{\range{\linA - \lambda}} \neq \hilH$.
This leads to the following

\begin{theorem}\label{POINT_SPECTRUM_IS_SET_OF_EIGENVALUES}
The point spectrum $\sigma_p(\linA)$ is exactly
the set of eigenvalues of $\linA$.
\end{theorem}

\subparagraph*{}
Third,
the \textit{purely continuous spectrum} is the set
\[
\sigma_{pc}(\linA) \eqdef
\{ \lambda \xeof \C:
\range{\linA - \lambda} \, \neq \, \overline{\range{\linA - \lambda} }
\= \hilH \}.
\]
We alternatively introduce the \textit{continuous spectrum}
$\sigma_c(\linA) := \sigma_{pec}(\linA) \sqcup \sigma_{pc}(\linA)$,
which of course is the set 
\[
\{ \lambda \xeof \C:
\range{\linA - \lambda} \, \neq \, \overline{\range{\linA - \lambda} } \}.
\]
In opposite to the former given definition of $\sigma_c(\linA)$
(see the paragraph on spectral theory),
this set may again contain elements $\lambda$,
for which $\linA-\lambda$ is not injective.

\paragraph*{}
For a self-adjoint operator $\linA$
it is also customary to denote $\sigma_{pp}(\linA)$ 
the \textit{discrete spectrum} 
and its complement relatively to
$\sigma(\linA)$ the \textit{essential spectrum}, {\ie}
\[
\sigma_{ess}(\linA) \eqdef
\sigma(\linA) \setminus \sigma_{pp}(\linA),
\]
so we have $\sigma_c(\linA) \xsubset \sigma_{ess}(\linA)$.

\subparagraph*{}
The \textit{discrete spectrum} of $\linA$ consists
of those values $\lambda \xeof \C$,
for which there is a terminating sequence $\{e_n\}$ in $\hilH$
consisting of pairwise disjoint unit elements $e_n$.
The essential spectrum of $\linA$ consists of those values $\lambda \xeof \C$,
for which there is a non-terminating sequence $\{e_n\}$ in $\hilH$
consisting of pairwise disjoint unit elements $e_n$,
called a \textit{Weyl sequence} for $\lambda \xeof \sigma_{ess}(\linA)$,
such that
\begin{equation}\label{WEYL_SEQUENCE}
\lim_{n \xto \infty} \, \norm{\linA e_n - \lambda e_n} \eq 0.
\end{equation}

\subparagraph*{}
In summary,
any number $\lambda \xeof \C$ belongs
to the spectrum of a self-adjoint operator $\linA$,
if it owns a (possibly terminating) Weyl sequence
fulfilling \eqref{WEYL_SEQUENCE}.

\section{Projection-Valued Measures and Spectral Integrals}
%----------------------------------------------------------
The {spectral theorem} is definitely the most important assertion
{\wrt} self-adjoint operators.
Another way to claim this theorem reads as follows:
each self-adjoint linear operator $\linA$ is a \textit{spectral operator},
whose underlying \textit{projection-valued measure} $\linP$
can be retrieved from $\linA$.
In the following we want to make this assertion precise
by going on step by step.

\subsubsection*{} % projections
As the projection theorem shows,
projection mappings acting in a Hilbert space $\hilH$
are in one-to-one relationship to closed linear subspaces in $\hilH$.
So it is nearby
to call two projections $\pi_1,\pi_2$ \textit{orthogonal} to each other,
if their related closed subspaces have this property,
and in this case the identity
$\pi_1 \circ \pi_2 \eq \pi_2 \circ \pi_1 \eq \0$ holds,
where $\0$ is the projection related
with the trivial subspace $\{\0\} \xsubset \hilH$.

\paragraph*{}
In general,
a projection mapping $\pi$ has the following properties:

\begin{itemize}[leftmargin=12pt]
\item[1.]
$\pi$ is idempotent: $\pi \circ \pi \eq \pi$.

\item[2.]
$\pi$ is linear on $\hilH$.

\item[3.]
$\pi$ is bounded with $\norm{\pi} \leq 1$,
hence $\pi \xeof \BO{\hilH}$.\\
If $\setA_{\pi} := \pi(\hilH) \eq \{\0\}$, then $\norm{\pi} \eq 0$,
else $\norm{\pi}=1$.
The only possible eigenvalues of $\pi$ are the real numbers $0$ and $1$.

\item[4.]
Let $\setA_1,\setA_2$ be closed subspaces of $\hilH$
with $\setA_1 \subseteq \setA_2$,
then for the associated projections the equation
$\pi_{\setA_1} \circ \pi_{\setA_2} \eq
 \pi_{\setA_2} \circ \pi_{\setA_1} \eq \pi_{\setA_1}$
holds.
In this case,
the operator $\pi_{\setA_2} - \pi_{\setA_1}$ is also a projection,
and related subspace is given by the set
$\setA_2 \ominus \setA_1 := \setA_2 \cap \setA_1^{\bot}$.

\item[5.]
$\pi \eq \pi^{*}$, that is, $\pi$ is symmetric.
Conversely,
each symmetric and idempotent linear operator
is a projection.

\item[6.]
If $\setA$ is a \textit{finite}-dimensional subspace of $\hilH$,
then $\setA$ is closed,
and the projection $\pi_{\setA}$ can be represented
by means of a set of base vectors $e_1,\ldots,e_n$ in $\setA$, namely
\[
\pi_{\setA} x \= \sum_{i=1}^n \scalar{x}{e_i} \, e_i.
\]
\end{itemize}

\subparagraph*{}
Finally,
let $\mathscr{P}(\hilH)$ be the set of projections in $\hilH$,
then there is a partial order relation $\geq$ on $\mathscr{P}(\hilH)$,
defined by means of inclusion between the related closed subspaces:
\[
\pi_{\setA_1} \geq \pi_{\setA_2} \quad
: \Longleftrightarrow \quad \setA_2 \subseteq \setA_1.
\]

\subparagraph*{}
By a \textit{projection-valued measure} we mean a mapping
$\linP:\Borel{\R} \xto \BO{\hilH}$,
where $\Borel{\R}$ is the ordinary Borel $\sigma$-algebra in $\R$
and $\linP$ fulfills the following properties:

\begin{itemize}[leftmargin=22mm]
\item[(P-1)]
$\linP(\Omega)$ is a projection operator for each $\Omega \xeof \Borel{\R}$;

\item[(P-2)]
$\linP(\R) \eq \1_{\hilH}$, the identity operator in $\hilH$;

\item[(P-3)]
$\linP$ is \textit{$\sigma$-additiv},
that is,
if $\Omega_1, \Omega_2,\ldots$ are measurable sets in the real line
and $\Omega := \bigsqcup_{\,n \xeof \N} \Omega_n$,
then
\[
\linP(\Omega) \= \lim_{n \xto \infty} \, \sum_{k=1}^n \, \linP(\Omega_k),
\]
where this kind of convergence has to be understood
in the sense of the strong operator topology in $\BO{\hilH}$.
\end{itemize}

\subparagraph*{}
The projection-valued measure $\linP$ is also called
a \textit{spectral measure} or a \textit{resolution of the identity}.
Notice that
intervals of the form $(-\infty,\lambda]$ are Borel-measurable,
so each projection $\linP((-\infty,\lambda])$ is well-defined.
In fact,
the one-parameter family of these projections constitutes
a projection-valued measure $\linP$.

\paragraph*{}
In the sequel
it is our aim to introduce a \textit{functional calculus}
for measurable functions $f\:\R \xto \C$
{\wrt} a projection-valued measure $\linP$.

\subparagraph*{}
Let $B(\R)$ be the linear space of \textit{bounded} measurable functions
$f\:\R \xto \C$ endowed with the norm
\[
\norm{f}_{\infty} \eqdef \sup_{x \xeof \R} \, \abs{f(x) } < \infty,
\]
which in fact is a Banach space.
Then a function $e \xeof B(\R)$ is called \textit{simple},
if it adopts only finitely many values $e_1,\ldots, e_N$.
Moreover, the set
\[
S(\R) := \{\,f\:\R \xto \C: f \mbox{ is simple}\,\}
\]
forms dense subspace of $B(\R)$,
and the supremum norm of $e \xeof S(\R)$ is given by
$\norm{e}_{\infty} \eq \max \, \{\abs{e_1} ,\ldots, \abs{e_N} \}$.

\paragraph*{}
After these preliminaries,
the setup of the spectral integral will be carried out
by the following five steps:

\begin{itemize}[leftmargin=12pt,itemsep=5pt]
\item[1.]
First of all,
each spectral measure $\linP$ gives rise
to a complex and real-valued Borel measure on the real line.
For arbitrary $x,y \xeof \hilH$
let the complex measure $\mu_{x,y}$ be defined by
\[
\begin{cases}[1.2]
\quad \mu_{x,y}: \Borel{\R} \xto \C,\\
\quad \Omega \mapsto \mu_{x,y}(\Omega) := \scalar{x}{\linP(\Omega)y}.\\
\end{cases}
\]
Due to the symmetry of $\linP(\Omega)$,
the derived measure
$\mu_x(\Omega) := \mu_{x,x}(\Omega)$ for $\Omega \xeof \Borel{\R}$ is real-valued.

\item[2.]
We next make sense of the expression $\int_{\R} f \, dP$
for measurable functions $f\:\R \xto \C$.
Let us start with the case, where $f$ is \textit{simple},
that is,
$f$ takes values in the finite set $\{f_1,\ldots,f_N\}$.
Define
\[
\int_{\R} f \, d\linP \eqdef \sum_{n=1}^N f_n \linP(\chi^{-1}(f_n)),
\]
where $\chi^{-1}(f_n)$ is the preimage of $f_n$ in $\Borel{\R}$.
We further mention
that
\[
I_{0,\linP}:f \xto \int_{\R} f \, d\linP
\]
is a continuous linear mapping from the normed space
$S(\R)$ consisting of simple functions
to the Banach space $\BO{\hilH}$ with operator norm $\norm{I_{0,\linP}}=1$.

\item[3.]
We shall continue with the integration
of \textit{bounded measurable functions} {\wrt} $\linP$.
Due to the B.L.T. theorem,
there exists a unique extension $I_P$ from $I_{0,\linP}$
to the closure $B(\R) = \overline{S(\R)}$,
where $\norm{I_P} \eq \norm{I_{0,P}}$.
The mapping $I_P$ is also linear and fulfills
\[
\begin{cases}[1.2]
\quad I_P(\1_{\R}) = \1_{\hilH}\\
\quad I_P(fg) = I_P(f) \circ I_P(g)\\
\quad I_P(\overline{f}) = I_P(f)^*\\
\end{cases}
\quad \quad f,g \xeof B(\R),
\]
where $\1_{\R}$ is the identity in $\R$,
thus $I_P$ even forms a C*-algebra homomorphism between
the function algebra $B(\R)$ and the operator algebra $\BO{\hilH}$.

\item[4.]
We will also define $I_P$
for an \textit{arbitrary measurable function} $f\:\R \xto \C$.
But this is only possible for those elements $x \xeof \hilH$ fulfilling
\begin{equation}\label{PVM_OPERATOR}
\int_{\R} \abs{f} ^2 \, d \mu_x < \infty.
\end{equation}
The set $D_f$ of elements $x \xeof \hilH$ satisfying \eqref{PVM_OPERATOR}
forms a linear subspace lying dense in $\hilH$.
For $f$ we introduce an approximating sequence of functions
$f_n \xeof B(\R)$ given by
\[
f_n \eqdef \chi_{ \{t \xeof \R: \abs{f(t)} \leq n\}} \cdot f
\quad (n \xeof \N).
\]
One can show that $f_n$ is Cauchy in $L^2(\R)$,
hence $I_P(f_n)$ is Cauchy in $\hilH$,
where $I_P$ is the previously defined operator acting on the space $B(\R)$.

\item[5.]
To the end,
for a fixed measurable function $f$
we may define an operator $S(\linP,f)$,
called the \textit{spectral operator} {\wrt} $\linP$,
by virtue of
\begin{equation}\label{FUNCTIONAL_CALCULUS}
\int_{\R} f \, d\linP: D_f \mapsto \hilH\,,
\quad \quad
x \xeof D_f \mapsto (\int_{\R} f \, d\linP\,) \, x \eqdef y\,,
\end{equation}
which enjoys the property
\[
(\int_{\R} f \, d\linP)^* \= \int_{\R} \overline{f} \, d\linP.
\]
Especially for $f \eq \1_{\R}$ 
the self-adjoint operator $\linS_P := \linS(\linP,\1)$ is defined by
\[
\domain{\linS_P} := D_{\1_{\R}},
\quad \linS_P x \eqdef (\int_{\R} \1_{\R} \, d\linP) \, x.
\]
\end{itemize}

\section{The Spectral Theorem}
%-----------------------------
In linear algebra
the \textit{main axis theorem} quotes
that each hermitean matrix is diagonalizable
and has purely real-valued eigenvalues.
The generalization of this assertion,
substituting hermitean matrices by densely defined self-adjoint operators
acting in abstract Hilbert space,
leads to the so-called \textit{spectral theorem}.

\subsubsection*{} % multiplication operators
Remember from integration theory
that for a given measure space $(X,\salgB,\mu)$
the linear space $L^2(X)$ of measurable and square-integrable functions
$f\:X \xto \C$ forms a Hilbert space with inner product
\[
\scalar{f}{g} \eqdef \int_X f \, \overline{g} \, d\mu.
\]
If $m:X \xto \C$ is a measurable function,
then the linear mapping $\linM_m$ defined by
\[
\begin{cases}[1.2]
\quad \domain{\linM_m} & := \,
\{\, f \xeof L^2(X): m \cdot f \xeof L^2(X)\,\}, \\
\quad (\linM_m f)(x) & := \, m(x) \cdot f(x) \;
\mbox{ for } x \xeof X \mbox{ and } f \xeof \domain{\linM_m} \\
\end{cases}
\]
is called a \textit{multiplication operator} with \textit{generator} $m$.
It turns out that the mapping $\linM_m$ is symmetric
iff
$m$ is real-valued,
and since the domains of definition
of the operators $\linM_m$ and $\linM_m^*$ coincide,
each multiplication operator with real-valued generator is self-adjoint.

\paragraph*{}
The \textit{spectral theorem for self-adjoint operators} claims
that the reverse assertion is also true:

\begin{theorem}\label{SPECTRAL_THEOREM_VERSION_I}
To each self-adjoint operator
$\linA\:\domain{\linA} \xsubset \hilH \xto \hilH$
there is assigned

\begin{itemize}[leftmargin=7mm]
\item[(1)]
a measure space $(X,\salgB,\mu)$,

\item[(2)]
a multiplication operator
$\linM_m$ in $L^2(X)$ with generator function $m:X \xto \R$,

\item[(3)]
a bijective and isometric linear operator $\linU:\hilH \xto L^2(X)$
\end{itemize}

such that $\linA \eq \linU^{-1} \circ \linM_m \circ \linU$ holds.
\end{theorem}

\paragraph*{}
The existence of the operator $\linU$ gives rise to speak of
\textit{unitary equivalence} between the operators $\linA$ and $\linM_m$.
If some operators $\linA_1$ and $\linA_2$ are unitary equivalent,
then their spectra $\sigma(\linA_1)$ and $\sigma(\linA_2)$ are identical.

\subparagraph*{}
In summary, 
theorem \ref{SPECTRAL_THEOREM_VERSION_I} implies
that an operator $\linA$ is self-adjoint in Hilbert space $\hilH$
iff
it is unitary equivalent to a real-valued multiplication operator $\linM_m$
on $L^2(X)$,
where $X$ is the base set of a suitable chosen measure space $(X,\salgB,\mu)$.
%For linear differential operators $\DO_c$ with constant real coefficients
%the unitary equivalence between
%the self-adjoint realization in Hilbert space $L^2(\Rd)$ and
%the polynomial multiplication operator $M_{\DO_c}$ is given
%by the Fourier-Plancherel operator $\FT$,
%that is,
%$\DO_c \eq \FT^{-1} \circ M_{\DO_c} \circ \FT$.

\subsubsection*{} % projection measure variant of the spectral theorem
Using the notion of projection-valued measure,
we are now able to quote a further variant of
the \textit{spectral theorem for self-adjoint operators}:

\begin{theorem}\label{SPECTRAL_THEOREM_VERSION_II}
Given a self-adjoint linear operator $\linA$,
there is a uniquely defined projection-valued measure $P$
such that $\linA$ is the spectral operator related to $\linP$:
\[
\linA \= \int_{\,\R} \1 \, d\linP.
\]
\end{theorem}

\subsubsection*{} % functional calculus for self-adjoint operators
This version of the spectral theorem is extremely useful
to define operator-valued mappings $f\:\sigma(\linA) \xto f(\linA)$ by
\[
f(\linA) \eqdef \int_{\,\R} f \, d\linP,
\]
where $f(\linA)$ will become into a densely defined operator
acting in $\hilH$.
Indeed,
passing from a self-adjoint linear operator
to its family of projection-valued measures
and then to the spectral operators defined by \eqref{FUNCTIONAL_CALCULUS}
is nothing but a \textit{functional calculus}
for the class of self-adjoint operators.
For example, 
if $\linA$ is self-adjoint and \textit{positive} at once,
that is,
$\scalar{\linA x}{x} \geq 0$ for all $x \xeof \domain{\linA}$,
then we may define operators like $\sqrt{\linA}$, 
$\linA^n$, $\exp(\linA)$, $\sin(\linA)$, $\cos(\linA)$ and so on.

%\input{PROOF_SPECTRAL_THEOREM_II}

\subsubsection*{} % spectral resolutions
It is also possible to formulate the spectral theorem by means of so-called
\textit{spectral resolutions} instead of projection-valued measures.
Preliminary,
to each projection-valued measure $\linP$ in Hilbert space $\hilH$
there is related the concept of \textit{spectral resolution},
which is nothing but a one-parameter family
$\{\linP_{\lambda}\}_{\lambda \xeof \R}$
of linear operators $\linP_{\lambda}:\hilH \xto \hilH$
such that

\begin{itemize}[leftmargin=11mm]
\item[(R-1)]
$\linP_{\lambda}$ is a projection for each $\lambda \xeof \R$.
\item[(R-2)]
For each element $x \xeof \hilH$ the mapping
$\lambda \mapsto \linP_{\lambda} x$ is continuous from the right.
\item[(R-3)]
$\lambda < \mu$ implies $\linP_{\lambda} \leq \linP_{\mu}$,
that is,
$\scalar{\linP_{\lambda} x}{x} \leq \scalar{\linP_{\mu} x}{x}$
for all $x \xeof \hilH$.
\item[(R-4)]
$\lim_{\lambda \xto -\infty} \linP_{\lambda} \, x \eq 0$
and
$\lim_{\lambda \xto  \infty} \linP_{\lambda}  \, x \eq x$
for all $x \xeof \hilH$.
\end{itemize}

\paragraph*{}
Each projection $\linP_{\lambda}$ of a given spectral resolution
is related to a closed linear subspace $\setA_{\lambda} \xsubset \hilH$
and the difference $\linP_{\mu} - \linP_{\lambda}$
for $\mu \geq \lambda$ also forms a projection related to the closed subspace
$\setA_{\mu+0} \ominus \setA_{\lambda}$.

\paragraph*{}
Now,
using spectral resolutions instead of projection-valued measures,
the spectral theorem claims
that each self-adjoint operator $\linA$
owns a spectral resolution $\{\linP_{\lambda}\}_{\lambda \xeof \R}$
such that
\[
\domain{\linA} \=
\{\, x \xeof \hilH:
    \int_{\,\R}
\lambda^2 \, d \scalar{\linP_{\lambda} \,x}{x} < \infty\,\}\,,
\]
\[
\scalar{\linA \, x}{y} \=
\int_{\,\R} \lambda \, d\scalar{\linP_{\lambda} \,x}{y}
\quad \mbox{for } x \xeof \domain{\linA} \mbox{ and } y \xeof \hilH.
\]

\subsubsection*{} % relationship between spectrum and spectral resolution
We next want to investigate the relationship between the shape of the spectrum
and the spectral resolution of a self-adjoint operator.

\subparagraph*{}
In fact,
the following theorem shows
that the resolvent set and the different kinds of spectra
of a self-adjoint linear operator $\linA$
are determined by its spectral resolution:

\begin{theorem}\label{SPECTRUM_RESOLUTION_THEOREM}
Let $\linA$ be self-adjoint in $\hilH$
and $\setA_{\mu}$ be the closed subspace in $\hilH$
associated with the projection $\linP_{\mu}$ for some $\mu \xeof \R$.
Then

\begin{itemize}[leftmargin=12pt]
\item[1.]
$\mu \xeof \R$ belongs to the \textit{point spectrum} $\sigma_p(\linA)$
iff
$0 < \dim \; (\setA_{\mu+0} \ominus \setA_{\mu}) < \infty$.
%Besides for some number $\epsilon > 0$
%the subspace
%$\setA_{\mu+\epsilon} \ominus \setA_{\mu-\epsilon}$ is finite-dimensional.
%Hence $\mu$ is not only an eigenvalue of $\linA$,
%but also belongs to the set $\sigma_d(\linA)$.

\item[2.]
$\mu \xeof \R$ belongs to the \textit{continuous spectrum} $\sigma_c(\linA)$
iff
for each $\epsilon > 0$ the subspace
$\setA_{\mu+\epsilon} \ominus \setA_{\mu-\epsilon}$ is infinite-dimensional.

\item[3.]
$\mu \xeof \R$ is a \textit{regular value} of $\linA$
iff
there is an open interval $(\alpha,\beta) \xsubset \R$
such that $\mu \xeof (\alpha,\beta)$
and $\linP_{\mu}$ is constant in $(\alpha,\beta)$.
\end{itemize}
\end{theorem}

%\input{PROOF_SPECTRUM_RESOLUTION_THEOREM}

\paragraph*{}
Especially the resolvents $R(\linA,\xi)$ for $\xi \xeof \rho(\linA)$ 
of a self-adjoint operator $\linA$
can be represented by means of the spectral resolution:
\[
R(\linA,\xi) \, x \= \int_{-\infty}^{\infty}
\frac{1}{\lambda-\xi} \, dP_{\lambda} x,
\quad (x \xeof \hilH, \, \im \xi \neq 0).
\]

\subparagraph*{}
On the other hand,
\textsc{Stone}\textit{'s formula} provides a constructive method
to describe the spectral resolution of $\linA$ in an interval $[a,b]$
by means of the resolvents $R(\linA,\xi)$.
Let $x \xeof \hilH$ be arbitrary, then
\[
\frac{1}{2} \{P[a,b] + P(a,b)\} \, x
\=
\lim_{\epsilon \downarrow 0} \, \frac{1}{2 \pi i} \,
\int_a^b
\{R(\linA,t + i \epsilon)\,x -
  R(\linA,t - i \epsilon)\,x \} \,
  dt.
\]

\section{Self-Adjoint Operators with Pure Point Spectrum}
%--------------------------------------------------------
We now turn to linear operators $\compA$
acting in a \textit{separable} Hilbert space $\hilH$,
which are compact and symmetric at the same time.
Since compactness implies continuity,
it may be assumed that $\compA$ is everywhere defined in $\hilH$.

\subparagraph*{}
Of course,
if an operator is compact-symmetric,
the properties of compactness and symmetry complement,
so we may quote
that the spectrum $\sigma(\compA)$ of such an operator consists
of a sequence $\{\lambda_n\}$ of \textit{real} eigenvalues of $\compA$,
which either terminates or converges to $0$.
Due to the compactness of $\compA$,
the related eigenspaces $E_{\lambda_n}$ are finite-dimensional.
By the \textsc{Gram-Schmidt} orthogonalization method,
one can define a subset $\orthO$ of the totality
of eigenvectors of $\compA$
such that $\orthO$ spans a dense subspace in $\hilH$
and consequently constitutes an \textit{orthonormal base} in $\hilH$.

\subparagraph*{}
It is the assertion of the \textsc{Hilbert-Schmidt} \textit{spectral theorem},
also often called the \textit{discrete spectral theorem},
that there is a representation of a symmetric-compact operator $\compA$
in a separable Hilbert space $\hilH$
as a \textit{countable} sum of operators $\compA_{\lambda_n}$,
where each of them is non-trivial on a finite-dimensional subspace
of $\hilH$ and acts there like a simple multiplication operator
$\compA_{\lambda_n} := \lambda_n \1$.
In other words,
the theorem provides an orthogonal decomposition of $\hilH$
into a sequence of $\compA$-invariant subspaces:

\begin{theorem}\label{DISCRETE_SPECTRAL_THEOREM}
Let the Hilbert space $\hilH$ be separable.
Given a symmetric-compact linear operator $\compA:\hilH \xto \hilH$,
there is a null sequence $\{\lambda_n\}$ consisting of real eigenvalues
and an ortho\-normal base in $\hilH$
consisting of eigenvectors $(x_n)$ of $\,\compA$
such that for all $x \xeof \hilH$ the following representation holds:
\begin{equation}\label{HILBERT_SCHMIDT_REPRESENTATION}
\compA \,x \=
\sum_{n=1}^{\infty} \, \lambda_n \, \scalar{x}{x_n} \, x_n.
\end{equation}
\end{theorem}

\subsubsection*{} % operators with compact resolvents
We finally consider linear operators $\linA$
having the property $\sigma(\linA) \eq \sigma_p(\linA)$,
that is,
the continuous and residual spectrum of $\linA$ are both empty
and $\sigma(\linA)$ only consists of eigenvalues.
Such operators are of outstanding meaning in applied analysis,
since they appear as linear differential operators
defined on \textit{bounded} domains $\G \xsubset \Rd$.

\subparagraph*{}
A closed operator $\linA$ with domain of definition $\domain{\linA}$
is called \textit{having compact resolvents},
if at least one resolvent $\linA -\xi$ with $\xi \xeof \rho(\linA)$ is compact.
The resolvent equation \eqref{RESOLVENT_EQUATION} shows
that in this case {every} other resolvent of $\linA$ is compact.

\paragraph*{}
There is a useful criterion to check
whether $\linA$ has this distinguished property:

\begin{theorem}\label{CRITERION_FOR_OPERATORS_HAVING_PURE_POINT_SPECTRUM}
Let the domain of definition $\domain{\linA} \xsubset \banX$
of a densely defined operator $\linA$ in $\banX$
be equipped with the graph norm defined by \eqref{GRAPH_NORM}.
Then $\linA$ has compact resolvents
\iff the natural embedding
$\iota:
( \domain{\linA} ,\norm{\cdot}_A )
\hookrightarrow
(\banX,\norm{\cdot})$
is compact.
\end{theorem}

\subparagraph*{}
One also says that an operator having compact resolvents
is an \textit{operator with pure point spectrum},
since each $\lambda \xeof \sigma(\linA )$
is an eigenvalue of finite multiplicity.
This follows from the fact
that $x \xeof \banX$ is a solution of the equation $\linA  x \eq \lambda x$
\iff
$x$ solves the equation \textit{in resolvent form}
\[
(\lambda - \xi)^{-1} x \= R(\linA ,\xi) x\,,
\quad \xi \xeof \rho(\linA ),
\]
so both equations are equivalent to other.
Then the left hand side of this equation  must be of the form
$(\lambda - \xi)^{-1} x  \eq \zeta_n(\xi) x$ for some $n \xeof \N$,
where $\,\{\zeta_n(\xi)\}_{n \xeof \N}$
is the sequence of eigenvalues of $R(\linA ,\xi)$.
This implies that $\sigma(\linA)$ is a subset of the set
$\{\,\xi + \zeta_n(\xi)^{-1}: n \xeof \N\,\}$.

\subparagraph*{}
It is quite natural to order the set of eigenvalues of $\linA$ by their moduli.
Hence one obtains a sequence $\{\lambda_n\}$ of eigenvalues $\lambda_n$
such that
\[
\abs{\lambda_1} \, \leq \, \abs{\lambda_1} \, \leq \, \abs{\lambda_3} \, \dots
\, \xto \infty.
\]

\subparagraph*{}
With respect to the solution theory of the equation $\linA x \eq y$,
an operator $\linA$ is called \textit{completely solvable},
if $\linA$ has compact resolvents and $\xi \eq 0$ is a regular value of $\linA$,
so $\linA^{-1}$ is bounded
and has range $\range{\linA^{-1}} \eq \domain{\linA}$.

\subsubsection*{}
Let $\linA$ be a self-adjoint linear operator with pure point spectrum.
Then each resolvent of $\linA$ is symmetric and compact, 
and by the \textsc{Hilbert-Schmidt} spectral theorem
\ref{DISCRETE_SPECTRAL_THEOREM} 
the set of eigenvectors forms an orthonormal base in $\hilH$.
But these eigenvectors are precisely those,
that also belong to the eigenvalues of $\linA$.

\subparagraph*{}
We summarize the preceeding considerations to the following
\textit{spectral theorem for self-adjoint operators with pure point spectrum}:

\begin{theorem}\label{SPECTRAL_THEOREM_FOR_OPERATORS_WITH_COMPACT_RESOLVENTS}
Let $\linA$ with domain of definition $\domain{\linA} \xsubset \hilH$
be self-adjoint and with pure point spectrum
$\sigma(\linA) \eq \{\, \lambda_n: n \xeof \N \,\} \xsubset \R$.
Then $\abs{\lambda_n} \xto \infty$ as $n \xto \infty$, and
\[
\linA x \= \sum_{n=1}^{\infty} \lambda_n \scalar{x}{x_n} \, x_n
\; \mbox{ for } \;
x \xeof \domain{\linA} \eqdef
\{ \, x \xeof \hilH:
\sum_{n=1}^{\infty} \lambda_n^2 \, \abs{\scalar{x}{x_n}} ^2 < \infty \, \}.
\]
\end{theorem}

\subparagraph*{}
This theorem is of fundamental meaning in the investigation of eigenvalue
problems on bounded domains {\wrt} formally self-adjoint 
linear differential operators.

\section{Evolution Equations Governed by Self-Adjoint Operators}
%---------------------------------------------------------------
The functional calculus for self-adjoint operators permits us to construct
solutions for the homogeneous Cauchy problem of first and second order.
For this we first have to generalize the concept of classical derivative
for functions with values in a Banach space $\banX$.

\subsubsection*{} % evolution equations governed by self-adjoint operators
The analysis of linear operator-differential equations
In the Hilbert space setting
can directly be attacked for the case,
where the differential equations is governed
by a positive self-adjoint operator.
The following existence and uniqueness theorems all rely on the spectral
theorem and the functional calculus for self-adjoint unbounded operators.

\paragraph*{}
In the sequel we deal with first order problems and shall show
how one can solve the homogeneous Cauchy problem \textit{of first order}
\begin{equation}\label{HOMOGENEOUS_CAUCHY_PROBLEM_1}
\begin{cases}[1.2]
\quad u'(t) \= \linA \, u(t) \\
\quad u(0)=u_0 \xeof \domain{\linA}, \\
\end{cases}
\end{equation}
where $\linA$ is self-adjoint and positive.

\subparagraph*{}
Remember that a linear operator $\linA$ acting in a Hilbert space
is called \textit{positive},
if $\linA$ is symmetric and $\scalar{\linA x}{x} \geq 0$
for all $x \xeof \domain{\linA}$.
If $\linA$ is positive,
then any solution of the problem \eqref{HOMOGENEOUS_CAUCHY_PROBLEM_1}
satisfies the estimate
\begin{equation}\label{ESTIMATION_FOR_NONNEGATIVE_OPERATOR}
\norm{u(t)} \, \leq \, \norm{u(0)}
\qfa t > 0\,,
\end{equation}
which is a consequence of the inequality
$\norm{u(t)}^2 \eq - 2 \, \scalar{\linA u}{u} \leq 0$
and causes the solution of the problem to be unique.

\subparagraph*{}
To verify the existence of a solution for \eqref{HOMOGENEOUS_CAUCHY_PROBLEM_1},
let $(P_{\lambda})_{\lambda \xeof \R}$ be the spectral resolution
of the self-adjoint operator $\linA$.
Notice that $P_{\lambda} \eq \0$ for $\lambda < 0$,
since $\linA$ is positive.
Then the functional calculus for the continuous function $t \mapsto e^{-t}$
gives rise to the spectral operator
\begin{equation}\label{SPECTRAL_OPERATOR_I}
\linL_t(x)
\eqdef
e^{-\linA t} x
\=
\int_{\,0}^{\,\infty} e^{-\lambda t} \, dP_{\lambda} x
\quad \quad (t \geq 0, \, x \xeof \hilH).
\end{equation}
Obviously,
we have $\linL_0 \eq \1$,
and \eqref{ESTIMATION_FOR_NONNEGATIVE_OPERATOR} implies
$\norm{\linL_t} \leq 1$ for all $t > 0$.

\subparagraph*{}
In conclusion,
the existence of the solution of \eqref{HOMOGENEOUS_CAUCHY_PROBLEM_1}
is ensured by the following

\begin{theorem}\label{EXISTENCE_THEOREM_FOR_ABSTRACT_HEAT_EQUATION}
The function $u\:t \xeof [0,\infty) \mapsto \linL_t(u_0)$ is continuous.
Moreover,
it is continuously differentiable in $(0,\infty)$.
and forms a solution of \eqref{HOMOGENEOUS_CAUCHY_PROBLEM_1}.
\end{theorem}

%\input{PROOF_EXISTENCE_THEOREM_FOR_ABSTRACT_HEAT_EQUATION}

\paragraph*{}
We also apply this method to the homogeneous Cauchy problem
\textit{of second order}
\begin{equation}\label{HOMOGENEOUS_CAUCHY_PROBLEM_2}
\begin{cases}[1.2]
\quad u''(t) \= \linA \, u(t) \\
\quad \; u(0)=u_0 \xeof \domain{\linA} \\
\quad u'(0) = u_1 \xeof \hilH,\\
\end{cases}
\end{equation}
where again the operator $\linA$ is assumed to be self-adjoint and positive.
Notice that under this condition
the square root operator $\sqrt{\linA}$ is well-defined.

\subparagraph*{}
Each solution of this Cauchy problem satisfies the energy conservation law
\[
\norm{u'(t)}^2 \, + \, \norm{\sqrt{\linA} \, u(t)}^2
\=
\norm{u'(0)}^2 \, + \, \norm{\sqrt{\linA} \, u(0)}^2
\quad (t > 0),
\]
which implies
that the solution of the problem must be unique.

\paragraph*{}
Rather similar to the first order Cauchy problem,
the existence of a solution for \eqref{HOMOGENEOUS_CAUCHY_PROBLEM_2}
can be solved by 
the functional calculus method for self-adjoint operators.

\subparagraph*{}
Let $(P_{\lambda})$ be the spectral resolution
of a self-adjoint operator $\linA$.
Then we define the function $u\:[0,\infty) \xto \hilH$
by applying the functional calculus to the functions
$x \mapsto \cos(\sqrt{x})$ and
$x \mapsto \sin(\sqrt{x})/\sqrt{x}$.
In fact, the function $u$ is built up 
as the sum of two spectral operators, namely
\begin{equation}\label{SPECTRAL_OPERATOR_II}
u(t)
\eqdef
\linL_t(u_0,u_1)
\eqdef
\int_{\,0}^{\,\infty} \cos{\sqrt{\lambda}t} \, dP_{\lambda} \, u_0
\, + \,
\int_{\,0}^{\,\infty} \frac{\sin{\sqrt{\lambda}t}}{\sqrt{\lambda}}
\, dP_{\lambda} \, u_1
\quad \quad (t > 0).
\end{equation}

\paragraph*{}
When defining the family of linear operators $\{\linL_t(u_0,u_1)\}$ this way,
we obtain the following existence and uniqueness theorem
for second order problems:

\begin{theorem}\label{EXISTENCE_THEOREM_FOR_ABSTRACT_WAVE_EQUATION}
The mapping $t \mapsto \linL_t(u_0,u_1)$ is continuous in $[0,\infty)$
and twice continuously differentiable in $(0,\infty)$.
Furthermore,
it solves problem \eqref{HOMOGENEOUS_CAUCHY_PROBLEM_2} in a unique manner.
\end{theorem}

%\input{PROOF_EXISTENCE_THEOREM_FOR_ABSTRACT_WAVE_EQUATION}

\paragraph*{}
The functional calculus for self-adjoint operators
can also be applied to further classes of initial-value problems.
By the \textit{abstract {\SCHROEDINGER} equation}
we mean the linear operator-differential equation
\begin{equation}\label{ABSTRACT_SCHROEDINGER_EQUATION}
u'(t) \, + \, i \linA \, u(t) = 0\,,
\end{equation}
where $\linA$ is a self-adjoint operator in Hilbert space $\hilH$
(the operator $i \linA$ is called \textit{skew-adjoint} in this case).

\subparagraph*{}
Each solution $u$ of \eqref{ABSTRACT_SCHROEDINGER_EQUATION}
satisfies the so-called \textit{energy equation}
\begin{equation}\label{ENERGY_CONSERVATION_LAW}
\norm{u(t)} = \norm{u(0)} \mbox{ for } t > 0\,,
\end{equation}
which immediately follows from
\[
\frac{d}{dt} \, \norm{u(t)}^2
\=
\scalar{-i\linA u(t)}{u(t)} + \scalar{u(t)}{-i \linA u(t)} \= 0.
\]
If an initial value $u(0) \eq u_0$ is given,
then \eqref{ENERGY_CONSERVATION_LAW} implies
that the solution of the initial-value problem 
for \eqref{ABSTRACT_SCHROEDINGER_EQUATION} is unique.

\subparagraph*{}
To construct a solution for problem \eqref{ABSTRACT_SCHROEDINGER_EQUATION},
we define the spectral operator
\begin{equation}\label{SPECTRAL_OPERATOR_III}
L_t(x) 
\eqdef
e^{-i \linA t} x
\=
\int_{\,-\infty}^{\,\infty} e^{-i \lambda t} dP_{\lambda} x
\quad \quad (t > 0),
\end{equation}
where $P_{\lambda}$ again is the spectral resolution of $\linA$.
Then the family of operators $L_t$ consitutes the solution
of the abstract {\SCHROEDINGER} problem:

\begin{theorem}\label{EXISTENCE_THEOREM_FOR_ABSTRACT_SCHROEDINGER_EQUATION}
The mapping $t \mapsto \linL_t(u_0)$ forms a solution
of the Cauchy problem for \eqref{ABSTRACT_SCHROEDINGER_EQUATION},
where each operator $\linL_t$ for $t > 0$ is unitary.
\end{theorem}

%\input{PROOF_EXISTENCE_THEOREM_FOR_ABSTRACT_SCHROEDINGER_EQUATION}

\section{Positive Self-Adjoint Operators Generate Semigroups}
%------------------------------------------------------------
There are strong relationships between the functional calculus
for half-bounded self-adjoint operators
and the generation of strongly continuous semigroups
by maximal accretive operators-
It is the aim of this section to elaborate the similarities
between these approaches.

\paragraph*{}
Given a Banach space $\banX$,
a \textit{C$_0$-semigroup} or
a \textit{strongly continuous semigroup of bounded linear operators}
is an operator-valued mapping $\linS:[0,\infty) \xto \BO{\banX}$,
shortly denoted $\{\linS(t)\}_{t \geq 0}$
and distinguished by the properties

\begin{itemize}[leftmargin=22mm]
\item[(G-1)]
$\linS(t+s) \eq \linS(t) \circ \linS(s)$ for all $t,s \geq 0$
(the \textit{semigroup property});

\item[(G-2)]
$\lim_{\,t \downarrow 0} \, \linS(t)\,x \eq x$ for all $x \xeof X$.
\end{itemize}

\subparagraph*{}
It follows from property (G-2),
which is nothing but strong continuity from the right-hand side in $t \eq 0$,
that $\linS(0)$ is the identity mapping in $\banX$.
Notice also
that the conditions (G-1) and (G-2) together
imply continuity from the right of the semigroup even in every $t > 0$.
In case of $\banX=\R$,
semigroups share both properties
with the family $\{\cal{E}(t)\}_{\, t \geq 0}$ of exponential functions
\[
\cal{E}(t): \lambda \mapsto e^{\lambda t} \quad \quad (\lambda \xeof \R),
\]
thus $\{\cal{E}(t)\}_{\, t \geq 0}$ serves as an elementary example
of a strongly continuous semigroup.

\subsubsection*{} % orbit maps
By the \textit{orbit map} or \textit{trajectory} $\eta_x$ of a semigroup
associated with a start element $x \xeof \banX$
we mean the vector-valued mapping $\eta_x:[\,0,\infty) \xto \banX$
given by $\eta_x(t) := S(t) \, x$.
Due to condition (G-1),
strong continuity of $\{\linS(t)\}_{t \geq 0}$
implies continuity of the functions $\eta_x$.
But $\eta_x$ is even differentiable under a very mild assumption:

\begin{theorem}\label{RIGHT_DIFFERENTIABLE}
Let $\{\linS(t)\}_{t \geq 0}$ be a C$_0$-semigroup.
Then for each $x \xeof \banX$
the mapping $\eta_x$ is differentiable in $(0,\infty)$
\iff
$\eta_x$ is right differentiable in $t \eq 0$.
\end{theorem}

\subsubsection*{} % infinitesimal generators
Given a C$_0$-semigroup $\{\linS(t)\}_{t \geq 0}$,
we define the \textit{infinitesimal generator}
(briefly denoted the \textit{generator} of the semigroup)
as a mapping $\linA\:\domain{\linA} \xsubset \banX \xto \banX$
by means of its graph
\[
\Gamma_{\linA}
\eqdef
\{\,(x,y) \xeof X \xtimes X: \;
y \= \lim_{\,h \downarrow 0} \, \frac{1}{h}(\linS(h)x-x) \exists \}.
\]
Hence the domain of definition $\domain{\linA}$ of $\linA$
consists of those elements $x \xeof \banX$,
for which the orbit map $\eta_x$ is differentiable from the right in $t \eq 0$.
By theorem \ref{RIGHT_DIFFERENTIABLE},
this is equivalent to
$\domain{\linA} \eq \{\,x \xeof \banX: \eta_x \mbox{ is differentiable}\,\}$.
It is easy to show that
$\linA$ is linear,
hence it is a linear operator acting in $\banX$.

\subsubsection*{} % Stone's theorem for unitary groups
In the following we indicate the relationship between self-adjoint operators
and strongly continuous semigroups.

\subparagraph*{}
Let $\{S(t)\}_{t \geq 0}$ be a strongly continuous semigroup.
If one defines $\linU(t) := \linS(t)$ for $t \geq 0$ and $\linU(t) := \linS(-t)$ for $t < 0$,
then a \textit{group} $\{U(t)\}_{t \xeof \R}$ of \textit{unitary} operators is defined,
which fulfills
the property $\linU(s+t) \eq \linU(s) \circ \linU(t)$ for all $s,t \xeof \R$.
Reversely, 
\textsc{Stone}\textit{'s theorem} claims
that the infinitesimal generator of a unitary group is skew-adjoint:

\begin{theorem}\label{STONE_THEOREM_FOR_UNITARY_GROUPS}
Each unitary group $\{\linU_t\}_{t \xeof \R}$ in a Hilbert space
is infinitesimally generated by a linear operator of type $i \linA$,
where $\linA$ is self-adjoint.
\end{theorem}

\paragraph*{} % Positive Self-Adjoint Operators as Generators of Contraction Semigroups
We finally want to show that a positive self-adjoint operator is the generator
of a contraction group.

\subparagraph*{}
It should be clear that if $\linA$ is positive and self-adjoint,
then it is m-accretive at the same time.
Hence $\linA$ being positive and self-adjoint is the infinitesimal generator
of a semigroup of contractions.
Of course, the family $\{L_t\}_{t \geq 0}$ of spectral operators
given by \eqref{SPECTRAL_OPERATOR_I} coincides with the semigroup
$\{\linS(t)\}_{t \geq 0}$ generated by $\linA$.
Thus the solution of the homogeneous Cauchy problem can be obtained either
by the functional calculus for self-adjoint operators or by the generator
theorem \ref{LUMER-PHILLIPS_THEOREM}.

\newpage
\thispagestyle{empty}
%------------------------------------------------------------------------------%
%                                   PART III                                   %
%------------------------------------------------------------------------------%
\part{Self-Adjoint Realizations of Differential Expressions}
%------------------------------------------------------------------------------%
This final section contains applications of the theory
of self-adjoint operators to linear partial differential operators
and the boundary and initial-value problems associated with them.

\subsubsection*{} % examples of Hilbert spaces
We first introduce the fundamental spaces of functions,
in which the analysis of linear partial differential operators is performed.
In fact,
the most important class of normed linear spaces
used in functional analysis, harmonic analysis and PDE theory
is constituted by so-called \textit{Lebesgue spaces}.

\subparagraph*{}
For $p \xeof [1,\infty)$,
let $(X,\Borel,\mu)$ be a measure space and
$\cal{L}^p(X)$ be the $\C$-linear space of measurable functions
$f\:X \xto \C$,
for which $\abs{f} ^p$ is Lebesgue-integrable.
For $f,g \xeof \cal{L}^p(X)$ we introduce the equivalence relation
$f \sim g$ according to $f=g$ $\mu$-almost everywhere.
Then the elements of $L^p(X)$ are
the equivalence classes in $\cal{L}^p(X)$ {\wrt} $\sim$,
that is,
$L^p(X) := \cal{L}^p(X) / \sim$.
Now,
it is possible to define a norm in $L^p(X)$ according to
\[
\norm{f}_p
\eqdef
\big( \int_{\,X} \, \abs{f} ^p \, d\mu \,\big) \, ^{\frac{1}{p}},
\]
and the \textsc{Fischer-Riesz} \textit{theorem} ensures 
that $(L^p(X),\norm{\cdot}_p)$ is a complete normed space.

\subparagraph*{} %- the spaces L2(X)
The case $p=2$ gives rise to a special attention:
here it is possible to derive the norm $\norm{\,\cdot\,}_2$ of $L^2(X)$
from the inner product
\[
\scalar{f}{g} \eqdef \int_{\,X} \, f \, \overline{g} \; d\mu
\]
by $\norm{f}_2 := \scalar{f}{f} ^{1/2}$,
which of course makes $L^2(X)$ into a Hilbert space.
Spaces of type $L^2(X)$,
where $X$ is the underlying set of a measure space $(X,\Borel{X},\mu)$,
form an important class of Hilbert spaces.
Especially in the case,
where $X=\G$ is an open subset of one of the Euclidean spaces $\Rd$,
we obtain the Hilbert spaces $\Lh$,
which proves to be separable.

\paragraph*{} %- examples of orthonormal bases in L2(I)
We finally enumerate the usual orthonormal bases
of some Hilbert spaces $L^2(\J)$,
where $\J \xsubset \R$ denotes a finite or infinite open interval:

\begin{itemize}[leftmargin=12pt]
\item[1.]
If $\J \eq (-\pi,\pi)$,
the \textit{trigonometrical functions} given by the functions
\[
\frac{1}{\sqrt{2\pi}}, \quad
\frac{1}{\sqrt{\pi}} \sin(x), \quad
\frac{1}{\sqrt{\pi}} \cos(x), \quad
\frac{1}{\sqrt{\pi}} \sin(2x), \quad
\frac{1}{\sqrt{\pi}} \cos(2x), \quad \dots
\]
form an orthonormal base in the Hilbert space $L^2((-\pi,\pi))$.

\item[2.]
In case of $\J=(-1,1)$,
an orthonormal base, founded on polynomials,
is given by the so-called \textit{Legendre polynomials}
\[
P_n(x)
\eqdef
\frac{1}{2^n n!} \, \frac{d^n}{dx^n} (x^2-1)^n,
\quad n \xeof \N,
\]
together with the constant function $P_0(x) := 1$.
These functions form the only orthogonal system in $L^2(-1,1)$
containing polynomials of any order.

\item[3.]
Let $\J$ be the half-line $(0,\infty)$.
Then the sequence of functions
\[
L_n^{\alpha}(x)
\eqdef
x^{\alpha/2} e^{-x/2}
\cdot
\cal{L}_n^{\alpha}(x),
\quad n \xeof \N
\]
forms an orthonormal base in the space $L^2(0,\infty)$,
where $\cal{L}_n^{\alpha}$ are called the \textit{Laguerre polynomials},
defined by the equations
\[
\cal{L}_n^{\alpha}(x)
\eqdef
\frac{x^{-\alpha}e^x}{n!} \, \frac{d^n}{dx^n} (e^{-x}x^{n+\alpha}),
\quad n \xeof \N.
\]

\item[4.]
In case of $\J$ being the entire set $\R$,
an orthonormal base of $L^2(\R)$ is given by the sequence
of \textit{Hermite functions}
\[
\cal{H}_n(x) \eqdef e^{-\frac{x^2}{2}} H_n(x),
\quad n \xeof \N,
\]
where $H_n$ are the \textit{Hermite polynomials}
\[
{H}_n(x)
\eqdef
(-1)^n \sum_{k_1+2k_2=n}
\frac{n!}{k_1! \, k_2!}
\,
(-1)^{k_1+k_2} \, (2x)^{k_1},
\quad n \xeof \N.
\]
\end{itemize}

\section{Fourier-Plancherel Transformation}
%------------------------------------------
The mapping known as \textit{Fourier transformation}
is an important tool in advanced analysis
and has many applications,
especially to the investigation of
linear differential operators with constant coefficients
considered in one of the Euclidean spaces $\Rd$.

\paragraph*{}
Rather Fourier transformation can be defined in various linear spaces,
it is customary to introduce this mapping first in Lebesgue space $L^1(\Rd)$.
Preliminary,
for elements $f,g \xeof L^1(\Rd)$ 
a multiplication operation $f \ast g$ is defined according to
\[
(f \ast g) \, (x) \eqdef \int_{\,\Rd} f(t) \cdot g(x-y) \, dy,
\]
called \textit{convolution} in $L^1(\Rd)$.
This mapping is commutative, associative and satisfies the estimate
\[
\norm{f \ast g}_1 \, \leq \, \norm{f} \cdot \norm{g}
\qfa f,g \xeof L^1(\Rd).
\]
Since $L^1(\Rd,+)$ is a Banach space,
the so-called \textit{convolution algebra} $(L^1(\Rd),+,\ast)$
constitutes a commutative but non-unital Banach algebra.

\subparagraph*{}
Now {Fourier transformation} is a mapping $\FT:f \xto \hat{f}$ in $L^1(\Rd)$,
which sends each $f \xeof L^1(\Rd)$ to $\hat{f} \xeof L^{\infty}(\Rd)$,
called the \textit{Fourier transform} of $f$ and defined by
\begin{equation}\label{L1_FOURIER_TRANSFORM}
\hat{f}(\xi) \eqdef
\int_{\,\Rd} \, e^{-2 \pi i \scalar{\xi}{x}} \, f(x) \, dx
\quad \quad (\xi \xeof \Rd),
\end{equation}
where $\scalar{\xi}{x}$ is the usual inner product of $\xi,x \xeof \Rd$.

\subsubsection*{} % properties of Fourier transformation
We collect some remarkable properties of the mapping $\FT$
to the following list:

\begin{itemize}[leftmargin=12pt,itemsep=5pt]
\item[1.]
For each $f \xeof L^1(\Rd)$
the function $\hat{f}$ is continuous on $\Rd$,
so $\hat{f} \xeof C(\Rd)$.

\item[2.]
$\hat{f}$ is vanishing at infinity,
that is,
Fourier transformation constitutes a mapping $\FT:L^1(\Rd) \xto C_0(\Rd)$
(\textit{Riemann-Lebesgue lemma}).

\item[3.]
$\FT$ is injective, {\ie} $\FT(f) \neq \FT(g)$ for $f \neq g$.

\item[4.]
$\norm{\FT(f)}_{\infty} \, \leq \,
\norm{f}_1$ for $f \xeof L^1(\Rd)$.

\item[5.]
$\FT(\alpha f + \beta g) \eq \alpha \, \FT(f) + \beta \, \FT(g)$
for $\alpha,\beta \xeof \C$,
thus $\FT$ is linear.

\item[6.]
$\FT(f \ast g) \eq \FT(f) \cdot \FT(g)$,
so $\FT$ is an algebra homomorphism between the Banach algebras
$L^1(\Rd,+,\ast)$ and $C_0(\Rd,+,\cdot)$.

\item[7.]
If $f \xeof L^1(\Rd)$ and $g_f\:x \xeof \Rd \mapsto \abs{x} f(x) \xeof L^1(\Rd)$,
then $\hat{f} \xeof C^1(\Rd)$ and
\[
\partial_j f(x) \= \widehat{-i \xi_j f(\xi)}
\quad \quad (j=1,\ldots,d).
\]
\end{itemize}

\paragraph*{} %- Fourier transformation in L2
It is possible
to extend $L^1$-Fourier transformation to the Hilbert space $L^2(\Rd)$.
If the functions $u,v:\Rd \xto \C$ are Lebesgue-integrable
then by definition they belong to the space $L^1(\Rd)$,
and their \textit{Fourier transforms} are well-defined
by formula \eqref{L1_FOURIER_TRANSFORM}.

\subparagraph*{}
However, 
if is not only a member of $L^1(\Rd)$, 
but also $u \xeof L^2(\Rd)$,
we define Fourier transformation a little bit different:
\begin{equation}\label{L2_FOURIER_TRANSFORM}
\hat{u}(\xi) \eqdef
\frac{1}{(2 \pi)^{d/2}}
\int_{\,\Rd} \, e^{-2 \pi i \scalar{\xi}{x}} \, u(x) \, dx.
\end{equation}
Then for functions $u,v \xeof L^2(\Rd)$ 
the inner product between $u$ and $v$
and the inner product of their Fourier transforms $\hat{u}$ and $\hat{v}$
are related by \textsc{Plancherel}\textit{'s formula}
\[
\scalar{u}{v} \= \scalar{ \hat{u} }{\hat{v} }.
\]
Hence \textit{Fourier transformation} $\FT:u \mapsto \hat{u} $
constitutes a partial isometry from the space
$L^{1,2}(\Rd) := L^1(\Rd) \, \cap \, L^2(\Rd)$
to the space $L^2(\Rd)$.

\subparagraph*{}
It is a crucial fact now
that $L^{1,2}(\Rd)$ proves to be dense in $L^2(\Rd)$.
Since the mapping $\FT$ is linear and continuous on $L^{1,2}(\Rd)$ ,
it is even \textit{uniformly continuous} on that space
and can be extended to an \textit{unitary},
that is,
a bijective and isometric linear operator
$\FT_P:L^2(\Rd) \xto L^2(\Rd)$,
called the \textsc{Fourier-Plancherel} \textit{operator}
or the \textit{$L^2$-Fourier transformation} in $L^2(\Rd)$.

\subparagraph*{}
But in general a function $u \xeof L^2(\Rd)$ is not an element of $L^1(\Rd)$.
Thus it is not possible to define the $L^2$-Fourier transform $\hat{u}$
of every $u \xeof L^2(\Rd)$ by means of formula \eqref{L2_FOURIER_TRANSFORM}.
However,
for $u \xeof L^2(\Rd)$ and $R > 0$ the function
\[
\hat{u} _R(\xi) \eqdef
\frac{1}{(2 \pi)^{d/2}}
\int_{\,\abs{x} \leq R}
e^{-2 \pi i \scalar{x}{\xi}} \, u(x) \, dx
\]
is well-defined,
and one may obtain the $L^2$-Fourier transform $\hat{u}$ 
as the limit of the functions $\hat{u} _R$ for $R \xto \infty$
{\wrt} the $2$-norm,
that is,
$\lim_{\,R \xto \infty} \, \norm{ \hat{u} - \hat{u} _R}_2 \eq 0$.

\section{Hilbert-Sobolev Spaces}
%-------------------------------
We continue with linear partial differential expressions and 
the notation of \textit{weak differentiability} for a function $f\:\G \xto \C$,
where $f$ is assumed to be \textit{locally integrable}
on a domain $\G \xsubset \Rd$,
that is,
$f$ is Lebesgue-integrable on each compact subset of $\G$.
Let $\Lc$ be the linear space of all such functions.

\subparagraph*{}
Preliminary,
by a $d$-dimensional \textit{multi-index} is meant a $d$-tuple
$\alpha \eq (\alpha_1,\ldots,\alpha_d)$ with values $\alpha_i \xeof \N_0$,
whose \textit{modulus} is the sum
$\abs{\alpha} := \alpha_1 + \ldots + \alpha_d$.
Multi-indices are used to shorten multiple partial derivatives
$
\partial^{\alpha} f := \partial_1^{\alpha_1}
                       \partial_2^{\alpha_2} \ldots
                       \partial_d^{\alpha_d} f
%\=
%\frac{\partial^{\abs{\alpha} }f}{\partial^{\alpha_1} x_1 \,
                                 %\partial^{\alpha_2} x_2 \, \ldots \,
                                 %\partial^{\alpha_d} x_d},
$
of a sufficiently often differentiable function $f\:\G \xto \C$,
but also to denote the multiple powers
$\xi^{\alpha} := \xi_1^{\alpha_1}
\cdot \xi_2^{\alpha_2} \cdot \ldots \cdot \xi_d^{\alpha_d}$
of a complex vector $\,\xi \eq (\xi_1,\xi_2,\ldots,\xi_d) \xeof \C^d$.

\paragraph*{}
Then $f \xeof \Lc$ is called \textit{weakly differentiable}
{\wrt} the multi-index $\alpha$,
if there is a further function $g \xeof \Lc$
such that 
\[
\int_{\,\G} g(x) \, \phi(x) \, dx
\=
(-1)^{ \abs{\alpha} } \,
\int_{\,\G} f(x) \, \PA \phi(x) \, dx
\qfa \phi \xeof \TF.
\]
We denote $\PA f := g$ in this case,
and weak differentiability arises as a linear mapping
$\PA:D_{\alpha} \xsubset \Lc \xto \Lc$,
where $D_{\alpha}$ is a suitable linear subspace of $\Lc$.

\subparagraph*{}
The concept of weak derivative is compliant with the classical derivative
in the following sense:
if $f\,$ has (classical) derivative $\PA f$,
then $f\,$ is also weakly differentiable
and both notions of derivatives coincide.
Clearly,
the reverse statement does not hold as the simple example
$f\:x \mapsto \abs{x} $ for $x \xeof (-1,1)$ shows.

Sobolev spaces play an important role in PDE theory,
since the solutions of partial differential equations
are frequently searched in these spaces.

\subparagraph*{}
We preliminary define for any integer $k \xeof \N$ 
the \textit{Hilbert-Sobolev space} $H^k(\G)$
as the set of those functions $f \xeof \Lh$,
whose weak derivatives $\PA f$ exist and belong to $\Lh$
for all $\alpha$ with $ \abs{\alpha} \leq k$,
briefly written
\[
H^k(\G)
\eqdef
\{\,f \xeof \Lh: \PA f \xeof \Lh \mbox{ for } \abs{\alpha} \leq k \, \}.
\]
The spaces $H^k(\G)$ are of course linear spaces,
and specifically for $k \eq 0$
the space $H^0(\G)$ is the Lebesgue space $\Lh$ itself.

\subparagraph*{}
For functions $f \xeof H^k(\G)$
let the mapping $f \mapsto \abs{f} _k$ be defined by
\begin{equation}\label{SOBOLEV_NORM}
\abs{f} _k \eqdef
\big(
\sum_{\abs{\alpha } \leq k}
\int_{\,\G} \, \abs{\PA f} ^2 \, dx
\big) ^{1/2}.
\end{equation}
It fulfills all the properties of a norm,
thus the space $H^k(\G)$ is actually a normed space,
and the norm defined by \eqref{SOBOLEV_NORM}
is usually called a \textit{Sobolev norm}.
With respect to the question of completeness,
it can be shown that
the spaces $H^k(\G)$ are complete.
But they also can be endowed with an inner product according to
\[
\scalar{f}{g}_k \eqdef
\sum_{\abs{\alpha} \leq k} \int_{\,\G} \;
\partial^{\alpha} f \; \partial^{\alpha} g \, dx\,,
\quad \quad f,g \xeof H^k(\G),
\]
so $H^k(\G)$ actually forms a Hilbert space.

\paragraph*{}
There is still another way to define Sobolev spaces.
Consider the subspace
\[
C_{k}^{\infty}(\G) \eqdef \{\, f \xeof \SF: \abs{f} _k < \infty \,\},
\]
and because this space is \textit{not} complete
{\wrt} the norm $\abs{\,\cdot\,} _k$,
we define
\[
\cal{H}^k(\G) := \overline{C_{k}^{\infty}(\G)}.
\]
But notice that the space $\cal{H}^k(\G)$ is defined by 
a purely topological operation,
{\ie} taking the closure of a smaller space,
so for elements in this space there is no instrinsic notion of derivative.
Hence one introduces the notion of \textit{strong derivative} for
$u \xeof \cal{H}^{k}(\G)$: 
given a multi-index $\alpha$ ($\abs{\alpha} \leq k$),
a further function $\PA u \xeof \cal{H}^k(\G)$
is defined to be the strong derivative of $u$ {\wrt} $\alpha$,
if there is a sequence $\{\phi_n\} \xeof C_k^{\infty}(\G)$
such that $\abs{\PA u - \PA \phi_n)} _k \xto 0$ for $n \xto \infty$.

\subparagraph*{}
The \textsc{Meyers-Serrin} \textit{theorem} claims
$
\cal{H}^k(\G) \cong H^k(\G)
$
in the sense of Banach space isomorphisms,
so the notions of strong and weak derivative are equivalent to each other.

\paragraph*{}
In order to model boundary conditions for partial differential equations,
a further class of Sobolev spaces is of interest.
Let C$_c^k(\G)$ be the space
of $k$-times continuously differentiable functions $f\:\G \xto \C$
with compact support in $\G$.
Since this space is not dense in $H^k(\G)$ {\wrt} the norm $\abs{\,\cdot\,} _k$,
it is nearby to define the closure
\[
\HN{k}{\G} \eqdef \overline{C_c^k(\G)} _{\norm{\,\cdot\,}_k},
\]
which is also called a \textit{Hilbert-Sobolev space},
since an inner product in $\HN{k}{\G}$ can be defined by
\[
\scalar{f}{g} _k \eqdef \sum_{\abs{\alpha}\leq k, \abs{\beta} \leq k}
 \int_{\,\G} \partial^{\alpha} f \; \partial^{\alpha} g \, dx\,,
\quad \quad f,g \xeof H^k(\G)
\]
such that $\abs{f} _1 \eq \scalar{f}{f}_k^{1/2}$.
The space $\HN{k}{\G}$ forms a closed subspace of $H^k(\G)$,
so we may regard $\HN{k}{\G}$ as a Hilbert space in its own right.

\paragraph*{}
The reason,
why compact operators are of outstanding meaning
in the analysis of linear problems,
is founded by \textsc{Rellich}'s \textit{embedding theorem} for Sobolev spaces:

\begin{theorem}\label{RELLICH_EMBEDDING_THEOREM}
If $\; \G \xsubset \Rd$ is \textit{bounded},
then for each $k \xeof \N$ there is a compact embedding
\[
\compA_m:\HN{k}{\G} \hookrightarrow \HN{k-1}{\G}.
\]
\end{theorem}

\subparagraph*{}
This theorem implies
that all Hilbert-Sobolev spaces $\HN{k}{\G}$ are compactly embedded
into the space $\Lh$.
Hence each bounded sequence within $\HN{k}{\G}$ owns a convergent subsequence
{\wrt} the norm $\norm{\,\cdot\,}_2$ of $\Lh$.
Notice that in theorem \ref{RELLICH_EMBEDDING_THEOREM}
there are no restrictions regarding the smoothness
or other properties of the boundary $\bndG$.

\section{Realizations of Linear Partial Differential Operators}
%--------------------------------------------------------------
In the sequel
we deal with the general \textit{linear partial differential expression}
\begin{equation}\label{LPDO}
\DO(x,u) \eqdef
\sum_{\abs{\alpha} \leq k} c_{\alpha}(x) \, \PA u(x)
\end{equation}
of order $k \xeof \N$
(also denoted a \textit{linear partial differential operator}),
whose coefficients $c_{\alpha}$ are real or complex-valued
functions $c_{\alpha}:\clG \xto \C$
satisfying certain regularity conditions
like measurability, boundedness, continuity or differentiability.

\subparagraph*{}
We assign to a linear differential expression $\DO$
two further possible features.
First,
$\DO$ is said to be \textit{essentially real}
if the top-order coefficients $c_{\alpha}$ with $\abs{\alpha} \eq k$
are real-valued.
Second,
if each function $c_{\alpha}$ is constant,
then $\DO$ is called a
\textit{differential expression with constant coefficients}.

\paragraph*{}
We denote a densely defined linear operator $\linA$
with domain of definition $\domain{\linA}$
and acting in a Banach space $\banX$ of functions $u\:\G \xto \C$
a \textit{realization} of the expression $\DO$,
if $\linA u \eq \DO u$ for all $u \xeof \cal{D}$,
where
\begin{itemize}[leftmargin=12pt]
\item[1.]
the set $\cal{D}$ is at least a dense subspace of $\domain{\linA}$,
\item[2.]
the expression $\DO u$ has a meaning either in the classical
or in the weak sense.
\end{itemize}
Both conditions are especially true for the choice $\cal{D} \eq \TF$,
so the space of test functions forms a \textit{core} of any
realization of $\DO$.

\subparagraph*{}
Since we are interested in operators acting in Hilbert spaces,
we define the linear operator $\DO_0$ acting in Hilbert space $\Lh$ 
to be the realization of $\DO$ with dense domain of definition
$\TF \xsubset \Lh$, {\ie}
\[
\begin{cases}[1.3]
\quad \domain{\DO_0} \eqdef \TF \subset \Lh \\
\quad \DO_0 \, \phi  \eqdef \DO \, \phi
\; \mbox{ for all } \phi \xeof \domain{\DO_0}.\\
\end{cases}
\]

\subparagraph*{}
Supposed the operator $\DO_0$ is closable,
then the closure $\DO_{min} := \closure{\DO_0 }$
is called the \textit{minimal operator} of $\DO$ in $\Lh$.

\paragraph*{}
The following theorem illustrates
that the theory of closable (and closed) linear operators
is applicable to a wide range of linear partial differential expressions:

\begin{theorem}\label{SMOOTH_DIFFERENTIAL_OPERATORS_ARE_CLOSABLE}
If a linear partial differential operator $\DO$ of order $n$
has smooth coefficients,
then the realization $\linA$ acting in $\Lh$
with domain of definition
$\domain{\linA} \eq C^n(\G) \cap \Lh$ is closable.
\end{theorem}

\Proof:
Let $(u_j)$ be a sequence with $u_j \xeof \domain{\linA}$ for $j \xeof \N$
and $u_j \xto 0$ for $j \xto \infty$.
Also,
let $(\linA u_j)$ be the sequence in $\Lh$
with $\linA u_j \xto f$ for $j \xto \infty$.
By the standard criterion for closability,
we have to show that $f \eq 0$.
Since the coefficients of $\A$ are sufficiently smooth,
there is the formal adjoint $\A^{\dag}$ associated with $\A$ such that
\[
\scalar{\A\,u_j}{\phi} \=
\int_{\,\G} \A\,u_j \,\phi\,dx \=
\int_{\,\G} u_j\,\A^{\dag}\,\phi\,dx \=
\scalar{u_j}{\A^{\dag}\,\phi}
\]
for all $\phi \xeof C_c^{n}(\G)$.
By the continuity of the inner product in $\Lh$,
performing the limit for $j \xto \infty$ leads to
\[
                      \scalar{f}      {\phi} \=
\lim_{\,j \xto \infty} \scalar{\A\,u_j}{\phi} \=
\lim_{\,j \xto \infty} \scalar{u_j}    {\A^{\dag}\,\phi} \=
                      \scalar{0}      {\A^{\dag}\,\phi} \= 0.
\]
Since $\phi \xeof C^{n}(\G)$ is arbitrary and $C^{n}(\G)$ is dense in $\Lh$,
we have $f \eq 0$.
\qed

\paragraph*{}
This theorem is especially true
if $\DO \eq \DO_c$ has constant coefficients.
In this case it is possible to determine
the domain of definition of the minimal operator $(\DO_c)_{min}$.
Indeed,
the definition of Sobolev spaces $H^k(\G)$ shows
that if $\DO_c$ is of order $k$,
then $(\DO_c)_{min}$ has domain of definition $\HN{k}{\G}$.

\section{Formally Self-Adjoint Differential Expressions}
%-------------------------------------------------------
It is often useful to combine the study of a linear differential operator
with that of the so-called \textit{formally adjoint} operator.
Let $C^m(\G)$ be the linear space
of at least $m$-times continuously differentiable complex-valued functions
and $\TF$ be the subspace of \textit{test functions}
consisting of indefinitely differentiable functions
$\phi:\G \xto \C$ \textit{with compact support},
that is,
the set
\[\supp{\phi} \eqdef \overline{\{\,x \xeof \G: \phi(x) \neq 0\,\} }\]
is bounded in $\Rd$,
so the distance between $\supp{\phi}$ and $\bndG$ is strictly positive.

\subparagraph*{}
Supposed the linear partial differential operator $\DO$ of order $2m$
\begin{equation}\label{LPDO2}
\DO(x,u) \eqdef
\sum_{\abs{\alpha} \leq 2m} c_{\alpha}(x) \; \PA u(x)
\end{equation}
has coefficients $c_{\alpha} \xeof C^{\abs{\alpha} }(\G)$,
then for $u \xeof C^m(\G)$ and $\psi \xeof \TF$
we find by repeated partial integration
\[
\int_{\,\G} \sum_{\abs{\alpha} \leq 2m}
(c_{\alpha} \, \PA u) \, \overline{\psi} \, dx
\=
\int_{\,\G} u \sum_{\abs{\alpha} \leq 2m}
\overline{\PA(\overline{c_{\alpha}} \, \psi)} \, dx.
\]
Hence there is a \textit{formally adjoint} differential expression
\begin{equation}\label{FAO}
\DO^{\dag}(x,u) \eqdef
\sum_{\abs{\alpha} \leq 2m} (-1)^{\abs{\alpha} }
\PA(\overline{c_{\alpha}(x)} \, u(x))
\end{equation}
such that the identity
\begin{equation}\label{GREEN_1}
\int_{\,\G} \DO \phi \, \overline{\psi} \, dx 
\=
\int_{\,\G} \phi \, \overline{\DO^{\dag} \psi} \, dx 
\end{equation}
holds for $\phi,\psi \xeof \TF$.
Using \textsc{Leibniz}'s rule to compute all the derivatives
$\PA(\overline{c_{\alpha}} \, u)$ in \eqref{FAO},
it can be shown
that there are functions $c_{\alpha}^*:\G \xto \C$
such that %$\DO^{\dag}$ can be represented by
\[
\DO^{\dag}(x,u) \=
\sum_{\abs{\alpha} \leq 2m}
%(-1)^{\abs{\alpha} } \,
c_{\alpha}^*(x) \; \PA u(x).
\]
Thus $\DO^{\dag}$ is uniquely defined by $\DO$
and also forms a linear differential expression of type \eqref{LPDO2}.
It is easy to verify that $\DO^{\dag}$ is elliptic \iff $\DO$ is so.
Assigning to each expression $\DO$
its formally adjoint $\DO^{\dag}$ is a self-dual operation,
so $\DO^{{\dag}{\dag}} \eq \DO$.
Especially
for all \textit{elementary} differential operators $\partial_j := d/dx_j$
of order one,
a simple computation yields
${\partial_j}^{\dag} \eq - \partial_j$ for $j=1,\ldots,d$.

\subparagraph*{}
Notice that formula \eqref{GREEN_1} holds
since at least one of the functions $\phi,\psi$ belongs to the space $\TF$. 
If we replace the test functions $\phi,\psi \xeof \TF$
by more general functions $u,v \xeof C^m(\clG)$,
the space of at least $m$-times continuously differentiable functions
$u\:\clG \xto \C$,
the identity \eqref{GREEN_1} becomes into
\begin{equation}\label{GREEN_2}
\DO u \, \overline{v} \, - \, u \, \overline{\DO^{\dag} v}
\= 
\sum_{j=1}^d \, \partial_j C_j(u, \overline{v}),
\end{equation}
where the functions $C_j$ are bilinear expressions
in $u,\overline{v}$ and their derivatives up to order $2m-1$. 
Then performing integration over $\G$,
\eqref{GREEN_2} leads to
\begin{equation}\label{GREEN_3}
\int_{\,\G}
(\DO u \, \overline{v} \, - \, u \, \overline{\DO^{\dag} v}) \, dx
\= 
\int_{\,\bndG} \, \sum_{j=1}^d \, \nu_j C_j(u,\overline{v}) \; d\sigma,
\end{equation}
where $d\sigma$ is the surface measure on the boundary $\bndG$.
Equation \eqref{GREEN_3} is called \textsc{Green}\textit{'s formula}
associated with the linear differential operator $\DO$.

\paragraph*{}
We denote a linear differential expression $\DO$ \textit{formally self-adjoint},
if $\DO^{\dag} \eq \DO$ holds.
If $\DO$ has constant coefficients,
then it is formally self-adjoint
\iff
$c_{\alpha} \xeof \R$ for $\abs{\alpha} $ even
and $c_{\alpha} \xeof i \, \R$ for $\abs{\alpha} $ odd.

\section{The Momentum Operator on Various Intervals}
%---------------------------------------------------
For an arbitrary interval $\J \xsubset \R$
we consider the first order differential operator
\[
\linA_1 \eqdef i \frac{d}{dx},
\]
known from theoretical physics as the \textit{momentum operator}.
Its classical domain of definition is the space $C^1(\J)$
consisting of continuously differentiable complex-valued functions.
The momentum operator is the simplest differential operator one can imagine,
so it is natural to study this operator first.
In advance,
the momentum operator defined on domains with homogeneous boundary conditions
proves to be symmetric but no half-bounded,
so we cannot apply the \textsc{Friedrichs} extension theory to obtain
self-adjoint extensions of $\linA_1$.

\subparagraph*{}
Basically,
for each $\lambda \xeof \C$ the function $f_{\lambda}(x) := e^{-i \lambda x}$
is a solution of the equation $\linA_1 f \eq \lambda f$,
hence $f_{\lambda}$ is an eigenvector of $\linA_1$
{\wrt} the eigenvalue $\lambda$.
Thus
the momentum operator has eigenvectors for any complex number,
but they do not necessarily belong to the Hilbert space $L^2(\J)$.

\subsubsection*{} % self-adjointness in the real line
We want to determine the spectrum and
the possible self-adjointness of the momentum operator
on various intervals in the real line and will start
with the case $\J \eq (-\infty,\infty)$.
If the domain of definition $\domain{\linA_1}$ is the space $C_c^{\infty}(\R)$
of test functions,
then $\linA_1$ is a densely defined operator in Hilbert space $L^2(\R)$. 
For $f,g \xeof \domain{\linA_1}$ we compute
\[
\scalar{\linA_1 f}{g} \=
i \int_{\R} f(t) \overline{g(t)}
\, dt \= 
i f \overline{g} \mid_{-\infty}^{+\infty} 
-i \int_{\R} f'(t) \overline{g(t)}\, dt
\= 
\scalar{f}{\linA_1 g},
\]
so $\linA_1$ is symmetric on $\domain{\linA_1}$.
Recall that the Sobolev spave $H^1(\R)$ is the closure of $C_c^{\infty}(\R)$
{\wrt} the Sobolev norm $\abs{\,\cdot\,} _1$,
which is also the graph norm of the momentum operator.
Hence the domain of definition of the operator $\overline{\linA_1}$
is the space $H^1(\R)$.

\subparagraph*{}
Now,
for $f \xeof H^1(\R)$
the Fourier transform of the function $\linA_1 f$ is given by
\[
\widehat{\linA_1 f}(\xi) 
\=
\int_{\R} e^{-it \xi} \, (i \frac{d}{dt}f(t)) \, dt
\=
\xi \, \hat{f}(\xi),
\]
so the momentum operator is unitarily equivalent to the multiplication
operator $M_{\1}$ generated by the identity mapping on $\R$.
As we have seen,
the operator $M_{\1}$ is self-adjoint on $L^2(\R)$ and has pure continuous
spectrum $\sigma(M_{\1})=\R$.
Due to the unitary equivalence of $M_{\1}$ and $\linA_1$,
both assertions must also be true for the momentum operator $\linA_1$
with domain of definition $H^1(\R)$.
Clearly,
each operator defined on the space $C_c^{\infty}(\R)$ is a core of its closure,
thus the momentum operator with domain of definition $C_c^{\infty}(\R)$
is essentially self-adjoint on the Hilbert space $L^2(\R)$. 

\paragraph*{}
Unfortunately,
the results for a bounded or half-bounded interval are not as satisfying as
in the case described above.
In the following,
we shall compute the deficiency spaces and deficiency indices $n_{\pm}$
of the closure of $\linA_1$ for the intervals $(0,\infty)$ and $(0,1)$.
Remember
that a symmetric linear operator can only have self-adjoint extensions
iff $n_+=n_-$.

\begin{itemize}[leftmargin=5mm]
\item[1.]
The interval $\J := (0,\infty)$.
The operator $\linA_1$ is symmetric
on the dense domain of definition $C_c^{\infty}(0,\infty)$.
The domain of definition of its closure $\overline{\linA_1}$ is
the Sobolev space $H_0^1(0,\infty)$,
and the domain of definition of the operator $\overline{\linA_1}^*$
is the Sobolev space $H^1(0,\infty)$.
We have $u \xeof D_+ := \kernel{\overline{\linA_0}^*-i}$, iff
$u$ is the solution of the differential equation $iu'=iu$,
whose solution $u(t)=e^{t}$ is not a $L^2$-function,
so we have $D_+ \eq \{\0\}$ and $n_+ \eq 0$, respectively.

\subparagraph*{}
Similarly,
$v \xeof D_- := \kernel{\overline{\linA_0}^{\,*} +i}$ iff
$v$ is the solution of the differential equation $iv'=-iv$.
But since its solution $v(t)=e^{-t}$ is $L^2$-integrable,
if follows that
$D_-$ is a non-trivial subspace of $L^2(0,\infty)$ and $n_-=1$.
In summary,
the deficiency indices of $\overline{\linA_1}$ are $(0,1)$,
thus there is no self-adjoint extension of this operator.

\item[2.]
The interval $\J := (0,1)$.
The operator $\linA_1$ is symmetric on the dense domain of definition
$C_c^{\infty}(0,1)$ and closed if $\domain{\linA_1}$
is the Sobolev space $H_0^1(0,1)$.
Moreover,
the domain of definition of the adjoint operator $\overline{\linA_1}^*$
is the Sobolev space $H^1(0,1)$.

\subparagraph*{}
Let us compute the deficiency spaces $D_{\pm}$ of $\overline{\linA_1}^*$.
We have $u \xeof D_{\pm}$
if and only if
$u$ satisfies the differential equation
\[
i \, \frac{du}{dt} \= \pm \, i \, u,
\]
which has the solutions $u_{\pm} \eq e^{\pm t}$.
Hence the deficiency indices are $n_+ \eq n_- \eq 1$,
so $\overline{\linA_1}$ owns self-adjoint extensions.

\subparagraph*{}
Indeed,
there is a one-parameter family $S_{\alpha}$ ($\alpha \xeof [0,2\pi)$)
of self-adjoint extensions for this operator, 
and the domain of definition of each extension is described by the set
\[
\domain{S_{\alpha}} \eqdef \set{f \xeof AC[0,1]}{f(1)=e^{i \alpha} f(0)}.
\]
Remember
that the set $AC[0,1]$ consists of almost everywhere differentiable functions, 
so we have $S_{\alpha}f := i f'$ for $f \xeof \domain{S_{\alpha}}$.

\subparagraph*{}
Clearly,
since there is no unique self-adjoint extension,
the operator $\linA_1$ with domain of definition
$\domain{\linA_1} \eq C_c^{\infty}(0,1)$
cannot be essentially self-adjoint in Hilbert space $L^2(0,1)$.
\end{itemize}

\section{Sturm-Liouville Operators}
%----------------------------------
A further example of a linear differential operator (in divergence form)
is given by the \textsc{Sturm-Liouville} \textit{operator}
\[
\DO_{SL}(u,x) \eqdef -(p(x)\, u'(x))' \, + \, q(x)\, u(x),
\quad \quad x \xeof (a,b)
\]
with $p(x) > 0$ and $q(x) \geq 0$ for $x \xeof [a,b]$,
where $-\infty < a < b < \infty$.
This operator models a vibrating string
being fixed at the endpoints of a finite interval $[a,b]$.

\subparagraph*{}
The corresponding Dirichlet form reads as
\[
\DF_{\DO_{SL}}(u,v) \=
\int_a^b \,
(p \, u'\, \overline{v}' \, + \, q \, u \, \overline{v}) \, dx,
\]
which shows that $\DO_{SL}$ is symmetric and half-bounded
on the domain of definition $\TF$.
The energetic space $H_{\DF_{SL}}$ is the Hilbert-Sobolev space $\HN{1}{a,b}$,
containing all nearly everywhere differentiable functions
$u \xeof C[a,b]$ with $u(a) \eq u(b) \eq 0$.

\subparagraph*{}
The self-adjoint extension $\DO_{SL,F}$ of the Sturm-Liouville operator
in Hilbert space $L^2(a,b)$ has domain
\[
\domain{\DO_{SL,F}} = H^2(a,b) \, \cap \, \HN{1}{a,b}.
\]
It is a linear operator with pure point spectrum,
since $(a,b)$ is a bounded subset of the real line.
Especially for $p(x)=1$ and $q(x) \eq 0$ for $x \xeof (a,b)$ it is possible
to compute explicitely the eigenvalues and eigenfunctions
of the operator $\DO_{SL,F} \eq - {d^2} / {dx^2}$.
In fact,
in this case the eigenvalue distribution is represented by
\[
\lambda_n \= \frac{n^2 \pi^2}{(b-a)^2} \quad \quad (n \xeof \N),
\]
where each $\lambda_n$ owns a one-dimensional eigenspace,
and the related eigenfunctions $u_n$ are the mappings
\[
u_n(x) \eqdef \sin \, (\frac{n \pi x}{b-a}) \; \for x \xeof (a,b)\,
\]
forming an orthonormal base in Hilbert space $L^2(a,b)$.
Thus,
the smallest eigenvalue is the number $\pi^2/(b-a)^2$,
and the asymptotical behaviour of the eigenvalues
is $\lambda_n \sim n^2$.

\section{The Laplacian in Bounded and Unbounded Domains}
%-------------------------------------------------------
This subsection is devoted to the \textit{Laplace operator},
which is the most important second order differential operator
in mathematical physics.
In general,
let $\G \xsubset \Rd$ be an open domain, 
then in the classical sense the Laplace operator
(briefly denoted \emph{Laplacian})
is the linear mapping $-\Delta: C^2(\G) \xto C(\G)$ defined by
\[
\Delta \, u \eqdef \frac{\partial^2 u}{\partial x_1^2} +
                     \frac{\partial^2 u}{\partial x_1^2} +
                     \ldots                              +
                     \frac{\partial^2 u}{\partial x_n^2}.
\]
This operator can be considered as a prototype
of a second order \textit{elliptic} differential operator.

\paragraph*{}
In order to apply the theory of self-adjoint operators
to the investigation of the Laplacian,
we firstly have to find a suitable dense subspace
of the Hilbert space $L^2(\G)$,
on which $\Delta$ can be defined and fulfills the symmetry property.
This will be the space $C_c^{\infty}(\G)$ of test functions.
Notice
that the Laplacian is an unclosed operator on this domain.

\subparagraph*{}
Using Green's second formula,
and due to the boundary condition $\phi(x) \eq 0$ for $x \xeof \partial \G$
and $\phi \xeof C_c^{\infty}(\G)$,
we compute for $\phi,\psi \xeof C_c^{\infty}(\G)$
\[
\int_{\G} \Delta \phi \, \overline{\psi} \, dx
\=
- \sum_{i=1}^n \int_{\G}
\frac{\partial^2 \phi}{\partial x_i^2} \, \overline{\psi} \, dx \=
- \sum_{i=1}^n \int_{\G}
u \, \overline{\frac{\partial^2 \phi}{\partial x_i^2}} \, dx \=
\int_{\G} \phi \, \overline{\Delta \psi} \, dx,
\]
so $\scalar{\Delta \phi}{\psi} \eq \scalar{\phi}{\Delta \psi}$,
which implies symmetry of the operator $\Delta$
on the space $C_c^{\infty}(\G)$.

\subparagraph*{}
Moreover,
we get for the operator $- \Delta$ the inequality
\[
\scalar{- \Delta \phi}{\phi} \=
\sum_{i=1}^n \int_{\G}
\frac{\partial^2 \phi}{\partial x_i^2} \, \overline{\phi} \, dx
\=
\sum_{i=1}^n \int_{\G} \abs{ \frac{\partial \phi}{\partial x_i}} ^2 dx 
\, \geq 0,
\]
so $-\Delta$ is positive.
Since the Laplacian is a half-bounded operator on the domain
$C_c^{\infty}(\G)$ with lower bound $c_{\linA} \eq 0$,
it fulfills all the prerequisites for
Friedrichs' extension method.
Thus we can state the following

\begin{theorem}
The operator $\linA := -\Delta$ with domain
$\domain{\linA} := C_c^{\infty}(\G)$ has at least
one self-adjoint extension.
\end{theorem}

It should be mentioned 
that when working with test functions $\phi \xeof C_0^{\infty}(\G)$
there was no need
to make a distinction between bounded and unbounded domains in the
computations above.

\subsubsection*{} % realizations of the Laplacian
In the following we consider some realizations of the Laplacian
in various domains:

\begin{itemize}[leftmargin=5mm]
\item[1.]
\textit{Laplacian in $\Rd$}.
Let $\G$ be the entire Euclidean space $\Rd$.
The analysis of the Laplacian will be as easy as for the momentum operator
in this case.
Since the operator $-\Delta$ is unitarily equivalent to
the multiplication operator $M_{-\Delta}$ generated by the function
\[
m_2:(x_1,\ldots,x_n) \mapsto (x_1^2 + \ldots + x_n^2), \quad x \xeof \Rd
\]
and the operator $M_{m_2}$ is real-valued and self-adjoint,
it follows
that the operator $-\Delta$ with domain of definition $H^2(\Rd)$
must be self-adjoint,
too.
The spectrum of $-\Delta$ is that of the operator $M_{m_2}$,
that is,
it is purely continuous with $\sigma_c(-\Delta)=[0,\infty)$.
Moreover,
since $-\Delta$ with domain of definition $H^2(\Rd)$
is the closure of $-\Delta$ with domain of definition $C_c^{\infty}(\Rd)$, 
the operator $-\Delta$ with domain of definition $C_c^{\infty}(\Rd)$ is
essentially self-adjoint on the Hilbert space $L^2(\Rd)$.

\subparagraph*{}
The energetic norm of the Laplacian for a function
$\phi \xeof \domain{-\Delta_F}$ is given by the formula
\[
\norm{\phi}_{-\Delta} \eqdef
\int_{\Rd} \abs{\phi} ^2 dx 
\, + \,
\sum_{i=1}^n \, \int_{\Rd} \frac{\partial \phi}{\partial x_i} \, dx.
\] 
Since it is the well-known Sobolev norm $\abs{\phi} _1$,
the domain of definition $\domain{-\Delta_F}$ is the Sobolev space $H_0^1(\Rd)$,
which in case of $\G=\Rd$ is identical to the space $H^1(\Rd)$.

\subparagraph*{}
The domain of definition of the adjoint operator $-\Delta^*$ is the subspace
\[
\domain{-\Delta^*} \= \set{f \xeof L^2(\Rd)}{-\Delta f \xeof L^2(\Rd)},
\]
which is equal to the space $H^2(\Rd)$ itself.
Recall
that a half-bounded operator $\linA$ is self-adjoint on the intersection set
$\domain{\linA^*} \cap \domain{T_F}$,
so we conclude
that the Laplacian must be self-adjoint on the domain
\[
\domain{-\Delta_F} \= H^2(\Rd) \cap H_0^1(\Rd) =H^2(\Rd)
\]
in compliance with the statement above.

\item[2.]
\textit{Dirichlet-Laplacian}.
Let $\G \xsubset \Rd$ be a \textit{bounded} domain.
Consider the Laplacian $-\Delta$ with domain
$\domain{-\Delta}:=C_c^{\infty}(\G)$.
The related sesquilinear form $Q$ to the Laplacian is given by 
\[
Q(\phi,\psi) \eqdef 
\int_{\G} \nabla \phi \, \nabla \psi \, dx,
\quad \quad
\phi,\psi \xeof \domain{-\Delta},
\]
which can continuously be extended to a sesquilinear form $\overline{Q}$
on the domain
$\domain{\overline{Q}} := \HN{1}{\G} \xtimes \HN{1}{\G}$.

\subparagraph*{}
Due to the \textsc{Friedrichs} extension procedure,
there is a self-adjoint operator $\linA_D$ on the Hilbert space $L^2(\G)$
related to $\Q$ and defined by
\[
\begin{cases}
\quad \domain{\linA_D} & := \,
\set{f \xeof \HN{1}(\G)}{\Delta f \xeof L^2(\G)} \\
\quad \linA_D f & := \,
- \Delta f, \quad f \xeof \domain{\linA_D}. \\
\end{cases}
\] 
The operator $\linA_D$ is called the
\textit{Laplace operator with Dirichlet boundary conditions},
often it is shortly called the \textit{Dirichlet-Laplacian}.
Obviously,
we have
\[
\HN{1}{(\G} \cap H^2(\G) \subset \domain{\linA_D}.
\]
One can prove that
if $\G$ is bounded and has a $C^2$-boundary or $\G \eq \R_+^n$,
then the reversed inclusion is also true,
and in this case $\domain{\linA_D} \eq \HN{1}{\G} \cap H^2(\G)$.
This equality still holds in the case in which $\G$ is convex and bounded.

\paragraph*{}
The main result concerning the Dirichlet-Laplacian is given by the following

\begin{theorem}
The operator $\linA_D$ as defined above is self-adjoint in Hilbert space $\Lh$.
The domain of definition $\domain{\linA_D}$ equipped with the graph norm of 
$\linA_D$ is continuously embedded in $\HN{1}{\G}$.
If $\G$ has $C^1$-boundary,
then $\linA_D$ is compactly embedded in $\Lh$.
\end{theorem}

From the last assertion it follows immediately that
for a bounded set with $C^1$-boundary the operator $\linA_D$ has
pure point spectrum.
It is an important property of the Dirichlet-Laplacian in this case
that $0$ does not belong to the spectrum,
so $\linA_D$ is invertible,
and the inverse operator $\linA_D^{-1}$ is bounded.

\item[3.]
\textit{The Neumann-Laplacian}.
Let $\G \xsubset \Rd$ be a domain with $C^1$-boundary $\partial \G$,
let $\hilH$ be the Hilbert space $\Lh$ and let the operator $\linA_N$
be defined by
\[
\begin{cases}
\quad \domain{\linA_N} & := \,
\{ f \xeof H^2(\G): f_v(x)= 0, x \xeof \partial \G \} \\
\quad \linA_N f & := \,
- \Delta f, \quad f \xeof \domain{\linA_N}. \\
\end{cases}
\]
The operator $\linA_N$ is called the
\textit{Laplace operator with Neumann boundary condition}
or briefly the \textit{Neumann-Laplacian}.
In distinction to the Dirichlet-Laplacian,
the Neumann-Laplacian has $0$ as an eigenvalue,
and all constant functions in $\domain{\linA_N}$
are corresponding eigenfunctions.
\end{itemize}

\section{{\SCHROEDINGER} Operators}
%----------------------------------
By the {\SCHROEDINGER} operator is meant
the linear partial differential operator given by
\[
\DO_q(u,x) \eqdef (-\Delta u)(x)\, + \, q(x),
\]
defined in the entire Euclidean space $\Rd$.
In the special case,
where $q \eq 0$,
self-adjointness of the operator $-\Delta$ acting in $L^2(\Rd)$
and with domain of definition $H^2(\Rd)$ is well-known from the last section.

\subparagraph*{}
As an application of theorem \ref{KATO-RELLICH_THEOREM_II}
coming from perturbation theory,
consider the case,
where $d=3$ and the function $q$ belongs to $L^2(\R^3) + L^{\infty}(\R^3)$,
that is,
\[
q = q_1 + q_2
\]
for some functions $q_1, q_2$
with $q_1 \xeof L^2(\R^3)$ and $q_2 \xeof L^{\infty}(\R^3)$.
Then the function $q$ gives rise to a multiplication operator 
with domain of definition $C_c^{\infty}(\R^3)$.
Since the negative Laplacian $-\Delta$ is essentially self-adjoint 
on that domain
and $q$ is $-\Delta$-bounded with bound $\theta < 1$,
the operator $\DO_q \eq -\Delta+q$
with domain of definition $C_c^{\infty}(\R^3)$ is also essentially self-adjoint.
The self-adjoint extension $\HAM_q$ of $\DO_q$ 
is the generator of the unitary group $\{U(t)\}_{t \xeof \R}$ with 
\[
U(t) := e^{it \HAM_q} \quad \quad (t \xeof \R),
\]
which uniquely solves the \textit{time-dependent} {\SCHROEDINGER} problem
\[
\begin{cases}[1.2]
\quad u_t = i \cdot \HAM_q u,\\
\quad u(0) = u_0.\\
\end{cases}
\]

\subsubsection*{} % spectral analysis of Schroedinger operators
The investigation of \textit{time-independent solutions}
of the {\SCHROEDINGER} equation leads to the eigenvalue problem
\begin{equation}\label{SCHROEDINGER_EIGENVALUE_PROBLEM}
\HAM_q u \eqdef -\Delta u + qu \= \lambda u,
\end{equation}
which means to find those real numbers $\lambda$,
for which there are non-trivial solutions $u$
of problem \eqref{SCHROEDINGER_EIGENVALUE_PROBLEM}.
Thus the discrete spectrum $\sigma_d(\HAM_q)$ of the {\SCHROEDINGER} operator
is of interest.
In fact,
the non-trivial eigenfunctions belonging to eigenvalues
$\lambda \xeof \sigma_d(\HAM_q)$ describe
the stationary states of a quantum mechanical system.

\subparagraph*{}
The presence of a non-trivial potential function $q$
can fairly change the spectrum of a {\SCHROEDINGER} operator 
like the example of the \textit{harmonic oscillator} shows:
the spectrum of the Laplacian equals the connected set $[0,\infty)$,
but the spectrum of the operator $-\Delta + q$ 
with $q:x \mapsto \frac{1}{2}x^2$ consists of discrete points only.

\paragraph*{}
It is a further and interesting problem in quantum mechanics
to determine lower and upper bounds for the discrete spectrum of the
operator $\HAM_q$.
In the following we present three assertions,
which contain information on the bounds of the set $\sigma_d(\HAM_q)$,
namely
the \textsc{Birman-Schwinger} \textit{bound}, 
the \textsc{Lieb-Cwikel-Rozenblum} \textit{bound}, and
    \textsc{Kato}\textit{'s theorem}.

\begin{itemize}[leftmargin=12pt]
\item[1.]
\textsc{Birman-Schwinger} \textit{bound.}
The following assertion presents bounds for the set $\sigma_d(\HAM_q)$,
which are valid for certain potentials $q$ in case of $d=3$.
Preliminary,
let $q \xeof L_{\infty}(\R^3)$ and 
\[
B_q \eqdef
\int_{\R^3} \int_{\R^3} \frac{q(x) q(y)}{\abs{x-y} ^2} \ dx \, dy
\, < \, + \infty.
\]
Then,
if $N(\HAM_q)$ is the total number of eigenvalues of the operator $\HAM_q$,
then 
\[
N(\HAM_q) \, \leq \, \frac{1}{16 \pi^2} \, B_q.
\]

\item[2.]
The \textsc{Lieb-Cwikel-Rozenblum} \textit{theorem}
provides an estimate on the number $N_(\HAM_q)$
of negative eigenvalues of the {\SCHROEDINGER} operator:
let $d > 3$ and $q$ be an element of the space $L_{loc}^{\infty}(\Rd)$,
then
\[
N(\HAM_q) \, \leq \, c_n \int_{\Rd} \min \{q(x),0\}^{n/2} \, dx,
\]
where $c_n$ is a constant depending only from the space dimension $d$.

\item[3.]
The subsequent result is often denoted as \textsc{Kato}\textit{'s theorem}.
Indeed, 
it does not describe bounds of the discrete spectrum,
but quotes the absence of positive eigenvalues under relative mild
assumptions on the potential function $q$:
if $q \xeof L_{loc}^{\infty}(\Rd)$ and
\[
\lim_{\abs{x} \xto \infty} \, \abs{x} \, q(x) \= 0,
\]
then the operator $-\Delta+q$ does not own any positive eigenvalue.
\end{itemize}

\section{Strongly Elliptic Differential Operators of Order $2m$}
%---------------------------------------------------------------
A \textit{linear elliptic partial differential operator} $\DO$
of even order $2m$ on a domain $\G \subseteq \Rd$ is a mapping of the form
$
\DO \, u \eq \sum_{\abs{\alpha} \leq 2m} a_{\alpha} D^{\alpha} u,
$
where the functions $a_{\alpha}$ are for examples
members of the space $C^{2m}(\G)$ and whose \textit{principal symbol}
\[
a(x,\xi) \, = \sum_{\abs{\alpha} = 2m} a_{\alpha}(x) \,
\xi^{\alpha}, \quad \xi \xeof \Rd
\]
satisfies $a(x,\xi) \neq 0$ for all $x \xeof \G$ and
$\xi \xeof \Rd \setminus \{\0\}$.
We will briefly call $\DO$ an \textit{elliptic operator}.
Furthermore,
$\DO$ is called \textit{uniformly strongly elliptic},
if there is a constant $c > 0$ such that
\[
\sum_{\abs{\alpha} = \abs{\beta} = m} a_{\alpha \beta}(x) \,
\xi^{\alpha} \xi^{\beta}
\, \geq \,
c \sum_{\abs{\alpha} =m} \xi^{2 \alpha}.
\]

\paragraph*{}
Elliptic operators are usually considered under the following assumptions:

\begin{itemize}[leftmargin=22mm]
\item[(E-1)]
The functions $a_{\alpha}$ are real-valued for $\abs{\alpha} \eq 2m$,
and either $a(x,\xi)< 0$ or $a(x,\xi) > 0$ holds for all $x \xeof \G$
and all $\xi \xeof \Rd \setminus \{\0\}$.
This implies
that for any compact set $\compK \xsubset \G$
there is a constant $C_{\compK} > 0$ such that
\[
a(x,\xi) \, \geq \, C_{\compK} \, \abs{\xi} ^{2m}
\quad \quad (x \xeof \G,\;\xi \xeof \Rd).
\]

\item[(E-2)]
The expression $\DO$ can be written in \textit{divergence form}
\[
\DO(x,\cdot) \= \sum_{\abs{\alpha} \leq m, \abs{\beta} \leq m}
(-1)^{\alpha} D^{\alpha}(a_{\alpha \beta}(x)\,D^{\beta}),
\]
where $a_{\alpha \beta} \eq \overline{a_{\beta \alpha}}$
for all $\alpha$ and $\beta$.
Notice that every linear differential operator with sufficiently smooth
coefficient functions can be written in divergence form.
\end{itemize}

\subparagraph*{}
As for the differential operators mentioned in the previous sections,
a formally self-adjoint and uniformly strongly elliptic differential operator
$\DO$ causes a symmetric and quasi-coercive sesquilinear form
defined on the space $\HN{m}{\G} \xtimes \HN{m}{\G}$,
so the \textsc{Friedrichs} extension theory gives rise
to a self-adjoint realization $\DO_F$ in $\Lh$
with domain of definition
\[
\domain{\DO_F} \= \{u \xeof \HN{m}{\G} \: \DO_F \, u \xeof \Lh\}.
\]
If the boundary of $\G$ is sufficiently smooth,
however,
then
\[
\domain{\DO_F} \eq H^{2m}(\G) \, \cap \, \HN{m}{\G}
\]
holds.
It is also true that $\DO_F$ has pure point spectrum, 
if $\G$ is bounded in $\Rd$.

%------------------------------------------------------------------------------%
%                              BIBLIOGRAPHY                                    %
%------------------------------------------------------------------------------%
\newpage
\thispagestyle{empty}
\begin{thebibliography}{99}
\begin{small}
\bibitem{AG}  N.I.\,Akhiezer, I.M.\,Glazman: \textit{Theory of Linear Operators in Hilbert Space}, Dover Publications, 1993.
\bibitem{BS}  M.S.\,Birman, M.Z.\,Solomjak:  \textit{Spectral Theory of Self-Adjoint Operators in Hilbert Space}, Springer Verlag, 1987.
\bibitem{Br2} F.\,Browder:                   \textit{{\FA} and Partial Differential Equations I}, Math. Annalen 138, pp. 55-79, 1959.
\bibitem{Br3} F.\,Browder:                   \textit{{\FA} and Partial Differential Equations II}, Math. Annalen 145, 81-226, 1962.
\bibitem{EN}  K.J.\,Engel, R.\,Nagel:        \textit{A Short Course on Operator Semigroups}, Universitext, Springer Verlag, 2006.
\bibitem{HP}  E.\,Hille, R.\,Phillips:       \textit{{\FA} and Semigroups}, AMS Vol.31, 1948.
\bibitem{Ka}  T.\,Kato:                      \textit{Perturbation Theory for Linear Operators}, Springer Verlag, Grundlehren der mathematischen Wissenschaften, Band 132, 1995.
\bibitem{Kr}  S.\,Krein: \textit{Linear Equations in Banach Spaces}, Birkh\"{a}user Verlag, 1982.
\bibitem{RN}  F.\,Riesz, B.\,Sz.-Nagy:       \textit{{\FA}}, Springer Verlag., 1956.
\bibitem{RS1} M.\,Reed, B.\,Simon:           \textit{Modern Methods of Mathematical Physics I: {\FA}}, Academic Press Inc., 1980.
\bibitem{RS4} M.\,Reed, B.\,Simon:           \textit{Modern Methods of Mathematical Physics IV: Analysis of Operators}, Academic Press Inc., 1978.
\bibitem{Tr}  H.\,Triebel:                   \textit{Higher Analysis}, Johann Ambrosius Barth Verlag, 1992.
\bibitem{We}  J.\,Weidmann:                  \textit{Linear Operators in Hilbert Spaces}, Springer Verlag, 1980.
\bibitem{Yo}  K.\,Yosida:                    \textit{{\FA}}, Springer Verlag, 1980.
\end{small}
\end{thebibliography}

%------------------------------------------------------------------------------%
%                             END OF DOCUMENT                                  %
%------------------------------------------------------------------------------%
\end{document}
